\documentclass[11pt,a4paper]{article}
\usepackage[T1]{fontenc}
\usepackage[utf8]{inputenc}
\usepackage[brazil]{babel}
\usepackage{amsmath,amssymb,textcomp}
\usepackage[margin=2.2cm]{geometry}
\usepackage{enumitem}
\usepackage{array}
\setlength{\parindent}{0pt}
\setlength{\parskip}{0.6em}

\begin{document}
\section*{Banco curado de questões}
Questões geradas pelo sistema multiagente e aprovadas pelo ciclo completo (verificação simbólica e avaliação didática). Versão do professor, com gabarito e resolução. Licença CC BY-SA 4.0.
\subsection*{EM13MAT302}
{\small\itshape Resolver e elaborar problemas cujos modelos são as funções polinomiais de 1º e 2º graus, em contextos diversos, incluindo ou não tecnologias digitais.}

\bigskip\noindent\textbf{Q001 · discursiva · fácil · aplicar · aplicada · gabarito conferido}\par\nopagebreak
Uma empresa de transporte por aplicativo cobra, em cada corrida, uma taxa fixa somada a um valor proporcional à distância percorrida. Um cliente fez uma corrida de 4 km e pagou R\$\,13,00. Em outro dia, o mesmo cliente percorreu 10 km e pagou R\$\,25,00. Sabe-se que o preço da corrida varia linearmente com a distância percorrida.

a) Determine a expressão algébrica que fornece o preço C, em reais, de uma corrida em função da distância d, em quilômetros.

b) Usando essa expressão, calcule quanto custará uma corrida de 15 km.

\textbf{Gabarito.} C(d) = 2d + 5; para d = 15 km, o custo é R\$\,35,00.

\textbf{Resolução.} \textbf{Modelando a situação}

Como o preço varia linearmente com a distância, o modelo é uma função afim:
$$C(d) = a \cdot d + b$$
onde $a$ é o valor cobrado por km e $b$ é a taxa fixa.

\textbf{Usando os dados para montar um sistema}

Com $d=4$, $C=13$: $\;13 = 4a + b$

Com $d=10$, $C=25$: $\;25 = 10a + b$

\textbf{Resolvendo o sistema}

Subtraindo a primeira equação da segunda:
$$25 - 13 = 10a - 4a \implies 12 = 6a \implies a = 2$$

Substituindo em $13 = 4a+b$:
$$13 = 4(2) + b \implies 13 = 8 + b \implies b = 5$$

\textbf{Modelo obtido}
$$C(d) = 2d + 5$$

\textbf{Item b)}

Para $d = 15$:
$$C(15) = 2(15) + 5 = 30 + 5 = 35$$

Logo, a corrida de 15 km custará \textbf{R\$\,35,00}.

\bigskip\noindent\textbf{Q003 · discursiva · média · aplicar · teórica · gabarito conferido}\par\nopagebreak
Uma empresa de transporte por aplicativo cobra uma tarifa composta por um valor fixo (bandeirada) somado a um valor proporcional à distância percorrida, medida em quilômetros. Um passageiro que percorreu 4 km pagou R\$\,18,00 pela corrida, e outro passageiro que percorreu 10 km pagou R\$\,36,00.

a) Determine a lei que expressa o custo $C(d)$ de uma corrida (em reais) em função da distância $d$ percorrida (em quilômetros), sabendo que essa relação é modelada por uma função polinomial de 1º grau.

b) Utilizando a lei obtida, calcule quanto pagará um passageiro que percorrer 15 km.

c) Determine a distância percorrida por um passageiro que pagou R\$\,81,00 pela corrida.

\textbf{Gabarito.} C(d) = 3d + 6; custo para 15 km = R\$\,51,00; distância para R\$\,81,00 = 25 km

\textbf{Resolução.} \textbf{1. Reconhecendo o modelo}

Como o custo é composto por uma parte fixa (bandeirada) e uma parte que cresce proporcionalmente à distância, a situação é modelada por uma função afim:
$$C(d) = a\,d + b$$
onde $a$ é o valor cobrado por quilômetro e $b$ é a bandeirada fixa.

\textbf{2. Montando o sistema com os dados do enunciado}

Para $d = 4$ km, $C = 18$:
$$4a + b = 18$$

Para $d = 10$ km, $C = 36$:
$$10a + b = 36$$

\textbf{3. Resolvendo o sistema}

Subtraindo a primeira equação da segunda:
$$(10a + b) - (4a + b) = 36 - 18$$
$$6a = 18 \implies a = 3$$

Substituindo em $4a + b = 18$:
$$4(3) + b = 18 \implies b = 6$$

Logo, a lei da função é:
$$C(d) = 3d + 6$$

\textbf{4. Item (b): custo para 15 km}

$$C(15) = 3(15) + 6 = 45 + 6 = 51$$

O passageiro pagará \textbf{R\$\,51,00}.

\textbf{5. Item (c): distância para custo de R\$\,81,00}

$$3d + 6 = 81$$
$$3d = 75$$
$$d = 25$$

O passageiro percorreu \textbf{25 km}.

\bigskip\noindent\textbf{Q005 · discursiva · difícil · criar · aplicada · conferido em parte}\par\nopagebreak
Uma pessoa comparou dois aplicativos de mototáxi usando cada um deles em meses diferentes, sempre para o mesmo tipo de deslocamento no trabalho. Em cada aplicativo, a fatura mensal é composta por uma tarifa fixa mais um valor constante por quilômetro total rodado no mês (isto é, o custo mensal é uma função afim da quilometragem mensal $x$, em km, com $0 \le x \le 150$). Os registros de fatura foram:

\begin{itemize}[nosep,leftmargin=*]
\item MotoVá: em um mês, rodou 40 km e pagou R\$\,62,00; em outro mês, rodou 100 km e pagou R\$\,122,00.
\item UrbanMoto: em um mês, rodou 40 km e pagou R\$\,70,00; em outro mês, rodou 90 km e pagou R\$\,100,00.
\end{itemize}

a) Determine a lei $c_A(x)$ que dá o custo mensal do MotoVá e a lei $c_B(x)$ que dá o custo mensal do UrbanMoto, em função da quilometragem mensal $x$.

b) Determine a quilometragem mensal para a qual os dois aplicativos cobram exatamente o mesmo valor, e qual é esse valor. Indique qual aplicativo é mais barato para quilometragens menores que esse valor e qual é mais barato para quilometragens maiores.

c) Um terceiro aplicativo, o RotaFácil, acaba de entrar no mercado, cobrando uma tarifa fixa de R\$\,50,00 mais R\$\,0,50 por quilômetro rodado no mês. Elabore uma única função $C(x)$, definida por partes para $0 \le x \le 150$, que forneça sempre o menor custo mensal possível entre os três aplicativos, para qualquer quilometragem mensal $x$ nesse intervalo (explicite os intervalos de $x$ de cada sentença). Em seguida, justifique com os cálculos por que o UrbanMoto não aparece em nenhuma sentença da função que você construiu, ou seja, mostre que ele nunca é a opção mais barata para nenhum valor de $x$ em $[0,150]$.

\textbf{Gabarito.} c\_A(x) = x + 22; c\_B(x) = 0,6x + 46; equilíbrio A–B em x = 60 km (custo R\$\,82,00), com MotoVá mais barato para x<60 e UrbanMoto mais barato para x>60. Com o RotaFácil, c\_C(x) = 0,5x + 50, e a função de menor custo é C(x) = x+22 para 0$\le$x$\le$56 e C(x) = 0,5x+50 para 56<x$\le$150; o UrbanMoto nunca é a opção mais barata em [0,150].

\textbf{Resolução.} \textbf{a) Construindo os modelos afins}

MotoVá: dois pontos $(40,62)$ e $(100,122)$.
$$a=\frac{122-62}{100-40}=\frac{60}{60}=1,\qquad b=62-1\cdot 40=22$$
$$c_A(x)=x+22$$

UrbanMoto: dois pontos $(40,70)$ e $(90,100)$.
$$a=\frac{100-70}{90-40}=\frac{30}{50}=\frac{3}{5}=0{,}6,\qquad b=70-0{,}6\cdot 40=46$$
$$c_B(x)=0{,}6x+46$$

\textbf{b) Ponto de equilíbrio entre MotoVá e UrbanMoto}

$$x+22=0{,}6x+46 \Rightarrow 0{,}4x=24 \Rightarrow x=60$$

Custo nesse ponto: $c_A(60)=60+22=82$ reais.

Como $c_A(0)=22<c_B(0)=46$, o MotoVá é mais barato para $x<60$; como as retas se cruzam apenas em $x=60$ e $c_B$ cresce mais devagar, para $x>60$ o UrbanMoto passa a ser mais barato (por exemplo, em $x=150$: $c_A=172$ e $c_B=136$).

\textbf{c) Incluindo o RotaFácil e construindo a função de menor custo}

Lei do RotaFácil: $c_C(x)=0{,}5x+50$.

Precisamos comparar as três retas duas a duas para descobrir, em cada faixa de $x$, qual delas é a menor.

*MotoVá × RotaFácil:*
$$x+22=0{,}5x+50 \Rightarrow 0{,}5x=28 \Rightarrow x=56,\ \text{custo }78$$

*UrbanMoto × RotaFácil:*
$$0{,}6x+46=0{,}5x+50 \Rightarrow 0{,}1x=4 \Rightarrow x=40,\ \text{custo }70$$

*MotoVá × UrbanMoto:* já obtido, $x=60$.

Agora testamos o custo das três funções em pontos representativos de cada intervalo delimitado pelos cruzamentos ($x=0,\,40,\,56,\,60,\,150$):

\begin{itemize}[nosep,leftmargin=*]
\item $x=20$: $c_A=42,\ c_B=58,\ c_C=60$ $\to$ menor é $c_A$.
\item $x=40$: $c_A=62,\ c_B=70,\ c_C=70$ $\to$ menor é $c_A$ (B e C empatam, mas A é menor que ambos).
\item $x=50$: $c_A=72,\ c_B=76,\ c_C=75$ $\to$ menor é $c_A$.
\item $x=56$: $c_A=78,\ c_B=79{,}6,\ c_C=78$ $\to$ menor é $c_A=c_C$.
\item $x=58$: $c_A=80,\ c_B=80{,}8,\ c_C=79$ $\to$ menor é $c_C$.
\item $x=100$: $c_A=122,\ c_B=106,\ c_C=100$ $\to$ menor é $c_C$.
\item $x=150$: $c_A=172,\ c_B=136,\ c_C=125$ $\to$ menor é $c_C$.
\end{itemize}

Os testes mostram que o mínimo é sempre dado por $c_A$ ou por $c_C$, nunca por $c_B$. A troca entre eles ocorre exatamente no cruzamento de $c_A$ com $c_C$, em $x=56$ (o cruzamento de $c_B$ com $c_C$, em $x=40$, não chega a valer, pois nesse ponto $c_A$ já é menor que os dois). Logo:

$$C(x)=\begin{cases}x+22, & 0\le x\le 56\\[2mm] 0{,}5x+50, & 56< x\le 150\end{cases}$$

\textbf{Por que o UrbanMoto nunca é o mais barato:}

\begin{itemize}[nosep,leftmargin=*]
\item Para $0\le x<60$, $c_A(x)<c_B(x)$ (pois só se cruzam em $x=60$ e em $x=0$ já vale $c_A<c_B$); logo, nessa faixa, $c_B$ nunca é menor que $c_A$.
\item Para $x\ge 40$, $c_B(x)>c_C(x)$, pois a inclinação de $c_B$ (0,6) é maior que a de $c_C$ (0,5) e elas partem empatadas em $x=40$; logo, para todo $x\ge 40$, $c_B$ é sempre maior que $c_C$.
\item Juntando as duas conclusões: para $x<60$ há sempre alguém mais barato que $c_B$ (o próprio $c_A$, já que $x<60<... $ cobre inclusive $0\le x<40$), e para $x\ge 40$ há sempre $c_C$ mais barato que $c_B$. Como esses dois intervalos cobrem todo $[0,150]$, em nenhum ponto o UrbanMoto é a opção de menor custo — no máximo empata com outra opção (em $x=40$, empata com $c_C$, mas $c_A$ ainda é menor que ambos).
\end{itemize}

Portanto, o UrbanMoto é sempre dominado por MotoVá ou por RotaFácil e não aparece na função $C(x)$.

\bigskip\noindent\textbf{Q007 · múltipla escolha · fácil · aplicar · teórica · gabarito conferido}\par\nopagebreak
Uma oficina mecânica cobra por seus serviços de revisão um valor fixo de mão de obra, mais uma taxa que depende do número de horas de trabalho. Uma revisão que durou 2 horas custou R\$\,150,00, e outra revisão, no mesmo padrão de cobrança, durou 5 horas e custou R\$\,300,00. Sabendo que o custo total varia linearmente com o número de horas trabalhadas, qual seria o custo de uma revisão que durasse 8 horas?
\begin{enumerate}[label=(\alph*)]
\item R\$\,480,00
\item R\$\,400,00
\item R\$\,600,00
\item R\$\,450,00
\end{enumerate}

\textbf{Gabarito.} R\$\,450,00

\textbf{Resolução.} Como o custo total $C$ varia linearmente com o número de horas $x$, podemos escrever $C(x) = ax + b$, em que $a$ é a taxa cobrada por hora e $b$ é o valor fixo de mão de obra.

\textbf{Passo 1 — Determinar a taxa por hora ($a$).}
Usando os dois pares de valores dados, $(2, 150)$ e $(5, 300)$:
$$a = \frac{300 - 150}{5 - 2} = \frac{150}{3} = 50$$

Ou seja, cada hora de trabalho custa R\$\,50,00.

\textbf{Passo 2 — Determinar o valor fixo ($b$).}
Substituindo o ponto $(2, 150)$ em $C(x) = 50x + b$:
$$150 = 50 \cdot 2 + b \;\Rightarrow\; 150 = 100 + b \;\Rightarrow\; b = 50$$

Logo, o modelo do custo é:
$$C(x) = 50x + 50$$

\textbf{Passo 3 — Calcular o custo para $x = 8$ horas.}
$$C(8) = 50 \cdot 8 + 50 = 400 + 50 = 450$$

Portanto, a revisão de 8 horas custaria \textbf{R\$\,450,00}.

\bigskip\noindent\textbf{Q009 · múltipla escolha · média · criar · aplicada · gabarito conferido}\par\nopagebreak
Uma locadora de bicicletas cobra pelo aluguel um valor que combina uma taxa fixa de retirada com um valor proporcional ao número de horas de uso. Um cliente que utilizou a bicicleta por 2 horas pagou R\$\,18,00, enquanto outro cliente, que a utilizou por 5 horas, pagou R\$\,33,00.

Buscando atrair clientes que alugam por períodos mais longos, a locadora lançou um novo plano de cobrança, mantendo a mesma taxa fixa de retirada, mas alterando o valor cobrado por hora de uso. Nesse novo plano, um cliente que utiliza a bicicleta por 8 horas paga R\$\,10,00 a menos do que pagaria, nesse mesmo período de uso, pelo plano antigo.

Determine a lei que descreve o custo do novo plano em função do tempo de uso $t$ (em horas) e avalie para quais valores de $t$ esse novo plano é financeiramente mais vantajoso do que o plano antigo.
\begin{enumerate}[label=(\alph*)]
\item D(t) = 4,75t + 8; o novo plano é mais vantajoso para todo t > 0.
\item D(t) = 5t - 2; o novo plano é mais vantajoso para todo t > 0.
\item D(t) = 3,75t + 8; porém o novo plano só se torna mais vantajoso para t > 8 horas.
\item D(t) = 3,75t + 8; o novo plano é mais vantajoso, com custo estritamente menor que o antigo, para todo tempo de uso t > 0 (os dois planos coincidem apenas em t = 0).
\end{enumerate}

\textbf{Gabarito.} D(t) = 3,75t + 8; o novo plano é mais vantajoso (custo estritamente menor) para todo t > 0, sendo os planos equivalentes apenas em t = 0.

\textbf{Resolução.} \textbf{Passo 1 — Modelar o plano antigo.}
Como o custo combina uma parte fixa com uma parte proporcional ao tempo, o modelo é afim: $C(t) = a\,t + b$.

Dos dados: $C(2)=18$ e $C(5)=33$:
$2a+b=18$
$5a+b=33$

Subtraindo: $3a=15 \Rightarrow a=5$. Substituindo: $b = 18-2(5)=8$.

Logo, $C(t) = 5t+8$.

\textbf{Passo 2 — Modelar o plano novo.}
O novo plano mantém a mesma taxa fixa, então $D(t) = a't + 8$, com $a'$ a determinar.

O custo do plano antigo em $t=8$ é:
$C(8) = 5(8)+8 = 48$

O novo plano cobra R\$\,10,00 a menos nesse tempo:
$D(8) = 48-10 = 38$

Como $D(8) = 8a'+8$:
$8a'+8 = 38 \Rightarrow 8a' = 30 \Rightarrow a' = 3{,}75$

Portanto, $D(t) = 3{,}75t + 8$.

\textbf{Passo 3 — Comparar os dois planos.}
Calcule a diferença:
$C(t)-D(t) = (5t+8)-(3{,}75t+8) = 1{,}25t$

Como $t \geq 0$ (tempo de uso não pode ser negativo), temos $1{,}25t \geq 0$, com igualdade apenas em $t=0$.

Logo, $C(t) \geq D(t)$ para todo $t\geq 0$, sendo $D(t) < C(t)$ estritamente para $t>0$.

\textbf{Conclusão:} $D(t) = 3{,}75t+8$, e o novo plano é sempre financeiramente mais vantajoso (ou igual) ao antigo, tornando-se estritamente mais vantajoso para qualquer tempo de uso $t>0$.

\bigskip\noindent\textbf{Q011 · múltipla escolha · difícil · aplicar · teórica · gabarito conferido}\par\nopagebreak
Uma função afim $f(x) = ax + b$, com $a, b \in \mathbb{R}$ e $a \neq 0$, satisfaz simultaneamente as condições $f(2) = 7$ e $f(f(1)) = 13$. Determine a soma de todos os valores possíveis de $f(0)$.
\begin{enumerate}[label=(\alph*)]
\item 4
\item 1
\item 3
\item 6
\end{enumerate}

\textbf{Gabarito.} 4

\textbf{Resolução.} \textbf{Passo 1 — Traduzir as condições.}

Como $f(x) = ax+b$, a condição $f(2)=7$ fornece:
$$2a+b=7 \implies b = 7-2a$$

\textbf{Passo 2 — Expressar $f(f(1))$.}

Primeiro, $f(1) = a+b$. Usando $b=7-2a$:
$$f(1) = a + (7-2a) = 7-a$$

Agora aplicamos $f$ novamente:
$$f(f(1)) = f(7-a) = a(7-a)+b = 7a - a^2 + (7-2a)$$

\textbf{Passo 3 — Impor $f(f(1))=13$.}

$$7a - a^2 + 7 - 2a = 13$$
$$-a^2 + 5a + 7 = 13$$
$$-a^2 + 5a - 6 = 0 \implies a^2 - 5a + 6 = 0$$

\textbf{Passo 4 — Resolver a equação quadrática (análise de casos).}

$$a^2-5a+6=(a-2)(a-3)=0 \implies a=2 \text{ ou } a=3$$

Ambas as raízes são não nulas, logo ambas geram funções afins válidas — é preciso considerar os dois casos.

\textbf{Caso $a=2$:} $b = 7-2(2)=3$, logo $f(x)=2x+3$ e $f(0)=3$.

\textbf{Caso $a=3$:} $b = 7-2(3)=1$, logo $f(x)=3x+1$ e $f(0)=1$.

\textbf{Passo 5 — Somar os valores possíveis de $f(0)$.}

$$f(0)_{\text{soma}} = 3+1 = 4$$
\subsection*{EM13MAT303}
{\small\itshape Resolver e elaborar problemas envolvendo porcentagens em diversos contextos e sobre juros compostos, destacando o crescimento exponencial.}

\bigskip\noindent\textbf{Q064 · discursiva · difícil · criar · teórica · gabarito conferido}\par\nopagebreak
Uma grandeza positiva $C$ sofre, em cada uma de $n$ etapas sucessivas, um acréscimo percentual constante de taxa $i$ (com $0<i<1$): ao final de cada etapa, o valor da grandeza fica multiplicado por $(1+i)$ em relação ao valor obtido na etapa anterior. Esse é exatamente o modelo de capitalização por juros compostos.

\textbf{a)} Deduza, a partir da regra acima, a expressão geral $C_n$ do valor da grandeza após $n$ etapas, em função de $C$, $i$ e $n$. Em seguida, explique por que essa dependência de $C_n$ em relação a $n$ caracteriza um crescimento \textbf{exponencial}, e não um crescimento linear (do tipo somar, a cada etapa, um valor fixo igual a $iC$).

\textbf{b)} Um colega afirma que, se em duas etapas sucessivas a grandeza recebe acréscimos percentuais de taxas $i_1$ e $i_2$ (não necessariamente iguais), então a taxa percentual única $I$ equivalente às duas etapas combinadas é simplesmente $I = i_1+i_2$. Usando a expressão obtida no item (a), construa a fórmula correta de $I$ em função de $i_1$ e $i_2$, e mostre algebricamente que, sempre que $i_1>0$ e $i_2>0$, tem-se $I > i_1+i_2$.

\textbf{c)} Aplicando o modelo do item (a) com uma taxa constante de $i=20\%$ por etapa, determine o menor número inteiro de etapas $n$ necessário para que a grandeza mais que dobre de valor em relação ao valor inicial $C$.

\textbf{Gabarito.} (a) $C_n = C(1+i)^n$, exponencial pois $n$ está no expoente e a razão entre termos consecutivos é constante e igual a $(1+i)$; (b) $I = i_1+i_2+i_1i_2$, que excede $i_1+i_2$ em $i_1i_2>0$; (c) $n=4$.

\textbf{Resolução.} \textbf{Item (a) — construção da fórmula geral}

Após a etapa 1: $C_1 = C(1+i)$.

Após a etapa 2, o acréscimo incide sobre $C_1$ (e não sobre $C$): $C_2 = C_1(1+i) = C(1+i)^2$.

Repetindo o raciocínio (indução): a cada nova etapa o valor anterior é \textbf{multiplicado} pelo mesmo fator $(1+i)$. Logo,
$$C_n = C(1+i)^n.$$

Esse crescimento é exponencial porque a variável $n$ aparece no \textbf{expoente}, não como fator multiplicativo de um incremento fixo. Os valores $C, C(1+i), C(1+i)^2, \dots$ formam uma progressão \textbf{geométrica} de razão $(1+i)>1$: cada termo é obtido multiplicando o anterior pela mesma constante. Se o crescimento fosse linear, teríamos $C_n = C + n\cdot(iC)$, em que os incrementos absolutos $iC, iC, iC,\dots$ seriam sempre iguais — o que \textbf{não} ocorre aqui, pois o incremento absoluto da etapa $n$ é $iC(1+i)^{n-1}$, que cresce a cada etapa. É essa multiplicação repetida pelo mesmo fator (e não a soma repetida da mesma parcela) que define o caráter exponencial do processo.

\textbf{Item (b) — taxa equivalente a duas etapas}

Aplicando duas etapas sucessivas de taxas $i_1$ e $i_2$ à fórmula de (a):
$$C_2 = C(1+i_1)(1+i_2).$$

A taxa única equivalente $I$ deve satisfazer $C(1+I) = C(1+i_1)(1+i_2)$, ou seja,
$$1+I = (1+i_1)(1+i_2) = 1 + i_1 + i_2 + i_1 i_2 \;\Rightarrow\; I = i_1+i_2+i_1 i_2.$$

Comparando com a afirmação do colega:
$$I - (i_1+i_2) = i_1 i_2.$$

Se $i_1>0$ e $i_2>0$, então $i_1 i_2>0$, logo $I > i_1+i_2$ estritamente. O erro de simplesmente somar as taxas percentuais ignora o termo $i_1 i_2$, que representa o efeito de o segundo acréscimo incidir também sobre o acréscimo já obtido na primeira etapa — a essência da composição exponencial.

\textbf{Item (c) — menor $n$ para mais que dobrar}

Com $i=20\% = \dfrac{1}{5}$, temos $1+i = \dfrac{6}{5}$, e queremos o menor inteiro $n$ tal que
$$C\left(\frac{6}{5}\right)^n > 2C \;\Longleftrightarrow\; \left(\frac{6}{5}\right)^n > 2.$$

Testando valores inteiros:
\begin{itemize}[nosep,leftmargin=*]
\item $n=3$: $\left(\frac{6}{5}\right)^3 = \dfrac{216}{125} = 1{,}728 < 2$ (ainda não dobrou).
\item $n=4$: $\left(\frac{6}{5}\right)^4 = \dfrac{1296}{625} = 2{,}0736 > 2$ (já mais que dobrou).
\end{itemize}

De fato, resolvendo a igualdade $\left(\frac{6}{5}\right)^n = 2$ obtemos $n = \dfrac{\log 2}{\log(6/5)} \approx 3{,}80$, um valor não inteiro entre 3 e 4; como o crescimento é exponencial (estritamente crescente em $n$), o menor inteiro que satisfaz a desigualdade é o primeiro inteiro maior que $3{,}80$.

Portanto, $n = 4$ etapas.

\bigskip\noindent\textbf{Q072 · múltipla escolha · difícil · aplicar · aplicada · gabarito conferido}\par\nopagebreak
Marina aplicou R\$\,8.000,00 em um investimento que rende juros compostos a uma taxa fixa de 1,5\% ao mês sobre o saldo do mês anterior. Ela decidiu que só fará o resgate no primeiro mês em que o montante acumulado for igual ou superior ao dobro do capital investido. Usando as aproximações $\log 2 \approx 0,301$ e $\log 1,015 \approx 0,00647$, determine o número mínimo de meses completos necessários para que Marina possa efetuar esse resgate.
\begin{enumerate}[label=(\alph*)]
\item 48 meses
\item 47 meses
\item 46 meses
\item 67 meses
\end{enumerate}

\textbf{Gabarito.} 47 meses

\textbf{Resolução.} \textbf{Passo 1 — Modelar o montante.}

Como os juros são compostos, o montante após $n$ meses é
$$M(n) = 8000\cdot(1{,}015)^n.$$
Essa é uma função exponencial (razão constante $1{,}015>1$ entre meses consecutivos), ou seja, o crescimento é exponencial, não linear.

\textbf{Passo 2 — Montar a condição do problema.}

Marina resgata no primeiro mês em que $M(n)\ge 2\cdot 8000 = 16000$:
$$8000\cdot(1{,}015)^n \ge 16000 \;\;\Rightarrow\;\; (1{,}015)^n \ge 2.$$

\textbf{Passo 3 — Aplicar logaritmo.}

$$\log\big((1{,}015)^n\big) \ge \log 2 \;\;\Rightarrow\;\; n\cdot \log 1{,}015 \ge \log 2 \;\;\Rightarrow\;\; n \ge \dfrac{\log 2}{\log 1{,}015}.$$

\textbf{Passo 4 — Substituir os valores aproximados.}

$$n \ge \dfrac{0{,}301}{0{,}00647} \approx 46{,}52.$$

\textbf{Passo 5 — Determinar o menor número inteiro de meses.}

Como $n$ representa um número de meses completos, é preciso o menor inteiro que satisfaça $n\ge 46{,}52$, ou seja, $n = 47$.

Verificação:
\begin{itemize}[nosep,leftmargin=*]
\item Para $n=46$: $46\times 0{,}00647 = 0{,}29762 < 0{,}301$, logo $(1{,}015)^{46}<2$ — ainda não dobrou.
\item Para $n=47$: $47\times 0{,}00647 = 0{,}30409 > 0{,}301$, logo $(1{,}015)^{47}>2$ — já dobrou.
\end{itemize}

Portanto, o menor número de meses é \textbf{47}.

\bigskip\noindent\textbf{Q077 · discursiva · média · criar · aplicada · gabarito conferido}\par\nopagebreak
Uma fintech oferece uma aplicação financeira em que o capital investido é corrigido mensalmente por juros compostos, a uma taxa fixa, sem aportes ou resgates. Marina investiu R\$\,6.000,00 nessa aplicação, com taxa de 4\% ao mês.

a) Elabore a expressão algébrica da função $M(t)$, que fornece o montante (em reais) acumulado por Marina após $t$ meses de aplicação.

b) Utilizando a função elaborada, calcule o montante ao final do 10º mês, arredondando o resultado para duas casas decimais.

c) Determine o menor número inteiro de meses necessário para que o montante ultrapasse o dobro do capital investido. Em seguida, explique por que esse tempo é bem menor do que o que seria necessário caso o capital rendesse mensalmente um valor fixo de R\$\,240,00 (4\% do capital inicial, sem incidência de juros sobre juros) — evidenciando, com essa comparação, o caráter exponencial (e não linear) do crescimento sob juros compostos.

\textbf{Gabarito.} a) $M(t) = 6000\cdot(1,04)^t$; b) $M(10) \approx R\$\,8.881{,}47$; c) menor número inteiro de meses: $t = 18$ (contra 25 meses no regime linear equivalente, evidenciando o crescimento exponencial).

\textbf{Resolução.} \textbf{a) Modelagem da função}

Como o capital é corrigido mensalmente por uma taxa percentual constante $i = 4\% = 0,04$ aplicada sobre o montante do mês anterior (e não sobre um valor fixo), cada mês o montante é multiplicado pelo mesmo fator $(1+i)$. Isso caracteriza um crescimento exponencial:

$$M(t) = 6000\cdot(1+0,04)^t = 6000\cdot(1,04)^t = 6000\cdot\left(\dfrac{26}{25}\right)^t$$

em que $t$ é o número de meses.

\textbf{b) Montante após 10 meses}

$$M(10) = 6000\cdot(1,04)^{10} = 6000\cdot 1{,}480244\ldots \approx R\$\,8.881{,}47$$

\textbf{c) Tempo para dobrar o capital}

Queremos o menor $t$ inteiro tal que $M(t) > 12000$ (o dobro de 6000):

$$6000\cdot(1,04)^t > 12000 \;\Rightarrow\; (1,04)^t > 2$$

Aplicando logaritmo em ambos os lados:

$$t > \dfrac{\log 2}{\log 1,04} \approx \dfrac{0,6931}{0,03922} \approx 17,67$$

Como $t$ deve ser um número inteiro de meses, o menor valor que satisfaz a desigualdade é $t = 18$. Conferindo:

\begin{itemize}[nosep,leftmargin=*]
\item $M(17) = 6000\cdot(1,04)^{17} \approx R\$\,11.687{,}40 < 12.000$
\item $M(18) = 6000\cdot(1,04)^{18} \approx R\$\,12.154{,}90 > 12.000$
\end{itemize}

Logo, \textbf{18 meses} é o menor número inteiro de meses necessário.

\textbf{Comparação com crescimento linear}

Se, em vez de juros compostos, o capital rendesse sempre um valor fixo de $0,04\cdot 6000 = R\$\,240,00$ por mês (juros simples, crescimento linear), o montante seria $M_{lin}(t) = 6000 + 240t$, e o tempo para dobrar o capital seria:

$$6000 + 240t = 12000 \;\Rightarrow\; t = \dfrac{6000}{240} = 25 \text{ meses}$$

Assim, no regime linear seriam necessários 25 meses, contra apenas 18 meses nos juros compostos. Essa diferença ocorre porque, nos juros compostos, os 4\% incidem sobre um montante que cresce mês a mês (juros sobre juros), fazendo os incrementos mensais aumentarem progressivamente — comportamento típico de uma função exponencial, e não de uma função linear, em que o incremento mensal é sempre o mesmo valor fixo.

\bigskip\noindent\textbf{Q085 · múltipla escolha · média · aplicar · teórica · gabarito conferido}\par\nopagebreak
Um capital $C_0$ é aplicado a juros compostos, com taxa fixa $i$ por período, de modo que o montante após $n$ períodos é dado por $M(n) = C_0\cdot(1+i)^n$. Sabe-se que, ao final de 3 períodos, o montante corresponde a 8 vezes o capital inicial, isto é, $M(3) = 8\,C_0$. Mantendo-se a mesma taxa $i$ por todo o tempo, qual é o montante $M(6)$, em função de $C_0$?
\begin{enumerate}[label=(\alph*)]
\item $64\,C_0$
\item $16\,C_0$
\item $32\,C_0$
\item $512\,C_0$
\end{enumerate}

\textbf{Gabarito.} 64 C0

\textbf{Resolução.} \textbf{Passo 1 — Traduzir o dado em uma equação.}

Como $M(n) = C_0\cdot(1+i)^n$, a condição $M(3) = 8\,C_0$ fornece
$$C_0\cdot(1+i)^3 = 8\,C_0 \;\Rightarrow\; (1+i)^3 = 8.$$

\textbf{Passo 2 — Usar a propriedade de potência para obter $(1+i)^6$ sem precisar isolar $i$.}

Como $6 = 3\cdot 2$, escrevemos
$$(1+i)^6 = \left[(1+i)^3\right]^2 = 8^2 = 64.$$

Esse passo é o núcleo do raciocínio: dobrar o número de períodos \textbf{eleva ao quadrado} o fator de crescimento acumulado (comportamento exponencial), e não o dobra (comportamento que seria linear).

\textbf{Passo 3 — Calcular o montante pedido.}

$$M(6) = C_0\cdot(1+i)^6 = C_0\cdot 64 = 64\,C_0.$$

Portanto, após 6 períodos o montante é $64\,C_0$, ou seja, 64 vezes o capital inicial — muito mais que o dobro dos 8 vezes obtidos em 3 períodos, evidenciando o crescimento exponencial (e não linear) do capital sob juros compostos.
\subsection*{EM13MAT304}
{\small\itshape Resolver e elaborar problemas com funções exponenciais nos quais é necessário compreender e interpretar a variação das grandezas envolvidas, em contextos como o da Matemática Financeira e o do crescimento de seres vivos microscópicos, entre outros.}

\bigskip\noindent\textbf{Q046 · discursiva · média · criar · teórica · gabarito conferido}\par\nopagebreak
Um biólogo monitora o crescimento de uma colônia de bactérias em um meio de cultura controlado. No início da observação ($t=0$), a colônia possui 500 bactérias. Verificou-se experimentalmente que essa população duplica a cada 3 horas, mantendo esse padrão de crescimento durante todo o experimento, que dura pelo menos 24 horas.

a) Elabore a lei da função exponencial $N(t)$ que fornece o número de bactérias em função do tempo $t$, em horas, decorrido desde o início da observação.

b) Utilizando a função obtida, determine o número de bactérias presentes após 9 horas de observação.

c) Um estudante afirmou que, como a população duplica a cada 3 horas, o crescimento percentual da colônia a cada hora deve ser de $\dfrac{100\%}{3}\approx 33{,}3\%$. Mostre que essa afirmação está incorreta, determinando a taxa de crescimento percentual horária real da colônia (isto é, encontre $i$ tal que $N(t) = 500\cdot(1+i)^t$ para todo $t$), e explique por que o raciocínio do estudante falha, considerando o comportamento multiplicativo da função exponencial.

d) Determine depois de quantas horas a população da colônia atingirá 32000 bactérias.

\textbf{Gabarito.} a) $N(t) = 500\cdot 2^{t/3}$; b) 4000 bactérias; c) $i = 2^{1/3}-1 \approx 26\%$ ao hora — o raciocínio do estudante é inválido pois o crescimento exponencial é multiplicativo, não pode ser dividido linearmente pelo tempo; d) 18 horas.

\textbf{Resolução.} \textbf{a) Construção da lei da função}

Como a população duplica a cada 3 horas, o número de bactérias é multiplicado por 2 sempre que $t$ aumenta em 3 unidades. Isso caracteriza uma função exponencial da forma
$$N(t) = N_0 \cdot 2^{t/3}$$
com $N_0 = 500$ (população inicial). Logo,
$$N(t) = 500\cdot 2^{t/3}.$$

Verificação: $N(0)=500$, $N(3)=500\cdot2=1000$, $N(6)=500\cdot4=2000$ — de fato duplica a cada 3 h.

\textbf{b) Valor após 9 horas}

$$N(9) = 500\cdot 2^{9/3} = 500\cdot 2^{3} = 500\cdot 8 = 4000.$$

Após 9 horas, há \textbf{4000 bactérias}.

\textbf{c) Taxa de crescimento horária real}

Queremos escrever a mesma função na forma $N(t) = 500\cdot(1+i)^t$. Para que essa forma coincida com $500\cdot 2^{t/3}$ para todo $t$, as bases devem ser iguais:
$$(1+i)^t = \left(2^{1/3}\right)^t \;\Rightarrow\; 1+i = 2^{1/3}.$$
Logo,
$$i = 2^{1/3} - 1 \approx 1{,}2599 - 1 = 0{,}2599,$$
ou seja, a colônia cresce aproximadamente \textbf{26\% a cada hora}, e não 33,3\%.

O erro do estudante está em tratar o crescimento como se fosse \textbf{distribuído de forma aditiva/linear} ao longo das 3 horas (dividindo os 100\% de aumento por 3). Mas o crescimento exponencial é \textbf{multiplicativo}: em cada hora a população é multiplicada pelo mesmo fator $2^{1/3}$, e não recebe uma fração fixa do aumento total. Prova disso: se a taxa fosse realmente $33{,}3\%$ ao h, após 3 horas o fator acumulado seria $(1{,}333)^3 \approx 2{,}37$, e não exatamente $2$, como exige o enunciado.

\textbf{d) Tempo para atingir 32000 bactérias}

$$500\cdot 2^{t/3} = 32000 \;\Rightarrow\; 2^{t/3} = 64 = 2^{6} \;\Rightarrow\; \frac{t}{3} = 6 \;\Rightarrow\; t = 18.$$

A população atingirá 32000 bactérias após \textbf{18 horas}.

\bigskip\noindent\textbf{Q054 · múltipla escolha · média · aplicar · aplicada · gabarito conferido}\par\nopagebreak
Um biólogo monitora o crescimento de uma cultura de bactérias em laboratório. No início da observação (t = 0), há 400 bactérias, e a cultura cresce de modo que a população dobra a cada 6 horas. A população N(t), em número de bactérias, t horas após o início da observação, é dada por N(t) = 400 · 2\textasciicircum{}(t/6). Considerando essa lei de crescimento, qual das afirmações a seguir compara corretamente o crescimento da população entre t = 0h e t = 6h com o crescimento entre t = 12h e t = 18h?
\begin{enumerate}[label=(\alph*)]
\item Em ambos os intervalos a população aumenta exatamente 400 bactérias, já que a taxa de crescimento de 2 vezes a cada 6 horas corresponde a um acréscimo fixo de indivíduos por hora.
\item A população dobra apenas uma vez em todo o período analisado, atingindo 1600 bactérias em 18 horas, pois o enunciado indica que o dobro ocorre a cada 12 horas, não a cada 6.
\item Em ambos os intervalos a população dobra (fator multiplicativo 2), mas o aumento absoluto de bactérias é maior no segundo intervalo: 400 entre 0h e 6h, contra 1600 entre 12h e 18h, pois a mesma taxa de crescimento multiplica quantidades cada vez maiores.
\item No segundo intervalo (12h a 18h) a população cresce mais do que o dobro, pois quanto mais tempo passa, maior fica o fator de multiplicação aplicado a cada período de 6 horas.
\end{enumerate}

\textbf{Gabarito.} Em ambos os intervalos a população dobra (fator 2), mas o aumento absoluto é de 400 bactérias no primeiro intervalo e de 1600 bactérias no segundo, pois a mesma razão de crescimento multiplica quantidades cada vez maiores.

\textbf{Resolução.} \textbf{Passo 1 — Calcular a população nos quatro instantes pedidos.}

$N(0) = 400 \cdot 2^{0/6} = 400 \cdot 1 = 400$

$N(6) = 400 \cdot 2^{6/6} = 400 \cdot 2 = 800$

$N(12) = 400 \cdot 2^{12/6} = 400 \cdot 4 = 1600$

$N(18) = 400 \cdot 2^{18/6} = 400 \cdot 8 = 3200$

\textbf{Passo 2 — Analisar o fator multiplicativo em cada intervalo de 6 horas.}

Entre 0h e 6h: $\dfrac{N(6)}{N(0)} = \dfrac{800}{400} = 2$.

Entre 12h e 18h: $\dfrac{N(18)}{N(12)} = \dfrac{3200}{1600} = 2$.

Como a função é exponencial com base fixa, todo intervalo de 6 horas multiplica a população pelo mesmo fator (2), \textbf{independentemente do instante em que ele começa} — essa é a característica que define o crescimento exponencial (taxa percentual constante).

\textbf{Passo 3 — Analisar o aumento absoluto em cada intervalo.}

Entre 0h e 6h: $N(6) - N(0) = 800 - 400 = 400$ bactérias.

Entre 12h e 18h: $N(18) - N(12) = 3200 - 1600 = 1600$ bactérias.

Embora o fator multiplicativo seja o mesmo (2), o aumento absoluto é diferente: como o segundo intervalo parte de uma população maior, a mesma multiplicação por 2 produz um acréscimo absoluto quatro vezes maior (1600 contra 400).

\textbf{Conclusão:} em ambos os intervalos a população dobra (razão constante = 2), mas o número de bactérias adicionadas é bem maior no segundo intervalo, pois o crescimento exponencial aplica a mesma taxa percentual sobre quantidades cada vez maiores.

\bigskip\noindent\textbf{Q067 · múltipla escolha · fácil · analisar · teórica · gabarito conferido}\par\nopagebreak
Uma colônia de bactérias cresce segundo a lei $N(t) = 100 \cdot 3^{t/4}$, em que $N(t)$ é o número de bactérias e $t$ é o tempo em horas, contado a partir do instante em que a colônia tinha 100 indivíduos. Um estudante quer saber quanto tempo leva para a população \textbf{triplicar} de tamanho, e resolve investigar isso em dois momentos diferentes: partindo de $t=0$ e partindo de $t=4$ horas (quando a população já triplicou uma vez). Ele calcula: (i) o tempo $\Delta t_1$ necessário para que $N(t)$ passe de $N(0)$ para $3\cdot N(0)$; (ii) o tempo $\Delta t_2$ necessário para que $N(t)$ passe de $N(4)$ para $3\cdot N(4)$. Assinale a alternativa que descreve corretamente a relação entre $\Delta t_1$ e $\Delta t_2$, e o motivo dessa relação.
\begin{enumerate}[label=(\alph*)]
\item $\Delta t_2 < \Delta t_1$, pois quanto maior a população, mais rápido ela cresce, então o tempo para triplicar diminui a cada rodada.
\item Não é possível comparar $\Delta t_1$ e $\Delta t_2$ sem conhecer o valor de $N(0)$, já que o tempo de triplicação depende do tamanho inicial da colônia.
\item $\Delta t_2 > \Delta t_1$, pois, como a população em $t=4$ já é maior, seria necessário um tempo maior para ela triplicar novamente.
\item $\Delta t_1 = \Delta t_2 = 4$ horas, pois na função exponencial a razão entre valores separados por um mesmo intervalo de tempo é sempre a mesma, independentemente do instante inicial ou do tamanho já atingido pela população.
\end{enumerate}

\textbf{Gabarito.} $\Delta$t1 = $\Delta$t2 = 4 horas — o tempo de triplicação é constante e independente do instante inicial, pois a variação relativa de uma função exponencial depende apenas do intervalo de tempo decorrido, não do valor absoluto da grandeza.

\textbf{Resolução.} \textbf{Passo 1 — Calcular $\Delta t_1$ (triplicação a partir de $t=0$):}

Queremos $t$ tal que $N(t) = 3\cdot N(0)$, ou seja:
$$100\cdot 3^{t/4} = 3\cdot 100$$
$$3^{t/4} = 3^1 \implies \frac{t}{4}=1 \implies t = 4$$
Logo $\Delta t_1 = 4$ horas.

\textbf{Passo 2 — Calcular $\Delta t_2$ (triplicação a partir de $t=4$):}

Primeiro, $N(4) = 100\cdot 3^{4/4} = 300$. Queremos $t'$ tal que $N(4+t') = 3\cdot N(4) = 900$:
$$100\cdot 3^{(4+t')/4} = 900$$
$$3^{(4+t')/4} = 9 = 3^2 \implies \frac{4+t'}{4} = 2 \implies 4+t' = 8 \implies t' = 4$$
Logo $\Delta t_2 = 4$ horas.

\textbf{Passo 3 — Interpretar o resultado:}

Como $N(t) = N(t_0)\cdot 3^{(t-t_0)/4}$ para qualquer instante inicial $t_0$, a razão $\dfrac{N(t_0+\Delta t)}{N(t_0)} = 3^{\Delta t/4}$ depende apenas de $\Delta t$ (a variação do tempo), e não do valor de $t_0$ nem do valor absoluto de $N(t_0)$. Por isso, o tempo necessário para a população \textbf{triplicar} é sempre o mesmo (4 horas), não importa em que instante o processo começa. Esse é um traço característico do crescimento exponencial: variações \textbf{relativas} (multiplicativas) de uma grandeza ocorrem em intervalos de tempo constantes, mesmo que as variações \textbf{absolutas} (em número de bactérias) sejam cada vez maiores.

Portanto, $\Delta t_1 = \Delta t_2 = 4$ horas, pois o tempo de triplicação depende só da razão de crescimento (base 3 na exponencial), não do instante inicial nem do tamanho já atingido pela população.

\bigskip\noindent\textbf{Q080 · discursiva · difícil · analisar · aplicada · gabarito conferido}\par\nopagebreak
Um laboratório de microbiologia monitora duas culturas de bactérias mantidas em condições ideais de nutrientes e temperatura, iniciando a contagem no instante $t=0$ (em horas).

\begin{itemize}[nosep,leftmargin=*]
\item A \textbf{Cultura A} parte de 500 bactérias e duplica sua população a cada 3 horas, sendo modelada por $P_A(t) = 500 \cdot 2^{t/3}$.
\item A \textbf{Cultura B} parte de 8000 bactérias e duplica sua população a cada 6 horas, sendo modelada por $P_B(t) = 8000 \cdot 2^{t/6}$.
\end{itemize}

Embora a Cultura B comece com uma população 16 vezes maior que a Cultura A, o pesquisador responsável afirma que, depois de certo tempo, a Cultura A ultrapassará a Cultura B e permanecerá com mais bactérias dali em diante, para sempre.

a) Determine o instante $t$ (em horas) em que as duas culturas apresentam exatamente a mesma quantidade de bactérias.

b) Sem calcular os valores das populações em instantes específicos, explique — comparando os tempos de duplicação (ou as taxas de crescimento) das duas culturas — por que, a partir do instante encontrado no item (a), a Cultura A permanecerá \textbf{para sempre} com mais bactérias que a Cultura B.

c) Confirme numericamente a afirmação do pesquisador calculando $P_A(30)$ e $P_B(30)$, isto é, as populações das duas culturas 30 horas após o início do experimento.

\textbf{Gabarito.} a) $t = 24$ horas. b) Como a taxa de crescimento relativo de A ($\ln 2/3$) é maior que a de B ($\ln 2/6$), a razão $P_A(t)/P_B(t) = \frac{1}{16}\cdot 2^{t/6}$ é uma exponencial crescente que vale 1 em $t=24$h e cresce indefinidamente para $t>24$h; logo A ultrapassa B nesse instante e permanece maior para sempre. c) $P_A(30)=512000$ e $P_B(30)=256000$, confirmando que $P_A(30) > P_B(30)$.

\textbf{Resolução.} \textbf{a) Igualando as populações}

$$500 \cdot 2^{t/3} = 8000 \cdot 2^{t/6}$$

Dividindo ambos os lados por $500 \cdot 2^{t/6}$:

$$2^{t/3 - t/6} = \frac{8000}{500} = 16$$

Como $\frac{t}{3} - \frac{t}{6} = \frac{t}{6}$, temos:

$$2^{t/6} = 16 = 2^4 \;\Rightarrow\; \frac{t}{6} = 4 \;\Rightarrow\; t = 24 \text{ horas}$$

Nesse instante, ambas as culturas têm $P_A(24)=500\cdot 2^8 = 128000$ bactérias, o mesmo valor de $P_B(24) = 8000\cdot 2^4 = 128000$.

\textbf{b) Interpretando as taxas de crescimento}

O tempo de duplicação da Cultura A é 3 horas, enquanto o da Cultura B é 6 horas — ou seja, A dobra sua população \textbf{duas vezes mais rápido} que B. Isso significa que a taxa exponencial de crescimento de A, $\frac{\ln 2}{3}$, é maior que a de B, $\frac{\ln 2}{6}$.

Para entender o comportamento da razão entre as populações, escrevemos:

$$\frac{P_A(t)}{P_B(t)} = \frac{500}{8000}\cdot 2^{t/3 - t/6} = \frac{1}{16}\cdot 2^{t/6}$$

Essa razão é uma função exponencial crescente de $t$ (pois a base $2>1$ e o expoente $t/6$ cresce com $t$). Em $t=24$ h, a razão vale exatamente 1 (populações iguais). Como a razão é estritamente crescente para todo $t$, para $t>24$ ela é sempre maior que 1, e continua aumentando indefinidamente conforme $t$ cresce. Logo, não é apenas que A ultrapassa B em $t=24$: como a taxa relativa de crescimento de A é permanentemente maior que a de B, a diferença entre as populações só tende a aumentar, e A permanece maior que B \textbf{para sempre}.

\textbf{c) Confirmação numérica em $t=30$}

$$P_A(30) = 500 \cdot 2^{30/3} = 500 \cdot 2^{10} = 500 \cdot 1024 = 512000$$

$$P_B(30) = 8000 \cdot 2^{30/6} = 8000 \cdot 2^{5} = 8000 \cdot 32 = 256000$$

De fato, $P_A(30) = 512000 > 256000 = P_B(30)$, confirmando que 6 horas após o cruzamento (que ocorreu em $t=24$h), a Cultura A já tem o dobro de bactérias da Cultura B, evidenciando o comportamento descrito no item (b).

\bigskip\noindent\textbf{Q088 · múltipla escolha · difícil · aplicar · teórica · gabarito conferido}\par\nopagebreak
Um biólogo monitora duas culturas de microrganismos mantidas em condições ideais de crescimento.

\begin{itemize}[nosep,leftmargin=*]
\item A cultura A começa com uma quantidade $A_0$ de microrganismos, e essa quantidade dobra a cada 2 horas.
\item A cultura B começa com o quádruplo da quantidade inicial da cultura A, ou seja, com $4A_0$ microrganismos, e essa quantidade dobra a cada 6 horas.
\end{itemize}

Sabendo que as duas populações crescem de forma exponencial desde o início da observação, determine depois de quantas horas as duas culturas terão exatamente a mesma quantidade de microrganismos.
\begin{enumerate}[label=(\alph*)]
\item 4 horas
\item 6 horas
\item 3 horas
\item 8 horas
\end{enumerate}

\textbf{Gabarito.} 6 horas

\textbf{Resolução.} \textbf{Modelando as populações}

Como a cultura A dobra a cada 2 horas, sua população no instante $t$ (em horas) é
$$P_A(t) = A_0 \cdot 2^{t/2}.$$

Como a cultura B começa com $4A_0$ e dobra a cada 6 horas,
$$P_B(t) = 4A_0 \cdot 2^{t/6}.$$

\textbf{Interpretando a situação}

No instante $t=0$, $P_B(0)=4A_0 > A_0 = P_A(0)$: a cultura B começa muito maior. Porém A cresce mais rapidamente (seu tempo de duplicação é menor), então a diferença entre as populações vai diminuindo até que, em algum instante, elas se igualam — e depois disso A ultrapassa B. É esse instante que queremos encontrar.

\textbf{Montando a equação}

$$A_0 \cdot 2^{t/2} = 4A_0 \cdot 2^{t/6}.$$

Como $A_0 > 0$, podemos cancelá-lo:
$$2^{t/2} = 4 \cdot 2^{t/6}.$$

Escrevendo $4 = 2^2$ e usando a propriedade de produto de potências de mesma base:
$$2^{t/2} = 2^{2 + t/6}.$$

\textbf{Igualando os expoentes}

Como as bases são iguais (e a função exponencial é injetora), os expoentes devem ser iguais:
$$\frac{t}{2} = 2 + \frac{t}{6}.$$

Multiplicando tudo por 6:
$$3t = 12 + t \implies 2t = 12 \implies t = 6.$$

\textbf{Verificação}

$P_A(6) = A_0 \cdot 2^{3} = 8A_0$.

$P_B(6) = 4A_0 \cdot 2^{1} = 8A_0$.

As populações coincidem em $8A_0$, confirmando que $t = 6$ horas é o instante correto.
\subsection*{EM13MAT305}
{\small\itshape Resolver e elaborar problemas com funções logarítmicas nos quais é necessário compreender e interpretar a variação das grandezas envolvidas, em contextos como os de abalos sísmicos, pH, radioatividade, Matemática Financeira, entre outros.}

\bigskip\noindent\textbf{Q049 · discursiva · difícil · analisar · teórica · gabarito conferido}\par\nopagebreak
Cientistas utilizam escalas logarítmicas para descrever a magnitude de fenômenos naturais cuja intensidade varia em muitas ordens de grandeza, como ocorre com a energia liberada em abalos sísmicos. Um modelo simplificado dessa relação define a magnitude $M$ de um terremoto, em função da energia $E$ (em joules) por ele liberada, pela expressão

$$M(E) = \frac{2}{3}\left(\log_{10}(E) - 4\right), \quad E > 0.$$

a) Um terremoto liberou energia $E_1 = 10^{13}$ J. Calcule sua magnitude $M_1$.

b) Um segundo terremoto ocorreu na mesma região, com magnitude exatamente uma unidade maior que a do primeiro, ou seja, $M_2 = M_1 + 1$. Determine, em função de $E_1$, a energia $E_2$ liberada por esse segundo terremoto, calcule a razão $E_2/E_1$ e mostre que essa razão não depende do valor de $E_1$ — isto é, que qualquer terremoto cuja magnitude seja uma unidade maior que a de outro libera sempre o mesmo número de vezes mais energia, independentemente do nível de energia inicial.

\textbf{Gabarito.} M1 = 6; E2/E1 = 10\textasciicircum{}(3/2) = 10$\surd$10 $\approx$ 31,62, valor constante e independente de E1 (todo acréscimo de 1 unidade na magnitude multiplica a energia por esse mesmo fator).

\textbf{Resolução.} \textbf{Passo 1 — Calcular $M_1$}

Com $E_1 = 10^{13}$ J, temos $\log_{10}(E_1) = 13$. Substituindo na fórmula:

$$M_1 = \frac{2}{3}\left(13 - 4\right) = \frac{2}{3}\cdot 9 = 6.$$

\textbf{Passo 2 — Inverter a função para expressar $E$ em termos de $M$}

Da definição:
$$M = \frac{2}{3}\left(\log_{10}E - 4\right) \;\Rightarrow\; \frac{3}{2}M = \log_{10}E - 4 \;\Rightarrow\; \log_{10}E = \frac{3}{2}M + 4.$$

Logo,
$$E = 10^{\,4 + \frac{3}{2}M} = 10^{4}\cdot 10^{\frac{3}{2}M}.$$

(Conferência: para $M_1=6$, $E_1 = 10^4\cdot10^{9} = 10^{13}$ J, coincidindo com o dado do item a.)

\textbf{Passo 3 — Determinar $E_2$ e a razão $E_2/E_1$}

Como $M_2 = M_1+1 = 7$:
$$E_2 = 10^{4}\cdot 10^{\frac{3}{2}\cdot 7} = 10^{4+10{,}5} = 10^{14{,}5}.$$

A razão entre as energias é:
$$\frac{E_2}{E_1} = \frac{10^{4}\cdot10^{\frac{3}{2}M_2}}{10^{4}\cdot10^{\frac{3}{2}M_1}} = 10^{\frac{3}{2}(M_2-M_1)} = 10^{\frac{3}{2}(1)} = 10^{3/2} = 10\sqrt{10} \approx 31{,}62.$$

\textbf{Passo 4 — Mostrar que a razão independe de $E_1$}

Para qualquer terremoto de magnitude $M$ (energia $E$) e outro de magnitude $M+1$ (energia $E'$), repetindo o cálculo anterior:
$$\frac{E'}{E} = 10^{\frac{3}{2}\left[(M+1)-M\right]} = 10^{3/2},$$
que é uma constante que \textbf{não depende de $M$ nem de $E$}. Ou seja, aumentar a magnitude em exatamente 1 unidade sempre multiplica a energia liberada pelo mesmo fator $10\sqrt{10}\approx 31{,}62$, seja o terremoto inicial fraco ou muito forte — a variação da grandeza física (energia) é sempre multiplicativa e constante quando a variação da grandeza logarítmica (magnitude) é aditiva e constante.

\bigskip\noindent\textbf{Q057 · múltipla escolha · difícil · analisar · aplicada · gabarito conferido}\par\nopagebreak
A magnitude $M$ de um terremoto na escala Richter relaciona-se com a energia $E$ (em joules) liberada pelo abalo por meio da expressão $M = \dfrac{2}{3}\log_{10}\left(\dfrac{E}{E_0}\right)$, em que $E_0 = 10^{4,4}$ J é uma constante de referência. Uma cidade foi atingida por um terremoto de magnitude 5,0. Meses depois, a mesma região sofreu outro abalo, de magnitude 6,5. Aproximadamente quantas vezes maior foi a energia liberada pelo segundo terremoto em comparação com o primeiro?
\begin{enumerate}[label=(\alph*)]
\item Aproximadamente 10 vezes maior, pois $10^{1} = 10$.
\item Aproximadamente 32 vezes maior, pois $10^{1,5} \approx 31{,}6$.
\item Aproximadamente 178 vezes maior.
\item Aproximadamente 1,3 vezes maior (30\% a mais), pois $6{,}5 \div 5{,}0 = 1{,}3$.
\end{enumerate}

\textbf{Gabarito.} Aproximadamente 178 vezes (10\textasciicircum{}2,25)

\textbf{Resolução.} \textbf{Passo 1 — Isolar a energia $E$ em função de $M$.}

Da expressão $M = \dfrac{2}{3}\log_{10}\left(\dfrac{E}{E_0}\right)$, multiplicamos ambos os lados por $\dfrac{3}{2}$:

$$\dfrac{3}{2}M = \log_{10}\left(\dfrac{E}{E_0}\right)$$

Aplicando a definição de logaritmo:

$$\dfrac{E}{E_0} = 10^{\frac{3}{2}M} \quad\Rightarrow\quad E = E_0\cdot 10^{\frac{3}{2}M}$$

Isso mostra que a energia \textbf{não varia linearmente} com $M$: cada acréscimo de magnitude corresponde a uma multiplicação da energia por uma potência de 10.

\textbf{Passo 2 — Calcular a energia para cada magnitude.}

Para $M_1 = 5{,}0$: $E_1 = E_0\cdot 10^{\frac{3}{2}\cdot 5} = E_0\cdot 10^{7{,}5}$

Para $M_2 = 6{,}5$: $E_2 = E_0\cdot 10^{\frac{3}{2}\cdot 6{,}5} = E_0\cdot 10^{9{,}75}$

\textbf{Passo 3 — Calcular a razão entre as energias.}

$$\dfrac{E_2}{E_1} = \dfrac{E_0\cdot 10^{9{,}75}}{E_0\cdot 10^{7{,}5}} = 10^{9{,}75-7{,}5} = 10^{2{,}25}$$

Note que o expoente $2{,}25$ é exatamente $\dfrac{3}{2}\cdot(M_2-M_1) = \dfrac{3}{2}\cdot 1{,}5 = \dfrac{9}{4}$.

\textbf{Passo 4 — Calcular o valor numérico.}

$$10^{2{,}25} = 10^2 \cdot 10^{0{,}25} \approx 100 \times 1{,}778 \approx 177{,}8$$

\textbf{Conclusão:} o segundo terremoto liberou aproximadamente \textbf{178 vezes} mais energia que o primeiro, mesmo a magnitude tendo aumentado apenas 1,5 unidade — evidenciando o crescimento exponencial da energia em relação à escala logarítmica de magnitude.

\bigskip\noindent\textbf{Q062 · discursiva · média · analisar · aplicada · gabarito conferido}\par\nopagebreak
A magnitude $M$ de um abalo sísmico, medida na escala Richter, está relacionada à amplitude $A$ das ondas sísmicas registradas por um sismógrafo pela função $M = \log_{10}\left(\dfrac{A}{A_0}\right)$, em que $A_0$ é uma amplitude de referência constante (a mesma para todos os registros de um mesmo aparelho).

Uma cidade foi atingida por um terremoto de magnitude 4,0. Semanas depois, um segundo terremoto, de magnitude 6,0, atingiu a mesma região e foi registrado pelo mesmo sismógrafo.

a) Determine quantas vezes a amplitude das ondas sísmicas do terremoto de magnitude 6,0 foi maior do que a amplitude das ondas do terremoto de magnitude 4,0.

b) Um morador da cidade afirmou: "Como a diferença de magnitude foi de apenas 2 pontos, numa escala que costuma ir até 10, o segundo terremoto deve ter sido só um pouco mais forte, talvez uns 20\% mais forte que o primeiro." Usando o valor obtido no item (a), explique por que essa interpretação está incorreta, destacando a natureza logarítmica (não linear) da escala Richter.

\textbf{Gabarito.} a) A amplitude do terremoto de magnitude 6,0 foi 100 vezes maior que a do terremoto de magnitude 4,0. b) A interpretação do morador está errada porque a escala Richter é logarítmica: cada unidade a mais em $M$ multiplica a amplitude por 10 (não soma uma fração fixa). Uma diferença de 2 pontos corresponde a uma amplitude 100 vezes maior, não 20\% maior.

\textbf{Resolução.} \textbf{a) Cálculo da razão entre as amplitudes}

Para os dois terremotos, temos:

$M_1 = \log_{10}\left(\dfrac{A_1}{A_0}\right) = 4{,}0$

$M_2 = \log_{10}\left(\dfrac{A_2}{A_0}\right) = 6{,}0$

Subtraindo as duas equações e usando a propriedade do logaritmo do quociente:

$M_2 - M_1 = \log_{10}\left(\dfrac{A_2}{A_0}\right) - \log_{10}\left(\dfrac{A_1}{A_0}\right) = \log_{10}\left(\dfrac{A_2/A_0}{A_1/A_0}\right) = \log_{10}\left(\dfrac{A_2}{A_1}\right)$

Logo:

$\log_{10}\left(\dfrac{A_2}{A_1}\right) = 6{,}0 - 4{,}0 = 2{,}0$

Aplicando a definição de logaritmo (elevando 10 a ambos os lados):

$\dfrac{A_2}{A_1} = 10^{2{,}0} = 100$

\textbf{A amplitude do terremoto de magnitude 6,0 foi 100 vezes maior que a do terremoto de magnitude 4,0.}

\textbf{b) Interpretação da variação}

O erro do morador é tratar a escala Richter como se fosse uma escala \textbf{linear}, na qual a diferença de magnitudes (2 pontos) corresponderia a um aumento percentual pequeno e proporcional a essa diferença (como 20\%).

No entanto, a função $M = \log_{10}(A/A_0)$ é \textbf{crescente}, mas de forma \textbf{logarítmica}: cada unidade a mais na magnitude corresponde a \textbf{multiplicar} a amplitude por 10, e não a somar uma parcela fixa. Por isso, uma diferença de 2 unidades na magnitude não corresponde a uma soma de 20\% na amplitude, mas sim a uma multiplicação por $10^2 = 100$, como mostrado no item (a).

Ou seja, o terremoto de magnitude 6,0 não foi "um pouco mais forte": suas ondas tiveram amplitude \textbf{cem vezes maior} que as do terremoto de magnitude 4,0. Isso ilustra por que pequenas variações no valor de $M$ (grandeza medida na escala logarítmica) representam variações enormes na grandeza física real (a amplitude $A$).

\bigskip\noindent\textbf{Q070 · múltipla escolha · média · criar · teórica · gabarito conferido}\par\nopagebreak
A magnitude de um abalo sísmico na escala Richter é definida por $M = \log_{10}\left(\dfrac{I}{I_0}\right)$, em que $I$ é a intensidade da onda sísmica registrada pelo sismógrafo e $I_0$ é uma intensidade de referência fixa, igual para todos os abalos.

Um sismólogo estuda uma situação hipotética em que dois abalos, de magnitudes $M_1$ e $M_2$, ocorrem na mesma falha geológica e suas ondas se sobrepõem exatamente no mesmo instante, de modo que a intensidade total registrada pelo sismógrafo passa a ser a soma das intensidades individuais dos dois abalos, isto é, $I = I_1 + I_2$.

Elabore a expressão geral que forneça a magnitude resultante $M$ em função de $M_1$ e $M_2$ nessa situação e utilize-a para calcular o valor de $M$ quando $M_1 = 5$ e $M_2 = 6$.

Assinale a alternativa que apresenta corretamente o valor de $M$ (aproximado a duas casas decimais).
\begin{enumerate}[label=(\alph*)]
\item $M = 11$
\item $M = 6$
\item $M = 5{,}5$
\item $M \approx 6{,}04$
\end{enumerate}

\textbf{Gabarito.} M $\approx$ 6,04

\textbf{Resolução.} \textbf{Passo 1 — Expressar as intensidades individuais.}

Da definição $M = \log_{10}(I/I_0)$, isolando $I$: $I = I_0 \cdot 10^{M}$.

Logo:
$$I_1 = I_0 \cdot 10^{M_1}, \qquad I_2 = I_0 \cdot 10^{M_2}$$

\textbf{Passo 2 — Somar as intensidades (não as magnitudes).}

Como a onda resultante tem intensidade $I = I_1 + I_2$:
$$I = I_0\left(10^{M_1} + 10^{M_2}\right)$$

\textbf{Passo 3 — Elaborar a expressão geral para a magnitude resultante.}

Substituindo em $M = \log_{10}(I/I_0)$:
$$M = \log_{10}\left(10^{M_1} + 10^{M_2}\right)$$

Essa é a expressão geral pedida: ela mostra que magnitudes \textbf{não se somam diretamente} quando as intensidades se somam — é necessário somar as intensidades (potências de 10) e só depois aplicar o logaritmo.

\textbf{Passo 4 — Aplicar aos valores dados.}

Com $M_1 = 5$ e $M_2 = 6$:
$$M = \log_{10}\left(10^{5} + 10^{6}\right) = \log_{10}(1{,}1 \times 10^{6}) = 6 + \log_{10}(1{,}1)$$

Como $\log_{10}(1{,}1) \approx 0{,}0414$:
$$M \approx 6{,}04$$

\textbf{Interpretação da variação:} embora a intensidade total ($1{,}1\times10^6\,I_0$) seja $10\%$ maior do que a intensidade do abalo mais forte sozinho ($10^6\,I_0$), a magnitude aumenta apenas de $6$ para $6{,}04$. Isso evidencia a natureza logarítmica da escala: grandes variações de intensidade correspondem a pequenas variações de magnitude, e vice-versa.

\bigskip\noindent\textbf{Q075 · discursiva · fácil · criar · teórica · gabarito conferido}\par\nopagebreak
O potencial hidrogeniônico (pH) de uma solução aquosa é definido por $pH = -\log_{10}[H^+]$, em que $[H^+]$ é a concentração de íons hidrogênio (em mol/L). Quando uma solução sofre diluição ou concentração, sua concentração de $H^+$ é multiplicada por um fator $k>0$ (por exemplo, $k = \tfrac{1}{10}$ significa que a concentração caiu para um décimo do valor original; $k=5$ significa que ela quintuplicou).

\textbf{(a)} Mostre, usando as propriedades dos logaritmos, que a variação do pH causada por esse fator, $\Delta pH = pH_{final} - pH_{inicial}$, pode ser escrita como uma fórmula que depende apenas de $k$ (isto é, não depende do valor inicial de $[H^+]$).

\textbf{(b)} Uma solução tem $pH_{inicial} = 4{,}0$. Ela é diluída de modo que a concentração de $H^+$ passa a ser $\tfrac{1}{10}$ da concentração original ($k = \tfrac{1}{10}$). Calcule o novo pH e a variação $\Delta pH$, e diga se a solução ficou mais ácida ou mais básica.

\textbf{(c)} Agora é sua vez de elaborar um problema: invente uma situação com um pH inicial diferente de $4{,}0$ e um fator $k$ diferente de $\tfrac{1}{10}$ (pode representar uma diluição, com $0<k<1$, ou uma concentração, com $k>1$). Escreva o enunciado completo do seu problema (informando claramente o pH inicial e o fator $k$ escolhidos) e, em seguida, resolva-o usando a fórmula geral obtida no item (a), calculando o novo pH, a variação $\Delta pH$ e interpretando se a solução se tornou mais ácida ou mais básica.

\textbf{Gabarito.} (a) $\Delta pH = -\log_{10}(k)$, independente do pH inicial. (b) $pH_{final}=5{,}0$; $\Delta pH = +1$; a solução ficou menos ácida. (c) Resposta aberta: será considerada correta qualquer situação-problema coerente (com pH inicial e fator $k$ escolhidos pelo estudante, diferentes dos do item b) cuja resolução aplique corretamente $\Delta pH=-\log_{10}(k)$ e $pH_{final}=pH_{inicial}-\log_{10}(k)$, com interpretação correta do sinal da variação (ácida/básica).

\textbf{Resolução.} \textbf{Item (a):}

Seja $[H^+]_i$ a concentração inicial e $[H^+]_f = k\cdot [H^+]_i$ a concentração final, com $k>0$.

$pH_{final} = -\log_{10}\big(k\cdot [H^+]_i\big)$

Usando a propriedade $\log(ab) = \log a + \log b$:

$pH_{final} = -\log_{10}(k) - \log_{10}([H^+]_i) = -\log_{10}(k) + pH_{inicial}$

Logo:

$$\Delta pH = pH_{final} - pH_{inicial} = -\log_{10}(k)$$

Essa expressão mostra que a \textbf{variação} do pH depende apenas do fator multiplicativo $k$ aplicado à concentração, e não do valor absoluto de $[H^+]_i$ — característica típica de grandezas relacionadas por escala logarítmica: variações \textbf{multiplicativas} na grandeza original (concentração) correspondem a variações \textbf{aditivas} na grandeza transformada (pH).

\textbf{Item (b):}

Com $pH_{inicial} = 4{,}0$ e $k = \tfrac{1}{10}$:

$$\Delta pH = -\log_{10}\left(\tfrac{1}{10}\right) = -(-1) = 1$$

$$pH_{final} = 4{,}0 + 1 = 5{,}0$$

Como o pH aumentou (de 4,0 para 5,0), a solução ficou \textbf{menos ácida} (mais próxima da neutralidade), confirmando que diluir a concentração de $H^+$ aumenta o pH.

\textbf{Item (c):}

Resposta pessoal (aberta). O estudante deve:

1. Escolher um $pH_{inicial}$ diferente de 4,0 e um fator $k$ diferente de $\tfrac{1}{10}$, escrevendo um enunciado coerente (por exemplo: "Uma solução tem pH inicial 6,0 e sua concentração de $H^+$ é multiplicada por 100 (k=100). Qual o novo pH?").
2. Aplicar corretamente a fórmula $\Delta pH = -\log_{10}(k)$.
3. Calcular $pH_{final} = pH_{inicial} - \log_{10}(k)$.
4. Interpretar corretamente o sinal da variação: se $k>1$ (concentração aumenta), $\Delta pH<0$ e a solução fica mais ácida; se $0<k<1$ (diluição), $\Delta pH>0$ e a solução fica menos ácida (mais básica).

Exemplo de resolução válida (com os dados sugeridos acima): $\Delta pH = -\log_{10}(100) = -2$, então $pH_{final} = 6{,}0 - 2 = 4{,}0$; a solução ficou mais ácida.
\subsection*{EM13MAT306}
{\small\itshape Resolver e elaborar problemas em contextos que envolvem fenômenos periódicos reais, como ondas sonoras, ciclos menstruais, movimentos cíclicos, entre outros, e comparar suas representações com as funções seno e cosseno, no plano cartesiano, com ou sem apoio de aplicativos de álgebra e geometria.}

\bigskip\noindent\textbf{Q044 · discursiva · fácil · aplicar · aplicada · gabarito conferido}\par\nopagebreak
As marés de um pequeno porto pesqueiro variam de forma periódica ao longo do dia, devido à atração gravitacional da Lua e do Sol. Um estudo local registrou que, tomando $t = 0$ como meia-noite (0h), a maré atinge seu nível máximo de $2{,}0$ m exatamente à meia-noite, cai até um nível mínimo de $1{,}0$ m às 6h da manhã, retorna ao máximo às 12h (meio-dia) e repete esse padrão a cada 12 horas ao longo de todo o dia.

a) Identifique a amplitude, o período e o deslocamento vertical (translação do eixo médio) desse fenômeno periódico.

b) Escreva uma função $h(t)$, em metros, que descreva a altura da maré em função do tempo $t$ (em horas), $t \in [0,24)$, usando a função cosseno. Compare essa função com o gráfico da função cosseno "pura" $y=\cos(t)$, indicando quais transformações (translação vertical e compressão/dilatação horizontal) foram necessárias para obter o modelo da maré.

c) Determine a altura da maré às 3h da manhã e explique, com base no gráfico da função obtida em (b), por que esse valor corresponde exatamente à média entre o nível máximo e o mínimo.

\textbf{Gabarito.} a) Amplitude = 0,5 m; período = 12 h; deslocamento vertical = 1,5 m. b) $h(t) = 1{,}5 + 0{,}5\cos\left(\frac{\pi}{6}t\right)$, obtida comprimindo horizontalmente (fator $\pi/6$), reduzindo a amplitude (fator 0,5) e deslocando verticalmente (+1,5) o gráfico de $y=\cos(t)$. c) $h(3) = 1{,}5$ m, que é a média entre o máximo e o mínimo.

\textbf{Resolução.} \textbf{a) Amplitude, período e deslocamento vertical}

O nível máximo é $2{,}0$ m e o mínimo é $1{,}0$ m. A amplitude é metade da diferença entre esses valores:
$$A = \frac{2{,}0 - 1{,}0}{2} = 0{,}5 \text{ m}$$

O deslocamento vertical (altura média, em torno da qual a maré oscila) é a média entre máximo e mínimo:
$$k = \frac{2{,}0 + 1{,}0}{2} = 1{,}5 \text{ m}$$

Como o padrão se repete a cada 12 horas (máximo em 0h e novamente em 12h), o período é:
$$T = 12 \text{ horas}$$

\textbf{b) Construção da função e comparação com $y=\cos(t)$}

A função cosseno padrão $y=\cos(t)$ tem amplitude 1, período $2\pi$, oscila entre $-1$ e $1$, e atinge seu máximo em $t=0$. Como aqui a maré também está no máximo em $t=0$ (meia-noite), o cosseno é a escolha natural (sem precisar de deslocamento horizontal de fase).

Para transformar $y=\cos(t)$ no modelo da maré:

\begin{itemize}[nosep,leftmargin=*]
\item \textbf{Compressão horizontal}: o período deve passar de $2\pi$ para 12 horas. Como o período de $\cos(bt)$ é $\dfrac{2\pi}{b}$, temos:
\end{itemize}
$$\frac{2\pi}{b} = 12 \;\Rightarrow\; b = \frac{2\pi}{12} = \frac{\pi}{6}$$

\begin{itemize}[nosep,leftmargin=*]
\item \textbf{Dilatação vertical (amplitude)}: multiplica-se por $A = 0{,}5$, reduzindo a oscilação de $[-1,1]$ para $[-0{,}5;\,0{,}5]$.
\end{itemize}

\begin{itemize}[nosep,leftmargin=*]
\item \textbf{Translação vertical}: soma-se $k=1{,}5$, deslocando o eixo de oscilação de $y=0$ para $y=1{,}5$.
\end{itemize}

Assim, o modelo é:
$$h(t) = 1{,}5 + 0{,}5\cos\!\left(\frac{\pi}{6}t\right), \quad t \in [0,24)$$

O gráfico de $h(t)$ é o gráfico de $y=\cos(t)$ comprimido horizontalmente (período 12 em vez de $2\pi$), "achatado" verticalmente (amplitude 0,5 em vez de 1) e deslocado para cima em 1,5 unidades.

\textbf{c) Altura às 3h}

Substituindo $t=3$:
$$h(3) = 1{,}5 + 0{,}5\cos\!\left(\frac{\pi}{6}\cdot 3\right) = 1{,}5 + 0{,}5\cos\!\left(\frac{\pi}{2}\right) = 1{,}5 + 0{,}5\cdot 0 = 1{,}5 \text{ m}$$

No gráfico da função cosseno, $t=\dfrac{\pi}{2}$ (aqui, $t=3$ h) é exatamente o ponto em que a curva cruza o eixo médio, entre o pico máximo (em $t=0$) e o pico mínimo (em $t=6$). Por isso o valor obtido coincide com a média entre o nível máximo (2,0 m) e o mínimo (1,0 m), confirmando que $t=3$h corresponde à passagem pelo nível médio da maré.

\bigskip\noindent\textbf{Q052 · múltipla escolha · fácil · analisar · teórica · gabarito conferido}\par\nopagebreak
Uma boia de monitoramento registrou a altura da água (em metros) em um cais, em intervalos regulares de 3 horas, ao longo de um dia. Observou-se que o padrão de subida e descida se repete a cada 12 horas (maré semidiurna). Os dados coletados foram:

\begin{center}
\begin{tabular}{|p{0.143\linewidth}|p{0.143\linewidth}|p{0.143\linewidth}|p{0.143\linewidth}|p{0.143\linewidth}|p{0.143\linewidth}|}
\hline
\textbf{t (horas)} & \textbf{0} & \textbf{3} & \textbf{6} & \textbf{9} & \textbf{12} \\
\hline
h (metros) & 8 & 5 & 2 & 5 & 8 \\
\hline
\end{tabular}
\end{center}

Esse tipo de fenômeno periódico pode ser representado no plano cartesiano por uma função seno ou por uma função cosseno. Observando em que instante ocorrem os valores máximo e mínimo da maré, identifique qual das duas representações (seno ou cosseno) é a mais adequada para modelar h(t) e, em seguida, utilize o modelo escolhido para determinar a altura da água no instante t = 2 horas.
\begin{enumerate}[label=(\alph*)]
\item 6,5 m
\item 9,5 m
\item 3,5 m
\item 7,6 m (aproximadamente)
\end{enumerate}

\textbf{Gabarito.} 6,5 m

\textbf{Resolução.} \textbf{Passo 1 — Amplitude e valor médio (deslocamento vertical).}
O valor máximo é $8$ m e o mínimo é $2$ m. Logo:
$$k=\frac{8+2}{2}=5,\qquad A=\frac{8-2}{2}=3$$

\textbf{Passo 2 — Período e frequência angular.}
O padrão se repete a cada $12$ h, então $T=12$ e
$$\omega=\frac{2\pi}{T}=\frac{2\pi}{12}=\frac{\pi}{6}$$

\textbf{Passo 3 — Escolha entre seno e cosseno.}
A função seno, sem deslocamento de fase, vale $0$ (o valor médio $k$) em $t=0$ e cresce a partir daí. Já a função cosseno atinge seu \textbf{valor máximo} exatamente em $t=0$. Como a tabela mostra $h(0)=8$ (o máximo), a representação adequada é a \textbf{cosseno}, sem necessidade de deslocamento de fase:
$$h(t)=5+3\cos\left(\frac{\pi}{6}t\right)$$

Conferindo com a tabela: $h(3)=5+3\cos(\pi/2)=5$ $\checkmark$; $h(6)=5+3\cos(\pi)=2$ $\checkmark$; $h(9)=5+3\cos(3\pi/2)=5$ $\checkmark$; $h(12)=5+3\cos(2\pi)=8$ $\checkmark$.

\textbf{Passo 4 — Cálculo em $t=2$.}
$$h(2)=5+3\cos\left(\frac{\pi}{6}\cdot 2\right)=5+3\cos\left(\frac{\pi}{3}\right)=5+3\cdot\frac{1}{2}=6{,}5\text{ m}$$

Portanto, a altura da água em $t=2$ h é \textbf{6,5 m}.

\bigskip\noindent\textbf{Q065 · discursiva · difícil · aplicar · aplicada · gabarito conferido}\par\nopagebreak
Um oceanógrafo registra a altura da maré (em metros) em um pequeno porto ao longo de um dia. Ele observa que a maré varia de forma periódica e simétrica em torno de um nível médio, atingindo seu valor máximo de 1,8 m exatamente à meia-noite ($t=0$, com $t$ em horas) e seu valor mínimo de 0,2 m exatamente às 6h da manhã. Esse padrão se repete a cada 12 horas.

a) Modele a altura da maré, em metros, por uma função do tipo $h(t) = A\cos(Bt) + D$, com $t$ em horas contado a partir da meia-noite. Determine explicitamente os valores de $A$, $B$ e $D$, justificando cada um a partir dos dados do fenômeno (amplitude, período e deslocamento vertical do gráfico em relação ao eixo $t$).

b) Determine todos os horários dentro das primeiras 12 horas do dia (isto é, para $0 \le t < 12$) em que a maré atinge exatamente 1,4 m de altura. Para cada horário encontrado, indique se a maré está subindo ou descendo naquele instante, justificando sua resposta apenas a partir do comportamento gráfico da função cosseno (sem usar derivadas).

\textbf{Gabarito.} h(t) = 0,8·cos($\pi$t/6) + 1,0 (A=0,8 m; B=$\pi$/6 rad/h; D=1,0 m). A maré atinge 1,4 m às 2h (descendo) e às 10h (subindo).

\textbf{Resolução.} \textbf{a) Construção do modelo}

O valor máximo da maré é $1{,}8$ m e o mínimo é $0{,}2$ m. Como $h(t)=A\cos(Bt)+D$ oscila entre $D-A$ e $D+A$:

$$D = \frac{1{,}8+0{,}2}{2} = 1{,}0 \qquad A = \frac{1{,}8-0{,}2}{2} = 0{,}8$$

O fenômeno se repete a cada 12 horas, logo o período é $T=12$. Como $T=\dfrac{2\pi}{B}$:

$$B = \frac{2\pi}{12} = \frac{\pi}{6}$$

Como a maré atinge seu \textbf{máximo} exatamente em $t=0$, e $\cos(0)=1$ é o valor máximo do cosseno, a escolha de cosseno (sem deslocamento horizontal) é coerente com o fenômeno. Assim:

$$h(t) = 0{,}8\cos\!\left(\frac{\pi t}{6}\right) + 1{,}0$$

Verificação: $h(0)=0{,}8(1)+1=1{,}8$ (máximo, ok) e $h(6)=0{,}8\cos(\pi)+1=0{,}8(-1)+1=0{,}2$ (mínimo, ok).

\textbf{b) Horários em que $h(t)=1{,}4$}

$$0{,}8\cos\!\left(\frac{\pi t}{6}\right)+1 = 1{,}4 \;\Rightarrow\; \cos\!\left(\frac{\pi t}{6}\right) = 0{,}5$$

As soluções gerais de $\cos\theta = 0{,}5$ são $\theta = \dfrac{\pi}{3}+2k\pi$ ou $\theta = -\dfrac{\pi}{3}+2k\pi$ (equivalente a $\dfrac{5\pi}{3}+2k\pi$).

Fazendo $\theta = \dfrac{\pi t}{6}$, temos $t = \dfrac{6\theta}{\pi}$:

\begin{itemize}[nosep,leftmargin=*]
\item Para $\theta=\dfrac{\pi}{3}$: $t = 6\cdot\dfrac{1}{3} = 2$
\item Para $\theta=\dfrac{5\pi}{3}$: $t = 6\cdot\dfrac{5}{3} = 10$
\end{itemize}

Dentro de $0\le t<12$, as soluções são $t=2$ e $t=10$.

\textbf{Análise de crescimento/decrescimento}

Como $h(0)=1{,}8$ é máximo e $h(6)=0{,}2$ é mínimo, no intervalo $(0,6)$ a função cosseno está descendo de seu valor máximo ao mínimo — ou seja, $h(t)$ é \textbf{decrescente} nesse trecho. Logo, em $t=2$ (que está entre 0 e 6), a maré está \textbf{descendo}.

No intervalo $(6,12)$, a função retorna do mínimo ($h(6)=0{,}2$) ao próximo máximo ($h(12)=1{,}8$), portanto $h(t)$ é \textbf{crescente} nesse trecho. Logo, em $t=10$ (entre 6 e 12), a maré está \textbf{subindo}.

\textbf{Conclusão:} às 02h00 a maré está a 1,4 m e descendo; às 10h00 a maré está a 1,4 m e subindo.

\bigskip\noindent\textbf{Q073 · múltipla escolha · difícil · aplicar · teórica · gabarito conferido}\par\nopagebreak
O nível da maré em certo ponto da costa varia de forma cíclica ao longo do dia e pode ser modelado pela função $h(t) = 2\cos\left(\dfrac{\pi}{6}(t-3)\right) + 3$, em que $t$ é o tempo em horas contado a partir da meia-noite ($t \ge 0$) e $h(t)$ é a altura da maré, em metros. Interpretando o gráfico dessa função como uma curva cossenoidal no plano cartesiano, determine: a altura máxima e a altura mínima da maré, o período do fenômeno e todos os instantes $t$, com $0 \le t < 12$, em que a maré atinge exatamente 4 metros de altura.
\begin{enumerate}[label=(\alph*)]
\item Altura máxima 2 m; altura mínima -2 m; período 12 h; maré de 4 m em $t=1$ h e $t=5$ h.
\item Altura máxima 5 m; altura mínima 1 m; período 12 h; maré de 4 m em $t=1$ h e $t=5$ h.
\item Altura máxima 5 m; altura mínima 1 m; período 6 h; maré de 4 m em $t=1$ h e $t=5$ h.
\item Altura máxima 5 m; altura mínima 1 m; período 12 h; maré de 4 m apenas em $t=1$ h.
\end{enumerate}

\textbf{Gabarito.} Altura máxima 5 m; altura mínima 1 m; período 12 h; maré de 4 m em $t=1$ h e $t=5$ h.

\textbf{Resolução.} \textbf{1. Identificando os parâmetros da função cossenoidal}

A função tem a forma $h(t) = A\cos(B(t-C)) + D$, com $A=2$, $B=\dfrac{\pi}{6}$, $C=3$ e $D=3$.

\textbf{2. Máximo e mínimo}

Como $-1 \le \cos(\cdot) \le 1$, o valor máximo ocorre quando o cosseno vale $1$:
$$h_{max} = 2(1) + 3 = 5 \text{ m}$$
e o valor mínimo ocorre quando o cosseno vale $-1$:
$$h_{min} = 2(-1) + 3 = 1 \text{ m}$$
Note que a amplitude ($A=2$) mede a distância entre a curva e a reta média $D=3$, e não o valor máximo em si — esse é um erro comum.

\textbf{3. Período}

O período de $\cos(Bt)$ é $T = \dfrac{2\pi}{B}$. Aqui $B = \dfrac{\pi}{6}$, logo:
$$T = \frac{2\pi}{\pi/6} = 12 \text{ horas}$$

\textbf{4. Instantes em que $h(t) = 4$}

$$2\cos\left(\frac{\pi}{6}(t-3)\right) + 3 = 4 \;\Rightarrow\; \cos\left(\frac{\pi}{6}(t-3)\right) = \frac{1}{2}$$

O cosseno vale $\tfrac{1}{2}$ quando o arco é $\pm\dfrac{\pi}{3}$ (mais múltiplos de $2\pi$, que aqui saem do intervalo pedido). Assim:
$$\frac{\pi}{6}(t-3) = \frac{\pi}{3} \;\Rightarrow\; t-3 = 2 \;\Rightarrow\; t = 5$$
$$\frac{\pi}{6}(t-3) = -\frac{\pi}{3} \;\Rightarrow\; t-3 = -2 \;\Rightarrow\; t = 1$$

Ambos os valores, $t=1$ e $t=5$, pertencem ao intervalo $0 \le t < 12$ (correspondendo à simetria da curva cossenoidal em torno do seu pico em $t=3$, onde $h=5$).

\textbf{Conclusão:} altura máxima $=5$ m; altura mínima $=1$ m; período $=12$ h; a maré atinge 4 m em $t=1$ h e $t=5$ h.

\bigskip\noindent\textbf{Q078 · discursiva · média · analisar · teórica · gabarito conferido}\par\nopagebreak
Em um porto, a altura da maré (em metros) varia de forma periódica ao longo do dia. Registros mostram que a maré atinge seu nível máximo de 5 m às 3h da manhã e seu nível mínimo de 1 m às 9h da manhã, repetindo esse padrão a cada 12 horas — sendo $t$ o tempo em horas contado a partir da meia-noite ($t=0$).

Um estagiário do porto, ao tentar modelar esse fenômeno, propôs a função
$$h(t) = 2\cos\left(\frac{\pi}{6}t\right) + 3,$$
com $t$ em horas, afirmando que ela representa corretamente a maré descrita.

a) Determine a amplitude, o período e o deslocamento vertical do fenômeno periódico descrito pelos registros.

b) Escreva a função $h(t)$, na forma $h(t) = A\cos(B(t-C)) + D$, que representa corretamente a maré observada.

c) Avalie se a proposta do estagiário está correta, comparando os valores previstos por sua função com os dados observados nos instantes $t=3$ h e $t=9$ h. Caso ela esteja incorreta, explique qual característica do gráfico (amplitude, período ou deslocamento de fase) foi tratada de forma equivocada.

\textbf{Gabarito.} a) Amplitude $A=2$ m, período $T=12$ h, deslocamento vertical $D=3$ m. b) $h(t) = 2\cos\left(\frac{\pi}{6}(t-3)\right)+3$. c) A proposta do estagiário está incorreta: amplitude e período estão certos, mas falta o deslocamento de fase ($C=3$); sua função dá $h(3)=3$ e $h(9)=3$, valores que não correspondem aos dados reais (5 m e 1 m).

\textbf{Resolução.} \textbf{a) Amplitude, período e deslocamento vertical}

A maré varia entre um máximo de $5$ m e um mínimo de $1$ m. Assim:

\begin{itemize}[nosep,leftmargin=*]
\item Amplitude: $A = \dfrac{5-1}{2} = 2$
\item Deslocamento vertical: $D = \dfrac{5+1}{2} = 3$
\end{itemize}

O intervalo entre o instante de máximo ($t=3$) e o instante de mínimo ($t=9$) corresponde a \textbf{meia} volta do ciclo (de crista a cale), ou seja, $T/2 = 9-3 = 6$ h. Logo:

$$T = 12 \text{ h}$$

\textbf{b) Construção da função correta}

O coeficiente $B$ (frequência angular) é dado por:
$$B = \frac{2\pi}{T} = \frac{2\pi}{12} = \frac{\pi}{6}$$

Como o cosseno atinge seu valor máximo quando o argumento é zero, e sabemos que o máximo ocorre em $t=3$, o deslocamento de fase deve ser $C = 3$. Assim, a função que representa corretamente o fenômeno é:

$$h(t) = 2\cos\left(\frac{\pi}{6}(t-3)\right) + 3$$

\textbf{Verificação:}
\begin{itemize}[nosep,leftmargin=*]
\item $h(3) = 2\cos(0) + 3 = 2 + 3 = 5$ $\checkmark$ (máximo)
\item $h(9) = 2\cos\left(\frac{\pi}{6}\cdot 6\right) + 3 = 2\cos(\pi) + 3 = -2+3 = 1$ $\checkmark$ (mínimo)
\end{itemize}

\textbf{c) Análise da proposta do estagiário}

A função proposta é $h(t) = 2\cos\left(\frac{\pi}{6}t\right)+3$, que tem a \textbf{mesma amplitude} ($A=2$) e o \textbf{mesmo período} ($T=12$, pois $B=\pi/6$ é igual) da função correta — porém \textbf{sem deslocamento de fase} (equivale a assumir $C=0$, ou seja, que o máximo ocorre à meia-noite).

Calculando os valores previstos pela proposta nos instantes observados:

\begin{itemize}[nosep,leftmargin=*]
\item $h(3) = 2\cos\left(\frac{\pi}{6}\cdot 3\right)+3 = 2\cos\left(\frac{\pi}{2}\right)+3 = 2(0)+3 = 3$
\item $h(9) = 2\cos\left(\frac{\pi}{6}\cdot 9\right)+3 = 2\cos\left(\frac{3\pi}{2}\right)+3 = 2(0)+3 = 3$
\end{itemize}

Os valores previstos ($3$ m em ambos os instantes) \textbf{não coincidem} com os valores observados ($5$ m em $t=3$ e $1$ m em $t=9$). Portanto, \textbf{a proposta do estagiário está incorreta}: embora amplitude e período estejam certos, ele \textbf{não aplicou o deslocamento horizontal (de fase)} necessário para que o pico do gráfico do cosseno coincida com o instante real do máximo da maré ($t=3$, e não $t=0$).

\bigskip\noindent\textbf{Q086 · múltipla escolha · média · aplicar · aplicada · gabarito conferido}\par\nopagebreak
Uma roda-gigante tem 40 m de diâmetro e seu centro está a 22 m de altura em relação ao solo. Ela gira com velocidade constante, completando uma volta a cada 12 minutos. Um passageiro embarca no ponto mais baixo da roda (rente ao chão de segurança, a 2 m de altura) no instante $t = 0$ minuto. A altura $h(t)$, em metros, desse passageiro em relação ao solo, $t$ minutos após o embarque, pode ser descrita por uma função do tipo cosseno. Considerando o gráfico de $h(t)$ ao longo de todo o primeiro giro completo ($0 \le t \le 12$ min), assinale a alternativa que apresenta corretamente a lei de $h(t)$ e o intervalo de tempo, nesse primeiro giro, em que a altura do passageiro é maior que 32 m.
\begin{enumerate}[label=(\alph*)]
\item $h(t) = 22 - 20\cos\left(\dfrac{\pi t}{6}\right)$, e a altura é maior que 32 m para $t \in (4, 8)$ minutos.
\item $h(t) = 22 - 20\cos\left(\dfrac{\pi t}{3}\right)$, e a altura é maior que 32 m para $t \in (2, 4)$ minutos.
\item $h(t) = 22 - 40\cos\left(\dfrac{\pi t}{6}\right)$, e a altura é maior que 32 m para $t \in (4, 8)$ minutos.
\item $h(t) = 22 + 20\cos\left(\dfrac{\pi t}{6}\right)$, e a altura é maior que 32 m para $t \in (0, 2) \cup (10, 12)$ minutos.
\end{enumerate}

\textbf{Gabarito.} h(t) = 22 - 20cos($\pi$t/6); altura > 32 m para t entre 4 e 8 minutos (alternativa A)

\textbf{Resolução.} \textbf{1. Amplitude:} o raio da roda é metade do diâmetro: $R = 40/2 = 20$ m. Essa é a amplitude do movimento (a amplitude é o raio, não o diâmetro).

\textbf{2. Deslocamento vertical:} o centro da roda está a 22 m do solo, então a função oscila em torno de $D = 22$.

\textbf{3. Período:} a roda completa uma volta a cada 12 minutos, logo $T = 12$ e o coeficiente angular do argumento é $b = \dfrac{2\pi}{T} = \dfrac{2\pi}{12} = \dfrac{\pi}{6}$.

\textbf{4. Fase inicial:} em $t=0$ o passageiro está no ponto \textbf{mais baixo}, ou seja, $h(0) = D - R = 22 - 20 = 2$. Como a função cosseno começa em seu valor máximo, usamos o \textbf{cosseno negativo} para que o ponto de partida seja o mínimo:
$$h(t) = 22 - 20\cos\left(\frac{\pi t}{6}\right)$$
Confira: $h(0) = 22 - 20(1) = 2$ $\checkmark$; $h(6) = 22 - 20(-1) = 42$ (ponto mais alto, na metade da volta) $\checkmark$.

\textbf{5. Leitura gráfica — altura maior que 32 m:} no gráfico de $h(t)$, essa é uma curva cosseno invertida oscilando entre 2 m (mínimo) e 42 m (máximo), com centro em 22 m. Queremos o intervalo em que a curva fica \textbf{acima da reta horizontal $y=32$}.

Resolvendo algebricamente para confirmar a leitura do gráfico:
$$22 - 20\cos\left(\frac{\pi t}{6}\right) > 32 \;\Rightarrow\; \cos\left(\frac{\pi t}{6}\right) < -\frac{1}{2}$$
No intervalo $0 \le t \le 12$ (ou seja, $0 \le \frac{\pi t}{6} \le 2\pi$), o cosseno é menor que $-\tfrac12$ quando o ângulo está entre $120^\circ $ e $240^\circ $:
$$\frac{\pi t}{6} = \frac{2\pi}{3} \Rightarrow t = 4 \qquad \frac{\pi t}{6} = \frac{4\pi}{3} \Rightarrow t = 8$$
Isso corresponde exatamente ao trecho do gráfico em que a curva está no seu 'pico' central, entre os instantes $t=4$ min e $t=8$ min — coerente com a simetria em torno de $t=6$ (o instante do máximo).

\textbf{Conclusão:} $h(t) = 22 - 20\cos\left(\dfrac{\pi t}{6}\right)$ e a altura é maior que 32 m para $t \in (4, 8)$ minutos.
\subsection*{EM13MAT401}
{\small\itshape Converter representações algébricas de funções polinomiais de 1º grau para representações geométricas no plano cartesiano, distinguindo os casos nos quais o comportamento é proporcional, recorrendo ou não a softwares ou aplicativos de álgebra e geometria dinâmica.}

\bigskip\noindent\textbf{Q025 · discursiva · fácil · entender · teórica · gabarito conferido}\par\nopagebreak
Considere as funções reais f(x) = 2x, g(x) = 2x + 3 e h(x) = -2x, definidas para todo x real.

a) Determine as coordenadas de dois pontos do gráfico de cada função e utilize-os para esboçar, num mesmo plano cartesiano, as retas correspondentes a f, g e h.

b) Observando os três gráficos esboçados, indique quais retas passam pela origem do plano cartesiano e quais não passam.

c) A partir do que você observou, enuncie uma regra geral: dada uma lei y = ax + b (com a $\neq$ 0), como decidir, apenas olhando os coeficientes a e b (sem construir o gráfico), se a reta correspondente passará ou não pela origem? Qual é o nome usual dado ao tipo de função cujo gráfico passa pela origem, e qual o nome dado ao caso em que isso não ocorre?

d) Compare os coeficientes angulares de f, g e h. Existe alguma relação geométrica de paralelismo entre duas dessas retas? Justifique sua resposta e descreva, em palavras, como o gráfico de g pode ser obtido a partir do gráfico de f.

\textbf{Gabarito.} b) f e h passam pela origem; g não passa. c) A reta y=ax+b passa pela origem se, e somente se, b=0; nesse caso a função é linear/proporcional, e quando b$\neq$0 é afim não proporcional. d) f e g são paralelas (mesmo a=2); h não é paralela a nenhuma delas; o gráfico de g é o gráfico de f transladado 3 unidades para cima.

\textbf{Resolução.} \textbf{a) Pontos e esboço}

Para $f(x)=2x$: $f(0)=0$ e $f(1)=2$, logo a reta passa pelos pontos $(0,0)$ e $(1,2)$.

Para $g(x)=2x+3$: $g(0)=3$ e $g(1)=5$, logo a reta passa pelos pontos $(0,3)$ e $(1,5)$.

Para $h(x)=-2x$: $h(0)=0$ e $h(1)=-2$, logo a reta passa pelos pontos $(0,0)$ e $(1,-2)$.

Ao marcar esses pontos no plano cartesiano e traçar as retas, obtém-se: $f$ crescente passando pela origem, $g$ crescente deslocada para cima, e $h$ decrescente passando pela origem.

\textbf{b) Quais passam pela origem}

$f(0)=0$ e $h(0)=0$, então as retas de $f$ e de $h$ passam pela origem $(0,0)$.

$g(0)=3\neq 0$, então a reta de $g$ \textbf{não} passa pela origem.

\textbf{c) Regra geral}

Uma reta $y=ax+b$ passa pela origem se, e somente se, o ponto $(0,0)$ satisfaz a equação, isto é, $a\cdot 0+b=0$, ou seja, $b=0$.

\begin{itemize}[nosep,leftmargin=*]
\item Se $b=0$, a lei é $y=ax$: a função é chamada de \textbf{função linear} (caso particular, também dito \textbf{proporcional}), e seu gráfico sempre passa pela origem.
\item Se $b\neq 0$, a função é chamada de \textbf{função afim} (não proporcional), e seu gráfico é uma reta paralela à de $y=ax$, mas deslocada verticalmente, não passando pela origem.
\end{itemize}

No exemplo: $f$ e $h$ são funções lineares (proporcionais), enquanto $g$ é afim não proporcional.

\textbf{d) Paralelismo e translação}

O coeficiente angular de $f$ é $a=2$, o de $g$ também é $a=2$, e o de $h$ é $a=-2$.

Como $f$ e $g$ têm o \textbf{mesmo coeficiente angular} ($a=2$), suas retas são \textbf{paralelas}. Já $h$ tem coeficiente angular diferente ($-2$), logo não é paralela nem a $f$ nem a $g$.

Geometricamente, como $g(x)=f(x)+3$, o gráfico de $g$ é obtido \textbf{transladando verticalmente} o gráfico de $f$ em $3$ unidades para cima — por isso as retas permanecem paralelas, mas apenas a de $f$ passa pela origem.

\bigskip\noindent\textbf{Q028 · discursiva · média · aplicar · aplicada · gabarito conferido}\par\nopagebreak
Duas empresas de transporte por aplicativo, TáxiJá e MoveCar, calculam o valor da corrida em função da distância percorrida, em quilômetros.

\begin{itemize}[nosep,leftmargin=*]
\item A TáxiJá cobra R\$\,2,50 por quilômetro rodado, sem qualquer taxa adicional.
\item A MoveCar cobra uma taxa fixa de R\$\,4,00 pela chamada do carro, mais R\$\,1,50 por quilômetro rodado.
\end{itemize}

Considere $x$ a distância percorrida, em quilômetros ($x \geq 0$), e $y$ o valor pago, em reais.

a) Escreva a lei de formação que representa o valor cobrado por cada empresa, em função de $x$.

b) Represente, num mesmo plano cartesiano, os gráficos das duas funções, indicando claramente o ponto em que cada reta corta o eixo $y$.

c) Uma dessas funções representa uma grandeza diretamente proporcional à distância percorrida, e a outra não. Identifique qual delas é a proporcional e explique, a partir do gráfico (ou da lei algébrica), como reconhecer esse comportamento.

d) Determine algebricamente a distância a partir da qual a corrida pela MoveCar passa a ser mais barata que pela TáxiJá, e indique as coordenadas do ponto em que os dois gráficos se cruzam.

\textbf{Gabarito.} a) $f_T(x)=2{,}5x$ e $f_M(x)=1{,}5x+4$. b) Retas crescentes: $f_T$ passa pela origem $(0,0)$; $f_M$ corta o eixo $y$ em $(0,4)$. c) $f_T$ é diretamente proporcional (passa pela origem, $b=0$); $f_M$ não é proporcional ($b=4\neq 0$). d) A partir de $x=4$ km a MoveCar fica mais barata; o ponto de intersecção das retas é $(4,10)$.

\textbf{Resolução.} \textbf{a) Leis de formação}

TáxiJá: valor proporcional à distância, sem taxa fixa:
$$f_T(x) = 2{,}5x$$

MoveCar: taxa fixa de R\$\,4,00 mais R\$\,1,50 por km:
$$f_M(x) = 1{,}5x + 4$$

\textbf{b) Representação gráfica}

Ambas são retas crescentes (coeficientes angulares positivos: $2{,}5$ e $1{,}5$).

\begin{itemize}[nosep,leftmargin=*]
\item $f_T$ passa pela origem $(0,0)$, pois $f_T(0)=0$.
\item $f_M$ corta o eixo $y$ no ponto $(0,4)$, pois $f_M(0)=4$.
\end{itemize}

Como $f_T$ tem inclinação maior, sua reta "sobe" mais rápido do que a de $f_M$, embora comece mais abaixo.

\textbf{c) Identificando o caso proporcional}

Uma função é diretamente proporcional quando tem a forma $y=ax$, ou seja, quando seu gráfico é uma reta que passa pela origem $(0,0)$.

\begin{itemize}[nosep,leftmargin=*]
\item $f_T(x)=2{,}5x$ tem $b=0$: é diretamente proporcional à distância.
\item $f_M(x)=1{,}5x+4$ tem $b=4 \neq 0$: sua reta não passa pela origem (corta o eixo $y$ em $4$), logo é apenas afim, não proporcional — mesmo que $x=0$, ainda se paga a taxa fixa de R\$\,4,00.
\end{itemize}

\textbf{d) Comparando os valores cobrados}

Igualando as duas expressões para achar o ponto de intersecção das retas:
$$2{,}5x = 1{,}5x + 4$$
$$2{,}5x - 1{,}5x = 4$$
$$x = 4$$

Substituindo em qualquer uma das leis:
$$f_T(4) = 2{,}5 \cdot 4 = 10$$

Logo, as retas se cruzam no ponto $(4, 10)$: para 4 km, ambas cobram R\$\,10,00.

Para $x>4$ (por exemplo, $x=5$): $f_T(5)=12{,}5$ e $f_M(5)=11{,}5$, ou seja, $f_M(x) < f_T(x)$.

Como o coeficiente angular de $f_T$ é maior, a partir de $x=4$ km a MoveCar torna-se mais barata que a TáxiJá.

\bigskip\noindent\textbf{Q031 · discursiva · difícil · aplicar · teórica · gabarito conferido}\par\nopagebreak
No plano cartesiano, as retas $r$ e $s$ representam, respectivamente, as funções polinomiais do 1º grau $f$ e $g$. A reta $r$ passa pela origem do sistema e pelo ponto $(3,6)$. A reta $s$ é paralela a $r$ e intercepta os eixos coordenados formando, com eles, um triângulo de área igual a 4 unidades de área, situado inteiramente no quarto quadrante (ou seja, com um vértice sobre o semieixo positivo das abscissas e outro sobre o semieixo negativo das ordenadas, além da origem).

a) Determine a lei de $f$ e classifique-a quanto à proporcionalidade, justificando pela posição geométrica da reta $r$.

b) Determine a lei de $g$. Nos seus cálculos aparecerão duas possibilidades para o coeficiente que desloca a reta; explique, usando a condição sobre o quadrante em que o triângulo deve estar, por que apenas uma delas é válida.

c) Classifique $g$ quanto à proporcionalidade e explique, em termos geométricos, por que a reta $s$ não passa pela origem.

\textbf{Gabarito.} f(x) = 2x (proporcional); g(x) = 2x - 4 (afim não proporcional, pois b = -4 $\neq$ 0)

\textbf{Resolução.} \textbf{a) Determinação de $f$}

Como $r$ passa pela origem $(0,0)$ e por $(3,6)$, o coeficiente angular é
$$a = \frac{6-0}{3-0} = 2.$$
Como a reta passa pela origem, o coeficiente linear é $b=0$, logo
$$f(x) = 2x.$$
Geometricamente, toda reta que passa pela origem representa uma função \textbf{proporcional} (caso particular da função afim com $b=0$).

\textbf{b) Determinação de $g$}

Como $s$ é paralela a $r$, tem o mesmo coeficiente angular $a=2$:
$$g(x) = 2x + b.$$

Os pontos onde $s$ corta os eixos são:
\begin{itemize}[nosep,leftmargin=*]
\item eixo $y$: $x=0 \Rightarrow y=b$, ponto $(0,b)$;
\item eixo $x$: $y=0 \Rightarrow x=-\dfrac{b}{2}$, ponto $\left(-\dfrac{b}{2},0\right)$.
\end{itemize}

Esses dois pontos, junto com a origem, formam um triângulo retângulo de área
$$A = \frac{1}{2}\cdot|b|\cdot\left|\frac{b}{2}\right| = \frac{b^2}{4}.$$

Impondo $A=4$:
$$\frac{b^2}{4}=4 \Rightarrow b^2=16 \Rightarrow b=4 \text{ ou } b=-4.$$

\textbf{Análise geométrica dos dois casos:}

\begin{itemize}[nosep,leftmargin=*]
\item Se $b=4$: os interceptos são $(0,4)$ e $(-2,0)$ — o triângulo fica no \textbf{primeiro} quadrante (vértices com $x\le 0$ e $y\ge 0$, ou seja, no segundo quadrante na verdade). Isso não corresponde ao quarto quadrante exigido.
\item Se $b=-4$: os interceptos são $(0,-4)$ e $(2,0)$ — o triângulo tem vértices $(0,0)$, $(2,0)$ e $(0,-4)$, com $x\ge 0$ e $y\le 0$, exatamente no \textbf{quarto quadrante}, como pedido.
\end{itemize}

Logo, a única solução compatível com a condição geométrica é $b=-4$, e
$$g(x) = 2x - 4.$$

Verificação da área: $A=\dfrac{1}{2}\cdot 2\cdot 4 = 4$ $\checkmark$.

\textbf{c) Classificação de $g$}

Como $b=-4\neq 0$, $g$ é uma função afim \textbf{não proporcional}. Geometricamente, a reta $s$ corta o eixo $y$ no ponto $(0,-4)$, diferente da origem — por isso ela não passa pelo ponto $(0,0)$, ao contrário de $r$, que representa a função proporcional $f$.

\bigskip\noindent\textbf{Q034 · múltipla escolha · fácil · aplicar · aplicada · gabarito conferido}\par\nopagebreak
Duas empresas de aplicativo de transporte, A e B, calculam o valor da corrida em função da distância percorrida $x$ (em km). A empresa A cobra R\$\,2,50 por quilômetro rodado, sem nenhuma taxa adicional. A empresa B cobra uma taxa fixa de embarque de R\$\,3,00 mais R\$\,2,00 por quilômetro rodado. Sejam $y_A(x)$ e $y_B(x)$, em reais, os custos totais das corridas das empresas A e B, respectivamente, em função de $x$. Se essas duas funções forem representadas por retas em um plano cartesiano (custo $y$ no eixo vertical e distância $x$ no eixo horizontal), assinale a alternativa que descreve corretamente essas retas.
\begin{enumerate}[label=(\alph*)]
\item A reta que representa a empresa B passa pela origem (0,0), pois seu coeficiente angular é 2; a reta da empresa A corta o eixo y no ponto (0,3), pois seu coeficiente angular é 2,5.
\item Ambas as retas passam pela origem (0,0), já que representam custo por quilômetro rodado; a reta de A tem coeficiente angular 2,5 e a de B tem coeficiente angular 2.
\item A reta que representa a empresa A passa pela origem (0,0) e tem coeficiente angular 2,5, sendo uma função proporcional; a reta da empresa B corta o eixo y no ponto (0,3) e tem coeficiente angular 2, não sendo proporcional.
\item A reta da empresa A corta o eixo y no ponto (0, 2,5); a reta da empresa B corta o eixo y no ponto (0,2) e tem coeficiente angular 3.
\end{enumerate}

\textbf{Gabarito.} A reta que representa a empresa A passa pela origem $(0,0)$ e tem coeficiente angular 2,5 (função proporcional); a reta da empresa B corta o eixo y no ponto $(0,3)$ e tem coeficiente angular 2 (função afim, não proporcional).

\textbf{Resolução.} \textbf{Passo 1 — Escrever as leis das funções.}

Empresa A: cobra só por km rodado, sem taxa fixa: $y_A(x) = 2{,}5x$.

Empresa B: cobra taxa fixa de R\$\,3,00 mais R\$\,2,00 por km: $y_B(x) = 2x + 3$.

\textbf{Passo 2 — Analisar a natureza de cada função.}

$y_A(x) = 2{,}5x$ é da forma $y = ax$ (sem termo constante), logo é uma função \textbf{linear/proporcional}. Seu gráfico é uma reta que passa pela origem $(0,0)$, já que $y_A(0) = 2{,}5 \cdot 0 = 0$. O coeficiente angular é $a = 2{,}5$.

$y_B(x) = 2x+3$ é da forma $y = ax+b$ com $b = 3 \neq 0$, logo é uma função \textbf{afim não proporcional}. Seu gráfico é uma reta que \textbf{não} passa pela origem, pois $y_B(0) = 3 \neq 0$. Ela corta o eixo $y$ no ponto $(0,3)$ e tem coeficiente angular $a = 2$.

\textbf{Passo 3 — Conclusão geométrica.}

A reta de A passa pela origem (comportamento proporcional, coeficiente angular 2,5); a reta de B corta o eixo $y$ em $(0,3)$ (comportamento afim, não proporcional, coeficiente angular 2). As duas retas são crescentes, pois ambos os coeficientes angulares são positivos, mas apenas a de A representa proporcionalidade direta entre custo e distância.

\bigskip\noindent\textbf{Q037 · múltipla escolha · média · entender · teórica · gabarito conferido}\par\nopagebreak
Duas gráficas de impressão, identificadas por R e S, cobram pelo serviço de acordo com o número de folhas impressas ($x$) e o valor total cobrado, em reais ($y$). Usando um aplicativo de geometria dinâmica, uma estudante registrou os seguintes pontos no plano cartesiano:

\begin{itemize}[nosep,leftmargin=*]
\item O gráfico de R passa pelos pontos $A(0,3)$ e $B(4,11)$.
\item O gráfico de S passa pelos pontos $C(2,4)$ e $D(5,10)$.
\end{itemize}

Determine as leis $y = f(x)$ que descrevem R e S e, a partir delas, analise o comportamento geométrico das duas retas no plano cartesiano, comparando-as com a reta de equação $y = 2x$. Assinale a alternativa que classifica corretamente cada gráfico quanto ao tipo de função e apresenta a justificativa geométrica adequada.
\begin{enumerate}[label=(\alph*)]
\item Nem R nem S representam relações de proporcionalidade direta, pois ambas as leis obtidas apresentam termo constante não nulo: $y=2x+3$ para R e $y=2x+2$ para S; portanto, nenhuma das retas passa pela origem.
\item R é uma função afim não proporcional ($y=2x+3$), cujo gráfico é uma reta paralela a $y=2x$, deslocada 3 unidades para cima, não passando pela origem; S é uma função linear (proporcional), pois sua lei é $y=2x$, coincidindo com a própria reta $y=2x$ e passando pela origem do plano cartesiano.
\item R é proporcional, pois seu gráfico foi definido a partir de um ponto com abscissa nula ($x=0$); S não é proporcional, pois nenhum dos pontos fornecidos tem abscissa igual a zero, logo sua reta não passa pela origem.
\item Como R e S têm o mesmo coeficiente angular ($a=2$) da reta $y=2x$, ambas representam relações de proporcionalidade direta, diferindo apenas na posição inicial da reta no plano.
\end{enumerate}

\textbf{Gabarito.} R é uma função afim não proporcional ($y=2x+3$), cujo gráfico é paralelo a $y=2x$ mas deslocado 3 unidades para cima, não passando pela origem; S é uma função linear (proporcional), pois $y=2x$ coincide com a própria reta $y=2x$, passando pela origem.

\textbf{Resolução.} \textbf{Passo 1 — Lei de R.}
O coeficiente angular é $a = \dfrac{11-3}{4-0} = \dfrac{8}{4} = 2$.
Usando o ponto $A(0,3)$: $3 = 2(0) + b \Rightarrow b = 3$.
Logo, $R: y = 2x + 3$.

\textbf{Passo 2 — Lei de S.}
O coeficiente angular é $a = \dfrac{10-4}{5-2} = \dfrac{6}{3} = 2$.
Usando o ponto $C(2,4)$: $4 = 2(2) + b \Rightarrow b = 4 - 4 = 0$.
Logo, $S: y = 2x$.

\textbf{Passo 3 — Comparação geométrica com $y=2x$.}
Ambas as retas têm o mesmo coeficiente angular ($a=2$) da reta $y=2x$, ou seja, têm a mesma inclinação. No entanto:
\begin{itemize}[nosep,leftmargin=*]
\item $R$ tem $b=3 \neq 0$: seu gráfico é uma reta \textbf{paralela} a $y=2x$, deslocada 3 unidades para cima, e \textbf{não passa pela origem} $(0,0)$.
\item $S$ tem $b=0$: sua lei coincide exatamente com $y=2x$, portanto seu gráfico \textbf{passa pela origem}.
\end{itemize}

\textbf{Passo 4 — Classificação.}
Como $S$ tem termo constante nulo e seu gráfico passa pela origem, $S$ representa uma \textbf{função linear (proporcionalidade direta)} entre número de folhas e valor cobrado. Já $R$, apesar de ter a mesma inclinação, possui um termo constante não nulo (uma espécie de "taxa fixa" de R\$\,3), sendo uma \textbf{função afim geral, não proporcional}.

\textbf{Conclusão:} a alternativa correta é a que afirma que R é afim não proporcional (paralela a $y=2x$, deslocada, sem passar pela origem) e S é linear/proporcional (coincide com $y=2x$, passando pela origem).

\bigskip\noindent\textbf{Q040 · múltipla escolha · difícil · entender · aplicada · gabarito conferido}\par\nopagebreak
Duas empresas de transporte por aplicativo, Moto Rápida e Carro Center, cobram uma corrida com base na distância percorrida. A tabela abaixo mostra o valor cobrado, em reais, para algumas distâncias, em quilômetros:

Moto Rápida: 2 km $\to$ R\$\,7,00; 5 km $\to$ R\$\,17,50; 8 km $\to$ R\$\,28,00

Carro Center: 2 km $\to$ R\$\,9,00; 5 km $\to$ R\$\,16,50; 8 km $\to$ R\$\,24,00

Um estudante decide representar, em um mesmo plano cartesiano, o custo da corrida (eixo y, em reais) em função da distância percorrida (eixo x, em km) para cada uma das duas empresas, obtendo duas retas.

Assinale a alternativa que descreve corretamente essas duas retas, indicando qual delas representa uma relação de proporcionalidade direta entre custo e distância, e em que ponto (distância, custo) as duas retas se cruzam.
\begin{enumerate}[label=(\alph*)]
\item Tanto a reta de Moto Rápida (inclinação 3,5) quanto a de Carro Center (inclinação 2,5) passam pela origem, sendo ambas proporcionais; por isso, as retas só têm um ponto em comum, a origem (0; 0).
\item A reta de Moto Rápida passa pela origem, com inclinação 3,5, representando proporcionalidade direta; a reta de Carro Center corta o eixo y em (0; 4) e tem inclinação 2,5; as retas se cruzam no ponto (4; 10).
\item A reta de Moto Rápida passa pela origem, com inclinação 3,5, representando proporcionalidade direta; a reta de Carro Center corta o eixo y em (0; 4) e tem inclinação 2,5; as retas se cruzam no ponto (4; 14).
\item A reta de Carro Center passa pela origem, com inclinação 2,5, representando proporcionalidade direta; a reta de Moto Rápida corta o eixo y em (0; 4) e tem inclinação 3,5; as retas se cruzam no ponto (4; 14).
\end{enumerate}

\textbf{Gabarito.} A reta de Moto Rápida passa pela origem, com inclinação 3,5, representando proporcionalidade direta; a reta de Carro Center corta o eixo y em (0; 4) e tem inclinação 2,5; as retas se cruzam no ponto (4; 14).

\textbf{Resolução.} \textbf{Passo 1 — Taxa de variação (inclinação) de cada reta}

Para a Moto Rápida, usando os pontos $(2;7)$ e $(5;17{,}5)$:
$$a_{MR}=\frac{17{,}5-7}{5-2}=\frac{10{,}5}{3}=3{,}5$$

Para o Carro Center, usando os pontos $(2;9)$ e $(5;16{,}5)$:
$$a_{CC}=\frac{16{,}5-9}{5-2}=\frac{7{,}5}{3}=2{,}5$$

\textbf{Passo 2 — Coeficiente linear (intercepto) de cada reta}

Para a Moto Rápida, testando $y=3{,}5x$ em $x=2$: $3{,}5\cdot 2 = 7$ $\checkmark$ (confere também em $x=5$ e $x=8$). Como não há termo independente, $b_{MR}=0$: a reta \textbf{passa pela origem}, logo $y=3{,}5x$ é uma função \textbf{proporcional} (linear no sentido estrito).

Para o Carro Center, usando $y=2{,}5x+b$ e o ponto $(2;9)$:
$$9=2{,}5\cdot 2+b \Rightarrow b=9-5=4$$
Verificando em $x=8$: $2{,}5\cdot 8+4=24$ $\checkmark$. Assim, $y=2{,}5x+4$: a reta \textbf{corta o eixo y em $(0;4)$}, não passando pela origem — é apenas afim, \textbf{não proporcional}.

\textbf{Passo 3 — Ponto de interseção das retas}

Igualando as duas leis:
$$3{,}5x = 2{,}5x+4$$
$$3{,}5x-2{,}5x=4 \Rightarrow x=4$$
$$y=3{,}5\cdot 4=14$$

As retas se cruzam no ponto $(4;14)$, ou seja, para 4 km de percurso, ambas cobrariam R\$\,14,00.

\textbf{Conclusão:} A reta da Moto Rápida ($y=3{,}5x$) passa pela origem e representa proporcionalidade direta; a do Carro Center ($y=2{,}5x+4$) corta o eixo y em $(0;4)$ e não é proporcional; as retas se cruzam em $(4;14)$.
\subsection*{EM13MAT402}
{\small\itshape Converter representações algébricas de funções polinomiais de 2º grau para representações geométricas no plano cartesiano, distinguindo os casos nos quais uma variável for diretamente proporcional ao quadrado da outra, recorrendo ou não a softwares ou aplicativos de álgebra e geometria dinâmica.}

\bigskip\noindent\textbf{Q047 · discursiva · média · entender · teórica · gabarito conferido}\par\nopagebreak
Considere as três funções quadráticas, definidas para todo $x \in \mathbb{R}$:

$f(x) = 3x^2$

$g(x) = 3x^2 - 12$

$h(x) = 3(x+2)^2$

a) Entre as três funções, identifique qual delas representa uma situação em que a variável $y$ é diretamente proporcional ao quadrado de $x$, isto é, pode ser escrita na forma $y = kx^2$ com $k$ constante. Justifique algebricamente por que as outras duas NÃO se enquadram nessa relação de proporcionalidade direta.

b) Determine as coordenadas do vértice da parábola associada a cada uma das três funções.

c) Descreva, em termos de translações no plano cartesiano (direção e número de unidades deslocadas), como os gráficos de $g$ e de $h$ podem ser obtidos a partir do gráfico de $f$.

\textbf{Gabarito.} Apenas $f(x)=3x^2$ é diretamente proporcional a $x^2$ (vértice em $(0,0)$). Vértices: $f\to(0,0)$, $g\to(0,-12)$, $h\to(-2,0)$. O gráfico de $g$ é o de $f$ deslocado 12 unidades para baixo; o de $h$ é o de $f$ deslocado 2 unidades para a esquerda.

\textbf{Resolução.} \textbf{a) Identificando a proporcionalidade direta}

Dizer que $y$ é diretamente proporcional ao quadrado de $x$ significa que existe uma constante $k$ tal que $y = kx^2$, \textbf{sem nenhum termo aditivo} (nem constante, nem linear).

\begin{itemize}[nosep,leftmargin=*]
\item $f(x) = 3x^2$: já está na forma $y = kx^2$, com $k=3$. \textbf{É} diretamente proporcional a $x^2$.
\item $g(x) = 3x^2 - 12$: ao dividir por $x^2$ obtemos $\dfrac{g(x)}{x^2} = 3 - \dfrac{12}{x^2}$, que \textbf{não é constante}. Logo, $g$ \textbf{não} representa proporcionalidade direta — a constante $-12$ quebra a razão constante entre $y$ e $x^2$.
\item $h(x) = 3(x+2)^2 = 3x^2 + 12x + 12$: expandindo, aparece o termo linear $12x$ e o termo constante $12$. Como há termos de grau menor que 2, $h$ \textbf{também não} é da forma $y=kx^2$.
\end{itemize}

Portanto, apenas $f$ representa uma relação de proporcionalidade direta entre $y$ e $x^2$.

\textbf{b) Vértices}

\begin{itemize}[nosep,leftmargin=*]
\item $f(x) = 3x^2 = 3(x-0)^2 + 0 \Rightarrow$ vértice $(0,0)$.
\item $g(x) = 3x^2 - 12 = 3(x-0)^2 + (-12) \Rightarrow$ vértice $(0,-12)$.
\item $h(x) = 3(x+2)^2 = 3(x-(-2))^2 + 0 \Rightarrow$ vértice $(-2,0)$.
\end{itemize}

\textbf{c) Translações no plano}

Comparando cada função com $f(x)=3x^2$ (que já está na forma canônica $a(x-x_v)^2+y_v$ com vértice na origem):

\begin{itemize}[nosep,leftmargin=*]
\item $g(x) = 3x^2 - 12$: mesmo coeficiente $a=3$, mesma abertura, mas o vértice passou de $(0,0)$ para $(0,-12)$. Logo, o gráfico de $g$ é o gráfico de $f$ \textbf{transladado 12 unidades para baixo}.
\item $h(x) = 3(x+2)^2$: mesmo coeficiente $a=3$, mesma abertura, vértice em $(-2,0)$. Logo, o gráfico de $h$ é o gráfico de $f$ \textbf{transladado 2 unidades para a esquerda}.
\end{itemize}

Geometricamente, apenas o gráfico de $f$ tem vértice na origem — condição necessária (mas equivalente, junto com a ausência de termos de grau menor na forma algébrica) para que a parábola represente $y$ diretamente proporcional a $x^2$; os deslocamentos de $g$ e $h$ tiram justamente essa propriedade.

\bigskip\noindent\textbf{Q055 · múltipla escolha · média · aplicar · aplicada · gabarito conferido}\par\nopagebreak
Em um teste de frenagem, a distância de frenagem $d$ (em metros) de um carrinho de brinquedo, em função de sua velocidade $v$ (em m/s, com $v \geq 0$), foi modelada por $d(v) = 0{,}5v^2$.

Em outro experimento, a altura $h$ (em metros) de uma bolinha lançada verticalmente para cima a partir de uma plataforma, em função do tempo $t$ (em segundos, com $t \geq 0$), foi modelada por $h(t) = -5t^2 + 20t + 1$.

Um estudante deseja esboçar, sem construir tabela de valores, o gráfico no plano cartesiano correspondente à situação em que a grandeza dependente é diretamente proporcional ao quadrado da grandeza independente.

Assinale a alternativa que descreve corretamente esse gráfico.
\begin{enumerate}[label=(\alph*)]
\item Parábola com concavidade voltada para cima, vértice no ponto $(0; 0{,}5)$, deslocada verticalmente para cima em relação à origem.
\item Parábola com concavidade voltada para cima, vértice na origem $(0,0)$, simétrica em relação ao eixo vertical e situada no primeiro quadrante para $v \geq 0$, pois $d(v)=0{,}5v^2$ é diretamente proporcional ao quadrado de $v$.
\item Reta crescente passando pela origem, com coeficiente angular $0{,}5$.
\item Parábola com concavidade voltada para baixo, vértice no ponto $(2,21)$, interceptando o eixo vertical em $(0,1)$.
\end{enumerate}

\textbf{Gabarito.} B

\textbf{Resolução.} \textbf{Passo 1 — Identificar a proporcionalidade direta ao quadrado.}

Uma grandeza $y$ é diretamente proporcional ao quadrado de $x$ quando $y = kx^2$, ou seja, quando a expressão \textbf{não possui} termo linear ($bx$) nem termo constante ($c$).

\begin{itemize}[nosep,leftmargin=*]
\item $d(v) = 0{,}5v^2$: aqui $b=0$ e $c=0$. Logo, $d$ \textbf{é} diretamente proporcional ao quadrado de $v$, com constante $k=0{,}5$.
\item $h(t) = -5t^2+20t+1$: aqui $b=20\neq0$ e $c=1\neq0$. Logo, $h$ \textbf{não é} diretamente proporcional ao quadrado de $t$.
\end{itemize}

Portanto, o gráfico pedido é o de $d(v)=0{,}5v^2$.

\textbf{Passo 2 — Converter a lei algébrica em características geométricas.}

\begin{itemize}[nosep,leftmargin=*]
\item Coeficiente do termo quadrático: $a=0{,}5>0$, logo a parábola tem \textbf{concavidade voltada para cima}.
\item Vértice: $x_v = -\dfrac{b}{2a} = -\dfrac{0}{2(0{,}5)} = 0$ e $y_v = d(0) = 0$. Logo o vértice é o ponto $(0,0)$ — a \textbf{origem} do plano cartesiano, o que é característico de toda função da forma $y=kx^2$ (sem deslocamentos).
\item Como o domínio físico exige $v\geq 0$, o gráfico corresponde apenas ao ramo direito da parábola, situado no primeiro quadrante (incluindo a origem).
\end{itemize}

\textbf{Passo 3 — Comparar com o gráfico de $h(t)$ (para contraste).}

Para $h(t)=-5t^2+20t+1$: $a=-5<0$ (concavidade para baixo), vértice em $t_v = -\dfrac{20}{2(-5)} = 2$, $h(2) = -20+40+1 = 21$, ou seja, vértice em $(2,21)$, e intersecção com o eixo vertical em $(0,1)$. Esse gráfico \textbf{não} representa uma proporcionalidade direta ao quadrado, pois o vértice não está na origem.

\textbf{Conclusão:} o gráfico correto é uma parábola com concavidade para cima, vértice na origem, restrita ao primeiro quadrante — alternativa (B).

\bigskip\noindent\textbf{Q068 · múltipla escolha · fácil · aplicar · teórica · gabarito conferido}\par\nopagebreak
Considere as quatro funções quadráticas a seguir, definidas para $x \in \mathbb{R}$:

$A) \; y = 3x^2$

$B) \; y = 3x^2 - 6$

$C) \; y = 3(x-2)^2$

$D) \; y = 3x^2 + 5x$

Para cada uma delas, determine algebricamente as coordenadas do vértice da parábola correspondente e, a partir dessa informação geométrica, identifique qual das quatro funções é tal que $y$ é diretamente proporcional a $x^2$ (isto é, pode ser escrita na forma $y = kx^2$). Justifique sua escolha relacionando a posição do vértice de cada parábola no plano cartesiano com a condição de proporcionalidade direta.
\begin{enumerate}[label=(\alph*)]
\item A função $y = 3x^2 + 5x$, pois seu gráfico passa pela origem (quando $x=0$, $y=0$), o que basta para garantir que $y$ é proporcional a $x^2$.
\item A função $y = 3(x-2)^2$, pois seu gráfico toca o eixo $x$ em um único ponto, o que caracteriza a proporcionalidade direta com $x^2$.
\item A função $y = 3x^2$, pois seu gráfico é uma parábola cujo vértice está exatamente em $(0,0)$, coincidindo com a origem do plano cartesiano.
\item A função $y = 3x^2 - 6$, pois seu vértice está sobre o eixo $y$ (em $x=0$), o que já garante a proporcionalidade direta entre $y$ e $x^2$.
\end{enumerate}

\textbf{Gabarito.} C) y = 3x², pois é a única cujo vértice está exatamente na origem (0,0) do plano cartesiano.

\textbf{Resolução.} \textbf{Condição geométrica-algébrica.} Uma função é da forma $y = kx^2$ (proporcionalidade direta entre $y$ e $x^2$) se, e somente se, não houver termo linear nem termo constante — o que equivale, geometricamente, a a parábola ter seu \textbf{vértice exatamente na origem} $(0,0)$ do plano cartesiano.

\textbf{Passo 1 — Vértice de A: $y=3x^2$.}
Aqui $a=3,\,b=0,\,c=0$. Logo $x_v = -\dfrac{b}{2a} = 0$ e $y_v = f(0) = 0$.
Vértice: $(0,0)$ — coincide com a origem.

\textbf{Passo 2 — Vértice de B: $y=3x^2-6$.}
$a=3,\,b=0,\,c=-6$. $x_v = 0$, $y_v = f(0) = -6$.
Vértice: $(0,-6)$ — está sobre o eixo $y$, mas \textbf{não} na origem (a parábola é a de A deslocada 6 unidades para baixo).

\textbf{Passo 3 — Vértice de C: $y=3(x-2)^2$.}
Expandindo: $y = 3x^2 - 12x + 12$, então $a=3,\,b=-12,\,c=12$.
$x_v = -\dfrac{-12}{2\cdot 3} = 2$, e $y_v = f(2) = 3(2-2)^2 = 0$.
Vértice: $(2,0)$ — a parábola toca o eixo $x$, mas em $x=2$, não na origem.

\textbf{Passo 4 — Vértice de D: $y=3x^2+5x$.}
$a=3,\,b=5,\,c=0$. $x_v = -\dfrac{5}{6}$.
$y_v = f\left(-\dfrac{5}{6}\right) = 3\cdot\dfrac{25}{36} + 5\cdot\left(-\dfrac{5}{6}\right) = \dfrac{25}{12} - \dfrac{25}{6} = -\dfrac{25}{12}$.
Vértice: $\left(-\dfrac{5}{6}, -\dfrac{25}{12}\right)$.
Observe que $f(0)=0$, ou seja, o gráfico de D *passa* pela origem como um ponto qualquer da curva, mas seu \textbf{vértice não está ali} — a presença do termo $5x$ desloca o vértice para fora da origem, o que confirma que $y$ não é proporcional a $x^2$ nesse caso.

\textbf{Conclusão.} Comparando os quatro vértices — $(0,0)$, $(0,-6)$, $(2,0)$ e $\left(-\frac56,-\frac{25}{12}\right)$ — apenas a parábola de \textbf{A} tem vértice exatamente na origem do plano cartesiano. Isso corresponde exatamente à condição algébrica $b=0$ e $c=0$, ou seja, $y=3x^2$, que é a única entre as quatro em que $y$ é diretamente proporcional a $x^2$.

\bigskip\noindent\textbf{Q081 · discursiva · difícil · aplicar · aplicada · gabarito conferido}\par\nopagebreak
Um professor de Física e um analista financeiro registraram duas situações reais.

\textbf{Situação 1 — Queda livre}
Um objeto é solto (sem velocidade inicial) do alto de um prédio. A distância percorrida $d$ (em metros) foi medida em função do tempo $t$ (em segundos):

\begin{center}
\begin{tabular}{|p{0.215\linewidth}|p{0.215\linewidth}|p{0.215\linewidth}|p{0.215\linewidth}|}
\hline
\textbf{$t$ (s)} & \textbf{1} & \textbf{2} & \textbf{3} \\
\hline
$d$ (m) & 5 & 20 & 45 \\
\hline
\end{tabular}
\end{center}

\textbf{Situação 2 — Lucro de uma loja}
Uma loja varia o preço unitário $p$ (em reais) de um produto e registra o lucro mensal $L$ (em milhares de reais):

\begin{center}
\begin{tabular}{|p{0.215\linewidth}|p{0.215\linewidth}|p{0.215\linewidth}|p{0.215\linewidth}|}
\hline
\textbf{$p$ (R\textbackslash\{\}\$)} & \textbf{10} & \textbf{20} & \textbf{30} \\
\hline
$L$ (mil R\textbackslash\{\}\$) & 800 & 1200 & 800 \\
\hline
\end{tabular}
\end{center}

Sabendo que, em ambos os casos, a grandeza da segunda linha é uma função polinomial do 2º grau da grandeza da primeira linha, resolva:

a) Determine a lei algébrica $d(t)$ e a lei algébrica $L(p)$ que descrevem cada situação.

b) Para cada função, identifique no plano cartesiano: as coordenadas do vértice da parábola, o sentido da concavidade e os pontos em que o gráfico intercepta os eixos coordenados.

c) Apenas uma das duas situações representa uma grandeza \textbf{diretamente proporcional ao quadrado} da outra. Diga qual é e justifique sua resposta usando dois argumentos: um algébrico (a forma da lei obtida em (a)) e um geométrico (a posição do gráfico em relação à origem do plano cartesiano).

\textbf{Gabarito.} $d(t)=5t^2$, vértice $(0,0)$, concavidade para cima, parábola passando pela origem — diretamente proporcional a $t^2$. $L(p)=-4p^2+160p-400$, vértice $(20,1200)$, concavidade para baixo, intercepta o eixo $L$ em $(0,-400)$ e o eixo $p$ em $p=20\pm10\sqrt3$ — quadrática geral, não é proporcional ao quadrado de $p$. A situação diretamente proporcional ao quadrado é a Situação 1 (queda livre).

\textbf{Resolução.} \textbf{Passo 1 — Montar o sistema para $d(t) = at^2+bt+c$}

Usando os pontos $(1,5)$, $(2,20)$ e $(3,45)$:

$a+b+c=5$
$4a+2b+c=20$
$9a+3b+c=45$

Subtraindo a 1ª da 2ª: $3a+b=15$. Subtraindo a 2ª da 3ª: $5a+b=25$. Subtraindo essas duas: $2a=10 \Rightarrow a=5$. Logo $b=15-3(5)=0$ e $c=5-5-0=0$.

$$d(t)=5t^2$$

\textbf{Passo 2 — Montar o sistema para $L(p) = ap^2+bp+c$}

Usando os pontos $(10,800)$, $(20,1200)$ e $(30,800)$:

$100a+10b+c=800$
$400a+20b+c=1200$
$900a+30b+c=800$

Subtraindo a 1ª da 2ª: $300a+10b=400 \Rightarrow 30a+b=40$.
Subtraindo a 2ª da 3ª: $500a+10b=-400 \Rightarrow 50a+b=-40$.
Subtraindo essas duas: $20a=-80 \Rightarrow a=-4$. Logo $b=40-30(-4)=160$ e $c=800-100(-4)-10(160)=800+400-1600=-400$.

$$L(p)=-4p^2+160p-400$$

Verificação: $L(10)=-400+1600-400=800$; $L(20)=-1600+3200-400=1200$; $L(30)=-3600+4800-400=800$. $\checkmark$

\textbf{Passo 3 — Elementos geométricos de $d(t)=5t^2$}

Como $b=c=0$, o vértice é $t_v=-\dfrac{b}{2a}=0$, $d(0)=0$, ou seja, vértice em $(0,0)$. Como $a=5>0$, a parábola tem concavidade voltada para cima. O gráfico passa pela origem e não intercepta o eixo $d$ em outro ponto nem o eixo $t$ em outro ponto (raiz dupla $t=0$).

\textbf{Passo 4 — Elementos geométricos de $L(p)=-4p^2+160p-400$}

Vértice: $p_v=-\dfrac{160}{2(-4)}=20$, $L(20)=-1600+3200-400=1200$, ou seja, vértice em $(20,1200)$. Como $a=-4<0$, a concavidade é voltada para baixo. Intersecção com o eixo $L$: $p=0 \Rightarrow L=-400$, ponto $(0,-400)$. Intersecções com o eixo $p$ (fazendo $L=0$): $-4p^2+160p-400=0 \Rightarrow p^2-40p+100=0 \Rightarrow p = 20\pm\sqrt{300} = 20\pm10\sqrt3$.

\textbf{Passo 5 — Distinguindo a proporcionalidade direta ao quadrado}

Uma grandeza $y$ é diretamente proporcional ao quadrado de $x$ quando $y=kx^2$, isto é, quando na forma $ax^2+bx+c$ temos $b=0$ e $c=0$. Isso ocorre apenas em $d(t)=5t^2$ (com $k=5$). Geometricamente, isso corresponde a uma parábola cujo vértice está exatamente na origem do plano cartesiano, o que é o caso de $d(t)$.

Já $L(p)=-4p^2+160p-400$ tem $b\neq 0$ e $c\neq 0$: é uma função quadrática geral, cujo vértice $(20,1200)$ está deslocado da origem — portanto $L$ \textbf{não} é diretamente proporcional a $p^2$.

\textbf{Conclusão:} apenas a Situação 1 (queda livre) representa uma grandeza diretamente proporcional ao quadrado da outra.

\bigskip\noindent\textbf{Q089 · múltipla escolha · difícil · entender · teórica · gabarito conferido}\par\nopagebreak
Uma função quadrática $y = ax^2 + bx + c$ (com $a \neq 0$) representa uma relação em que $y$ é diretamente proporcional ao quadrado de $x$ quando existe uma constante $k \neq 0$ tal que $y = kx^2$ para todo $x$ real — ou seja, quando $b = 0$ e $c = 0$ simultaneamente.

Considere quatro parábolas, todas com concavidade voltada para cima, esboçadas no plano cartesiano e descritas apenas por características geométricas de seus gráficos:

\begin{itemize}[nosep,leftmargin=*]
\item \textbf{Parábola A}: intercepta o eixo $x$ nos pontos $(0,0)$ e $(4,0)$, e passa pelo ponto $(1,-3)$.
\item \textbf{Parábola B}: tangencia o eixo $x$ exatamente no ponto $(0,0)$ — esse ponto é o vértice da parábola — e passa pelo ponto $(3,18)$.
\item \textbf{Parábola C}: intercepta o eixo $x$ nos pontos $(-3,0)$ e $(3,0)$, e passa pelo ponto $(1,-8)$.
\item \textbf{Parábola D}: tangencia o eixo $x$ exatamente no ponto $(2,0)$ — esse ponto é o vértice da parábola — e passa pelo ponto $(0,8)$.
\end{itemize}

Assinale a alternativa que identifica corretamente qual das parábolas descritas representa uma função em que $y$ é diretamente proporcional ao quadrado de $x$.
\begin{enumerate}[label=(\alph*)]
\item Parábola C, pois seu gráfico é simétrico em relação ao eixo y.
\item Parábola A, pois seu gráfico passa pela origem do plano cartesiano.
\item Parábola D, pois seu gráfico tangencia o eixo x em um único ponto.
\item Parábola B, pois seu vértice está na origem e ela não possui termo linear nem constante.
\end{enumerate}

\textbf{Gabarito.} Alternativa (d): Parábola B, pois $y=2x^2$, com vértice na origem, $b=0$ e $c=0$.

\textbf{Resolução.} \textbf{Ideia-chave:} $y$ é diretamente proporcional a $x^2$ somente quando a equação da parábola é $y=kx^2$, isto é, quando $b=0$ \textbf{e} $c=0$ ao mesmo tempo — geometricamente, isso significa que o \textbf{vértice está exatamente na origem}. Passar pela origem, ou ser simétrica em relação ao eixo $y$, ou tangenciar o eixo $x$ em um único ponto \textbf{não bastam isoladamente}.

\textbf{Parábola A} — raízes em $x=0$ e $x=4$: $y=a\,x(x-4)$. Usando $(1,-3)$: $a(1)(1-4)=-3a=-3 \Rightarrow a=1$. Logo $y=x^2-4x$, com $b=-4\neq 0$. Embora passe pela origem, \textbf{não} é proporcional (a outra raiz, $x=4$, mostra que o vértice não está em $(0,0)$).

\textbf{Parábola B} — vértice (tangência) em $(0,0)$, logo raiz dupla em $x=0$: $y=a x^2$. Usando $(3,18)$: $9a=18 \Rightarrow a=2$. Logo $y=2x^2$, com $b=0$ e $c=0$. \textbf{É} diretamente proporcional, com $k=2$.

\textbf{Parábola C} — raízes em $x=-3$ e $x=3$: $y=a(x-3)(x+3)=ax^2-9a$. Usando $(1,-8)$: $a-9a=-8a=-8 \Rightarrow a=1$. Logo $y=x^2-9$: aqui $b=0$ (é simétrica em relação ao eixo $y$), mas $c=-9\neq 0$, então o vértice é $(0,-9)$, \textbf{não} a origem. Não é proporcional.

\textbf{Parábola D} — vértice (tangência) em $(2,0)$: $y=a(x-2)^2$. Usando $(0,8)$: $4a=8 \Rightarrow a=2$. Logo $y=2x^2-8x+8$, com $b=-8\neq0$ e $c=8\neq0$. Apesar de tangenciar o eixo $x$ em um único ponto (raiz dupla), essa tangência ocorre em $(2,0)$ e não na origem, então \textbf{não} é proporcional.

\textbf{Conclusão:} apenas a Parábola B tem simultaneamente $b=0$ e $c=0$ (vértice na origem), logo é a única cuja lei é $y=kx^2$, isto é, $y$ diretamente proporcional a $x^2$.
\subsection*{EM13MAT403}
{\small\itshape Comparar e analisar as representações, em plano cartesiano, das funções exponencial e logarítmica para identificar as características fundamentais (domínio, imagem, crescimento) de cada uma, com ou sem apoio de tecnologias digitais, estabelecendo relações entre elas.}

\bigskip\noindent\textbf{Q013 · discursiva · fácil · analisar · teórica · conferido em parte}\par\nopagebreak
Um professor propõe que os alunos investiguem, usando um software gráfico (ou construindo os gráficos manualmente em papel milimetrado), as funções $f(x) = 2^x$ e $g(x) = \log_2 x$.

a) Determine o domínio e o conjunto-imagem de $f$ e de $g$. Compare esses dois conjuntos (domínio de $f$ com imagem de $g$, e domínio de $g$ com imagem de $f$) e explique, matematicamente, por que essa relação ocorre.

b) Calcule $f(3)$ e $g(8)$; calcule também $f(0)$ e $g(1)$. Observando os pares de valores obtidos, descreva que tipo de simetria os gráficos de $f$ e $g$ apresentam no plano cartesiano.

c) Classifique $f$ e $g$ quanto ao crescimento (crescente ou decrescente) em todo o domínio de cada uma. Em seguida, utilizando os valores calculados no item (b), explique por que o fato de $f$ e $g$ serem ambas crescentes é compatível com a simetria identificada — ou seja, justifique por que essa simetria não transforma o crescimento de uma função em decrescimento na outra.

\textbf{Gabarito.} a) $D(f)=\mathbb{R}=Im(g)$ e $D(g)=(0,+\infty)=Im(f)$, pois $g$ é a inversa de $f$. b) $f(3)=8, g(8)=3, f(0)=1, g(1)=0$: os gráficos de $f$ e $g$ são simétricos em relação à reta $y=x$. c) $f$ e $g$ são ambas crescentes; a simetria por $y=x$ preserva a ordem dos valores, por isso não inverte o crescimento.

\textbf{Resolução.} \textbf{a) Domínio e imagem}

Para $f(x) = 2^x$: como a base $2>0$ e $2 \neq 1$, a expressão $2^x$ está definida para todo $x$ real e sempre resulta em um número positivo. Logo:
$$D(f) = \mathbb{R}, \qquad Im(f) = (0, +\infty)$$

Para $g(x) = \log_2 x$: o logaritmo só é definido para argumentos positivos, e pode resultar em qualquer número real. Logo:
$$D(g) = (0, +\infty), \qquad Im(g) = \mathbb{R}$$

Comparando: $D(f) = Im(g) = \mathbb{R}$ e $D(g) = Im(f) = (0,+\infty)$. Essa troca ocorre porque $g$ é a \textbf{função inversa} de $f$ (pois $g(x)=\log_2 x$ desfaz a operação $f(x)=2^x$, já que $\log_2(2^x)=x$). Para duas funções inversas entre si, o domínio de uma é sempre igual à imagem da outra, pois inverter uma função significa trocar o papel de $x$ e $y$.

\textbf{b) Valores e simetria}

$f(3) = 2^3 = 8$ e $g(8) = \log_2 8 = 3$. O ponto $(3,8)$ pertence ao gráfico de $f$, e o ponto $(8,3)$ pertence ao gráfico de $g$ — as coordenadas estão trocadas.

$f(0) = 2^0 = 1$ e $g(1) = \log_2 1 = 0$. O ponto $(0,1)$ pertence a $f$, e o ponto $(1,0)$ pertence a $g$ — novamente coordenadas trocadas.

Como, para todo ponto $(a,b)$ do gráfico de $f$, o ponto $(b,a)$ pertence ao gráfico de $g$, os dois gráficos são \textbf{simétricos em relação à reta $y = x$} (reflexão que troca as coordenadas de cada ponto).

\textbf{c) Crescimento e compatibilidade com a simetria}

$f(x)=2^x$ é crescente em $\mathbb{R}$ (base $2>1$: quanto maior $x$, maior $2^x$).

$g(x)=\log_2 x$ é crescente em $(0,+\infty)$ (base $2>1$: quanto maior $x$, maior $\log_2 x$).

Com os valores do item (b): $f(0)=1 < f(3)=8$ mostra que $f$ cresce de $x=0$ para $x=3$. Ao refletir esses pontos trocando as coordenadas, obtemos $g(1)=0 < g(8)=3$, ou seja, $g$ também cresce de $x=1$ para $x=8$. A ordem entre os valores foi preservada na troca de coordenadas — isso acontece porque a reflexão em torno da reta $y=x$ preserva a relação de ordem quando a função original é crescente (se $x_1<x_2 \Rightarrow f(x_1)<f(x_2)$, então trocando as coordenadas, $f(x_1) < f(x_2)$ vira exatamente a condição de que $g$ também cresce nesse intervalo). Por isso, a simetria em relação a $y=x$ \textbf{não inverte o sentido de crescimento}: uma função crescente tem inversa também crescente (o mesmo não ocorreria se $f$ fosse decrescente, caso em que $g$ também seria decrescente).

\bigskip\noindent\textbf{Q015 · discursiva · média · analisar · aplicada · gabarito conferido}\par\nopagebreak
Uma colônia de bactérias em um experimento de laboratório tem seu número de indivíduos, a partir do início da observação, dado por $N(t) = 100 \cdot 2^{t}$, em que $t$ é o tempo em horas ($t \geq 0$) e $N(t)$ é o número de bactérias. Um pesquisador quer saber, ao contrário, quanto tempo é necessário para que a colônia atinja um determinado número $N$ de bactérias, e por isso utiliza a função inversa de $N(t)$, que ele chama de $t(N)$.

a) Determine o domínio e a imagem de $N(t)$, considerando o contexto do experimento.

b) Encontre a expressão de $t(N)$ (a função inversa de $N(t)$) e determine seu domínio e sua imagem.

c) Compare o crescimento das duas funções: as duas são crescentes, mas uma cresce 'cada vez mais rápido' e a outra cresce 'cada vez mais devagar' à medida que a variável aumenta. Identifique qual é qual e justifique observando como varia o incremento de cada função.

d) Explique por que o domínio de $N(t)$ coincide com a imagem de $t(N)$, e a imagem de $N(t)$ coincide com o domínio de $t(N)$, relacionando esse fato com a posição dos gráficos das duas funções no plano cartesiano.

\textbf{Gabarito.} a) $D_N=[0,+\infty)$, $Im_N=[100,+\infty)$. b) $t(N)=\log_2(N/100)$, com $D_t=[100,+\infty)$, $Im_t=[0,+\infty)$. c) $N(t)$ cresce cada vez mais rápido (convexa); $t(N)$ cresce cada vez mais devagar (côncava). d) Por serem funções inversas, seus gráficos são simétricos em relação à reta $y=x$, o que troca domínio e imagem entre as duas funções.

\textbf{Resolução.} \textbf{a) Domínio e imagem de $N(t) = 100\cdot 2^t$}

Como $t$ representa tempo a partir do início da observação, $t \geq 0$. Logo:
$$D_N = [0, +\infty)$$

Quando $t=0$, $N(0) = 100 \cdot 2^0 = 100$, o menor valor possível. Como $2^t$ é crescente e ilimitada quando $t \to +\infty$, $N(t)$ também cresce sem limite. Logo:
$$Im_N = [100, +\infty)$$

\textbf{b) Função inversa $t(N)$}

De $N = 100\cdot 2^t$, isolamos $t$:
$$\frac{N}{100} = 2^t \implies t = \log_2\left(\frac{N}{100}\right)$$

Como funções inversas trocam domínio e imagem entre si:
$$D_t = Im_N = [100, +\infty), \qquad Im_t = D_N = [0, +\infty)$$

(De fato, o menor $N$ possível no experimento é 100, e para $N=100$ temos $t=\log_2(1)=0$, consistente.)

\textbf{c) Comparando o crescimento}

Ambas as funções são estritamente crescentes, pois a base $2$ é maior que $1$ tanto na exponencial quanto no logaritmo.

Porém, o \textbf{modo} de crescer é diferente. Observe os incrementos de $N(t)$ a cada hora:
\begin{itemize}[nosep,leftmargin=*]
\item de $t=0$ a $t=1$: $N$ vai de $100$ a $200$ (aumenta 100)
\item de $t=1$ a $t=2$: $N$ vai de $200$ a $400$ (aumenta 200)
\item de $t=2$ a $t=3$: $N$ vai de $400$ a $800$ (aumenta 400)
\end{itemize}

Ou seja, o incremento absoluto de $N(t)$ \textbf{aumenta} a cada intervalo — a função exponencial cresce cada vez mais rápido (é convexa, sua taxa de variação aumenta).

Já para $t(N)$: cada vez que $N$ \textbf{dobra}, $t$ aumenta sempre a mesma quantidade fixa, 1 hora (pois $\log_2(2N/100) - \log_2(N/100) = \log_2 2 = 1$). Mas para produzir esse mesmo aumento de 1 unidade em $t$, é preciso um incremento absoluto de $N$ cada vez maior (de $100\to200$, depois $200\to400$, depois $400\to800$...). Isso significa que, para incrementos iguais de $N$, o aumento de $t$ fica cada vez menor — a função logarítmica cresce cada vez mais devagar (é côncava, sua taxa de variação diminui).

\textbf{d) Relação entre domínio e imagem das duas funções}

Como $t(N)$ é a função inversa de $N(t)$, seus gráficos são simétricos em relação à reta $y=x$ no plano cartesiano. Essa simetria troca os eixos: o que era eixo do domínio (horizontal) de uma passa a ser o eixo da imagem (vertical) da outra. Por isso, necessariamente:
$$D_{N} = Im_{t} = [0,+\infty) \qquad \text{e} \qquad Im_{N} = D_{t} = [100,+\infty)$$

Essa é uma propriedade geral de qualquer par de funções inversas, e em particular da relação entre exponencial e logaritmo de mesma base: uma 'desfaz' a outra, e seus papéis de domínio e imagem se invertem.

\bigskip\noindent\textbf{Q017 · discursiva · difícil · analisar · teórica · conferido em parte}\par\nopagebreak
Um biólogo estuda o crescimento de uma colônia de bactérias em laboratório. Ele constata que, em um modelo matemático simplificado, a quantidade de indivíduos na colônia, em milhares, pode ser descrita pela função $f(x) = 2^{x}$, em que $x$ representa o número de horas decorridas desde um instante de referência (podendo $x$ assumir valores negativos, representando horas anteriores a esse instante, dentro do intervalo em que o modelo é válido).

Para resolver o problema inverso — estimar quantas horas se passaram até que a colônia atingisse uma determinada quantidade $y$ de milhares de indivíduos — o biólogo utiliza a função $g(x) = \log_{2}(x)$.

Com base nessas duas funções, resolva:

\textbf{a)} Determine o domínio e a imagem de $f$ e de $g$. Compare os dois pares de conjuntos e explique, em termos matemáticos, por que essa relação entre eles não é uma coincidência.

\textbf{b)} Analise o crescimento de $f$ e de $g$ (cada uma é crescente ou decrescente?). Embora ambas cresçam no mesmo sentido, explique, sem usar tabelas numéricas, por que o crescimento de $f$ é qualitativamente muito mais rápido que o de $g$ à medida que $x$ aumenta.

\textbf{c)} Calcule as composições $f(g(x))$ e $g(f(x))$, indicando para quais valores de $x$ cada composição está definida. O que esses resultados revelam sobre a relação entre as funções $f$ e $g$?

\textbf{d)} Sabendo que $f(1) = 2$ e $f(3) = 8$, calcule $g(2)$ e $g(8)$. Em seguida, esboce em um mesmo plano cartesiano os gráficos de $f$, de $g$ e da reta $y = x$, situando nele os pontos $(1,2)$, $(2,1)$, $(3,8)$ e $(8,3)$. Explique, a partir desse esboço, que relação geométrica existe entre o gráfico de $f$ e o gráfico de $g$, e como as respostas dos itens (a), (b) e (c) justificam essa relação.

\textbf{Gabarito.} a) $D_f=\mathbb{R}$, $Im_f=(0,+\infty)$; $D_g=(0,+\infty)$, $Im_g=\mathbb{R}$ — domínio e imagem se invertem entre f e g. b) f e g são ambas crescentes, mas f cresce de forma multiplicativa (cada vez mais rápido) e g cresce de forma cada vez mais lenta. c) $f(g(x))=x$ para $x>0$ e $g(f(x))=x$ para todo x real, mostrando que g é a função inversa de f. d) $g(2)=1$, $g(8)=3$; os gráficos de f e g são simétricos em relação à reta $y=x$.

\textbf{Resolução.} \textbf{a) Domínio e imagem}

Para $f(x) = 2^{x}$: como toda potência de base positiva está definida para qualquer expoente real, o domínio de $f$ é $D_f = \mathbb{R}$. Como $2^x > 0$ para todo $x$, e $f$ assume todos os valores positivos, a imagem é $Im_f = (0, +\infty)$.

Para $g(x) = \log_2(x)$: o logaritmo só está definido para argumentos positivos, logo $D_g = (0, +\infty)$. Como o logaritmo assume qualquer valor real quando o argumento varia em $(0,+\infty)$, temos $Im_g = \mathbb{R}$.

Comparando: $D_f = Im_g = \mathbb{R}$ e $Im_f = D_g = (0,+\infty)$. Essa troca exata entre domínio e imagem não é coincidência: ela é a marca registrada de duas funções que são \textbf{inversas uma da outra} — o domínio de uma função sempre coincide com a imagem de sua inversa, e vice-versa.

\textbf{b) Crescimento}

Como a base $2 > 1$, $f(x) = 2^x$ é \textbf{crescente} em todo o seu domínio, e $g(x) = \log_2(x)$ também é \textbf{crescente} em todo o seu domínio (base do logaritmo maior que 1).

Apesar de ambas serem crescentes, o crescimento é qualitativamente diferente: em $f$, cada aumento de uma unidade em $x$ \textbf{multiplica} o valor da função por 2 (crescimento multiplicativo/exponencial), fazendo $f$ crescer cada vez mais depressa. Em $g$, para que o valor da função aumente uma unidade é preciso que o argumento $x$ seja \textbf{multiplicado} por 2 (pois $\log_2(2x) = \log_2(x)+1$); ou seja, é preciso um acréscimo cada vez maior em $x$ para obter o mesmo ganho em $g$. Por isso $g$ cresce cada vez mais devagar, enquanto $f$ cresce cada vez mais rápido.

\textbf{c) Composições}

$f(g(x)) = 2^{\log_2(x)} = x$, definida para $x \in (0,+\infty)$ (domínio de $g$).

$g(f(x)) = \log_2(2^{x}) = x$, definida para todo $x \in \mathbb{R}$ (domínio de $f$).

Como aplicar $g$ depois de $f$ (e $f$ depois de $g$) devolve sempre o valor original de $x$, isso mostra que $f$ e $g$ \textbf{desfazem uma a ação da outra}: são funções inversas, $g = f^{-1}$ (e $f = g^{-1}$).

\textbf{d) Pontos e gráfico}

Como $f(1)=2$, temos $g(2) = 1$ (pois $g$ desfaz o que $f$ fez). Como $f(3)=8$, temos $g(8) = 3$.

Observe que o ponto $(1,2)$ de $f$ corresponde ao ponto $(2,1)$ de $g$ — as coordenadas trocaram de posição. O mesmo ocorre com $(3,8)$ e $(8,3)$. Trocar as coordenadas $(a,b) \to (b,a)$ é exatamente a operação de \textbf{refletir um ponto em torno da reta $y=x$}.

No esboço, a curva $y = 2^x$ parte de valores próximos de $0$ (para $x$ muito negativo), passa por $(0,1)$, $(1,2)$ e $(3,8)$, subindo cada vez mais rápido. A curva $y=\log_2(x)$ passa por $(1,0)$, $(2,1)$ e $(8,3)$, subindo cada vez mais devagar, com uma assíntota vertical em $x=0$. As duas curvas são simétricas uma da outra em relação à reta $y=x$, que corta o plano na diagonal.

Essa simetria gráfica é a representação geométrica exata do que foi mostrado algebricamente: como $D_f = Im_g$, $Im_f = D_g$ (item a) e $f(g(x))=g(f(x))=x$ (item c), $f$ e $g$ são funções inversas, e por isso seus gráficos são reflexos um do outro através da reta $y=x$. O item (b) explica por que, apesar dessa simetria, uma curva ($f$) dispara para cima muito mais rapidamente que a outra ($g$) se afasta lentamente do eixo — reflexo do fato de que uma é exponencial e a outra, logarítmica.

\bigskip\noindent\textbf{Q019 · múltipla escolha · fácil · entender · aplicada · conferido em parte}\par\nopagebreak
Um biólogo estuda o crescimento de uma cultura de bactérias em laboratório. O número de bactérias, em milhares, após $t$ horas (com $t$ podendo assumir qualquer valor real, inclusive negativo, para fins de modelagem matemática) é dado por $f(t) = 2^t$. Para descobrir quanto tempo é necessário para que a cultura atinja uma quantidade $N$ (em milhares) de bactérias, o biólogo utiliza a função inversa $g(N) = \log_2(N)$. As duas funções são representadas no mesmo plano cartesiano. Com base na análise do domínio, da imagem e do comportamento de crescimento de $f$ e $g$, assinale a alternativa correta.
\begin{enumerate}[label=(\alph*)]
\item Como f e g são inversas, ambas têm exatamente o mesmo domínio (os reais) e a mesma imagem (os reais positivos), sendo ambas crescentes.
\item O domínio de f é o conjunto dos reais positivos e a imagem de g é o conjunto de todos os reais, sendo f decrescente e g crescente.
\item O domínio de f coincide com a imagem de g, mas g é decrescente enquanto f é crescente.
\item O domínio de f (todos os reais) é igual à imagem de g, e a imagem de f (reais positivos) é igual ao domínio de g; ambas as funções são crescentes em seus domínios.
\end{enumerate}

\textbf{Gabarito.} O domínio de f (todos os reais) é igual à imagem de g, e a imagem de f (reais positivos) é igual ao domínio de g; ambas as funções são crescentes em seus domínios.

\textbf{Resolução.} \textbf{Passo 1 — Domínio e imagem de $f(t) = 2^t$.}

A função exponencial $f(t) = 2^t$ está definida para qualquer valor real de $t$, logo seu domínio é $D_f = \mathbb{R}$. Como $2^t > 0$ para todo $t$, sua imagem é $Im_f = (0, +\infty)$.

\textbf{Passo 2 — Domínio e imagem de $g(N) = \log_2(N)$.}

O logaritmo só é definido para argumentos positivos, então $D_g = (0, +\infty)$. Como o logaritmo assume qualquer valor real (positivo, negativo ou zero), sua imagem é $Im_g = \mathbb{R}$.

\textbf{Passo 3 — Relação entre domínio e imagem.}

Como $f$ e $g$ são funções inversas uma da outra, o domínio de uma corresponde exatamente à imagem da outra: $D_f = Im_g = \mathbb{R}$ e $D_g = Im_f = (0, +\infty)$.

\textbf{Passo 4 — Crescimento.}

Como a base $2 > 1$, a função exponencial $f(t) = 2^t$ é crescente em todo o seu domínio. Do mesmo modo, como a base do logaritmo é $2 > 1$, a função $g(N) = \log_2(N)$ também é crescente em todo o seu domínio.

\textbf{Conclusão.} O domínio de $f$ (os reais) coincide com a imagem de $g$, e a imagem de $f$ (reais positivos) coincide com o domínio de $g$; ambas as funções são crescentes.

\bigskip\noindent\textbf{Q021 · múltipla escolha · média · entender · teórica · conferido em parte}\par\nopagebreak
No mesmo plano cartesiano estão representadas duas curvas, $C_1$ e $C_2$, associadas a duas funções reais $f$ e $g$, cada uma definida por uma única sentença (sem sentenças por partes). Sabe-se que:

\begin{itemize}[nosep,leftmargin=*]
\item $C_1$ passa pelos pontos $\left(-1,\ \dfrac{1}{2}\right)$, $(0,1)$, $(1,2)$ e $(2,4)$;
\item $C_2$ passa pelos pontos $\left(\dfrac{1}{2},\ -1\right)$, $(1,0)$, $(2,1)$ e $(4,2)$;
\item $C_1$ e $C_2$ são simétricas uma em relação à outra em torno da reta $y=x$.
\end{itemize}

Com base na leitura desses pontos e na relação de simetria entre as curvas, assinale a alternativa que descreve corretamente o domínio, a imagem e o comportamento de crescimento de $f$ e de $g$.
\begin{enumerate}[label=(\alph*)]
\item $f$ tem domínio $(0,+\infty)$ e imagem $\mathbb{R}$; $g$ tem domínio $\mathbb{R}$ e imagem $(0,+\infty)$; ambas são crescentes.
\item $f$ é crescente, mas $g$ é decrescente, pois $C_2$ passa por pontos com ordenada negativa.
\item $f$ tem domínio $\mathbb{R}$ e imagem $(0,+\infty)$; $g$ tem domínio $\mathbb{R}$ e imagem $\mathbb{R}$; ambas são crescentes.
\item $f$ tem domínio $\mathbb{R}$ e imagem $(0,+\infty)$; $g$ tem domínio $(0,+\infty)$ e imagem $\mathbb{R}$; ambas são crescentes em todo o seu domínio.
\end{enumerate}

\textbf{Gabarito.} Alternativa D: $f$ tem domínio $\mathbb{R}$ e imagem $(0,+\infty)$; $g$ tem domínio $(0,+\infty)$ e imagem $\mathbb{R}$; ambas são crescentes em todo o seu domínio.

\textbf{Resolução.} \textbf{Passo 1 — Reconhecer o padrão de $C_1$.}

Os pontos de $C_1$ são $(-1,\tfrac12), (0,1), (1,2), (2,4)$. A cada aumento de $1$ unidade em $x$, o valor de $y$ dobra: $\tfrac12 \to 1 \to 2 \to 4$. Esse é o comportamento característico de uma função exponencial de base $2$: $f(x) = 2^x$. Verificando: $f(-1)=2^{-1}=\tfrac12$, $f(0)=1$, $f(1)=2$, $f(2)=4$. $\checkmark$

\textbf{Passo 2 — Reconhecer o padrão de $C_2$.}

Os pontos de $C_2$ são exatamente os pontos de $C_1$ com as coordenadas trocadas: $(-1,\tfrac12)\to(\tfrac12,-1)$, $(0,1)\to(1,0)$, $(1,2)\to(2,1)$, $(2,4)\to(4,2)$. Trocar as coordenadas de todos os pontos de uma curva é exatamente a operação que gera a \textbf{função inversa}. Logo $g = f^{-1}$, isto é, $g(x) = \log_2(x)$. Verificando: $g(\tfrac12)=\log_2(\tfrac12)=-1$, $g(1)=0$, $g(2)=1$, $g(4)=2$. $\checkmark$

\textbf{Passo 3 — Usar a simetria em relação a $y=x$.}

Como $g$ é a inversa de $f$, o gráfico de $g$ é o reflexo do gráfico de $f$ em torno da reta $y=x$ — o que é coerente com a informação dada no enunciado. Uma consequência direta dessa relação é que \textbf{domínio e imagem se trocam} entre a função e sua inversa:
$$\text{Dom}(f) = \text{Im}(g) \quad \text{e} \quad \text{Im}(f) = \text{Dom}(g).$$

\textbf{Passo 4 — Determinar domínio e imagem de cada uma.}

\begin{itemize}[nosep,leftmargin=*]
\item $f(x)=2^x$: está definida para todo $x$ real, e como $2^x>0$ sempre, $\text{Dom}(f)=\mathbb{R}$ e $\text{Im}(f)=(0,+\infty)$.
\item $g(x)=\log_2(x)$: só está definida para $x>0$, e assume qualquer valor real, logo $\text{Dom}(g)=(0,+\infty)$ e $\text{Im}(g)=\mathbb{R}$.
\end{itemize}

Isso confirma a relação do Passo 3: $\text{Dom}(f)=\mathbb{R}=\text{Im}(g)$ e $\text{Im}(f)=(0,+\infty)=\text{Dom}(g)$.

\textbf{Passo 5 — Analisar o crescimento.}

Como a base $2>1$ em ambas as expressões, a exponencial $f(x)=2^x$ é \textbf{crescente} em $\mathbb{R}$, e sua inversa $g(x)=\log_2(x)$ também é \textbf{crescente} em $(0,+\infty)$ (funções inversas de funções crescentes são crescentes). Note que $C_2$ conter pontos com ordenada negativa não indica decrescimento: apenas significa que, para $x$ entre $0$ e $1$, $g(x)$ é negativo — o valor da função, não seu sentido de variação.

\textbf{Conclusão:} $f$ tem domínio $\mathbb{R}$ e imagem $(0,+\infty)$; $g$ tem domínio $(0,+\infty)$ e imagem $\mathbb{R}$; ambas são crescentes.

\bigskip\noindent\textbf{Q023 · múltipla escolha · difícil · analisar · aplicada · gabarito conferido}\par\nopagebreak
Um biólogo estuda o crescimento de uma cultura de bactérias em laboratório. A quantidade de bactérias na cultura, em milhares, $t$ horas após o início da observação, é dada por $N(t) = 100 \cdot 2^{t/3}$, para $t \geq 0$. Para planejar a coleta de amostras, o biólogo também utiliza uma função que fornece o tempo necessário, em horas, para que a cultura atinja uma quantidade $N$ de bactérias (em milhares): $T(N) = 3 \cdot \log_2\left(\dfrac{N}{100}\right)$, para $N \geq 100$. Analisando o domínio, a imagem e o comportamento de crescimento das duas funções, bem como a relação entre seus gráficos, assinale a alternativa correta.
\begin{enumerate}[label=(\alph*)]
\item O domínio de $N(t)$ é $[0,+\infty)$ e sua imagem é $[100,+\infty)$; o domínio de $T(N)$ é $[100,+\infty)$ e sua imagem é $[0,+\infty)$. A função $N(t)$ é crescente, mas $T(N)$ é decrescente, pois o logaritmo cresce mais lentamente que a exponencial.
\item O domínio de $N(t)$ é $[0,+\infty)$ e sua imagem é $[100,+\infty)$; como toda função logarítmica está definida para qualquer número real, o domínio de $T(N)$ é $\mathbb{R}$, e sua imagem é $[0,+\infty)$. Ambas as funções são crescentes.
\item O domínio de $N(t)$ é $[0,+\infty)$ e sua imagem é $[100,+\infty)$; o domínio de $T(N)$ é $[100,+\infty)$ e sua imagem é $[0,+\infty)$. Ambas as funções são crescentes, e seus gráficos são simétricos em relação à reta $y=x$, pois $T$ é a função inversa de $N$.
\item O domínio de $N(t)$ é $[100,+\infty)$ e sua imagem é $[0,+\infty)$; o domínio de $T(N)$ é $[0,+\infty)$ e sua imagem é $[100,+\infty)$. Ambas as funções são crescentes, e seus gráficos são simétricos em relação à reta $y=x$.
\end{enumerate}

\textbf{Gabarito.} Domínio de $N(t)$: $[0,+\infty)$; Imagem de $N(t)$: $[100,+\infty)$; Domínio de $T(N)$: $[100,+\infty)$; Imagem de $T(N)$: $[0,+\infty)$; ambas crescentes; gráficos simétricos em relação a $y=x$ (funções inversas).

\textbf{Resolução.} \textbf{Passo 1 — Domínio de $N(t)$:} o contexto exige $t \geq 0$ (tempo não pode ser negativo), logo $D(N) = [0,+\infty)$.

\textbf{Passo 2 — Imagem de $N(t)$:} em $t=0$, $N(0) = 100 \cdot 2^0 = 100$; como a base $2>1$, $N(t)$ cresce sem limite quando $t \to +\infty$. Logo $Im(N) = [100,+\infty)$.

\textbf{Passo 3 — Crescimento de $N(t)$:} como a base da exponencial é $2>1$, $N(t)$ é estritamente crescente em todo o seu domínio.

\textbf{Passo 4 — Domínio de $T(N)$:} o problema restringe $N \geq 100$ (tempo para atingir pelo menos a população inicial), o que também garante que o argumento do logaritmo, $N/100$, seja $\geq 1 > 0$. Logo $D(T) = [100,+\infty)$.

\textbf{Passo 5 — Imagem de $T(N)$:} em $N=100$, $T(100) = 3\log_2(1) = 0$; quando $N \to +\infty$, $T(N) \to +\infty$. Logo $Im(T) = [0,+\infty)$.

\textbf{Passo 6 — Crescimento de $T(N)$:} como $\log_2$ é crescente (base $2>1$) e o fator $3$ é positivo, $T(N)$ é estritamente crescente.

\textbf{Passo 7 — Relação entre as funções:} substituindo, $T(N(t)) = 3\log_2\!\big(2^{t/3}\big) = 3\cdot\dfrac{t}{3} = t$, e analogamente $N(T(N))=N$. Logo $T$ é a função inversa de $N$. Por isso, $D(N)=Im(T)=[0,+\infty)$ e $Im(N)=D(T)=[100,+\infty)$, e os gráficos de $N$ e $T$ são simétricos em relação à reta $y=x$.

\textbf{Conclusão:} a alternativa correta é a que descreve exatamente essas correspondências de domínio, imagem, crescimento e simetria.
\subsection*{EM13MAT404}
{\small\itshape Identificar as características fundamentais das funções seno e cosseno (periodicidade, domínio, imagem), por meio da comparação das representações em ciclos trigonométricos e em planos cartesianos, com ou sem apoio de tecnologias digitais.}

\bigskip\noindent\textbf{Q026 · discursiva · fácil · entender · aplicada · gabarito conferido}\par\nopagebreak
Uma roda-gigante tem raio de 10 m e seu eixo central está a 12 m de altura em relação ao solo. Ela gira em velocidade constante, completando uma volta a cada 40 segundos. No instante $t = 0$ (em segundos), um determinado carrinho está exatamente no ponto mais baixo do movimento circular. A altura $h(t)$, em metros, desse carrinho em relação ao solo, $t$ segundos após o início da contagem do tempo, é dada por:

$$h(t) = 12 - 10\cos\left(\dfrac{\pi t}{20}\right)$$

Imagine o movimento do carrinho representado simultaneamente de duas formas: como um ponto percorrendo o ciclo trigonométrico (com ângulo variando de $0$ a $2\pi$ radianos e voltando a se repetir) e como uma curva no plano cartesiano, com o tempo $t$ no eixo horizontal e a altura $h(t)$ no eixo vertical.

Com base nessa comparação entre as duas representações, responda:

a) Qual é o período da função $h$, ou seja, o tempo necessário para que o carrinho complete uma volta e a altura volte a se repetir? Explique como esse valor se relaciona com o fato de o ciclo trigonométrico ter $2\pi$ radianos.

b) Qual é o domínio da função $h$, considerando que $t$ representa o tempo decorrido desde o início do funcionamento da roda-gigante (que continua girando indefinidamente)?

c) Qual é a imagem da função $h$, ou seja, quais são as menores e maiores alturas atingidas pelo carrinho? Relacione essas alturas com os pontos do ciclo trigonométrico em que $\cos(\theta) = 1$ e $\cos(\theta) = -1$.

\textbf{Gabarito.} a) Período $T = 40$ s, pois o ângulo $\frac{\pi t}{20}$ percorre $2\pi$ radianos (uma volta completa no ciclo trigonométrico) exatamente quando $t=40$. b) Domínio: $[0,+\infty)$. c) Imagem: $[2,22]$ metros, correspondendo a $\cos(\theta)=1$ (altura mínima, 2 m) e $\cos(\theta)=-1$ (altura máxima, 22 m).

\textbf{Resolução.} \textbf{a) Período}

No ciclo trigonométrico, uma volta completa corresponde a um ângulo de $2\pi$ radianos. Na função dada, o ângulo (argumento do cosseno) é $\theta = \dfrac{\pi t}{20}$.

O período $T$ é o valor de $t$ para o qual esse ângulo completa exatamente $2\pi$:

$$\frac{\pi T}{20} = 2\pi \implies T = \frac{2\pi \cdot 20}{\pi} = 40$$

Ou seja, cada volta completa no ciclo trigonométrico ($2\pi$ radianos) corresponde, no plano cartesiano, a um intervalo de $40$ segundos no eixo $t$ após o qual o gráfico de $h(t)$ se repete. Assim, o \textbf{período é $T = 40$ segundos}.

\textbf{b) Domínio}

Como $t$ representa o tempo decorrido desde o início do funcionamento, e a roda-gigante continua girando indefinidamente (repetindo o ciclo trigonométrico voltas após voltas), o tempo pode assumir qualquer valor real não negativo. Logo:

$$D = \{t \in \mathbb{R} \mid t \ge 0\} = [0, +\infty)$$

\textbf{c) Imagem}

No ciclo trigonométrico, o cosseno de um ângulo varia entre $-1$ e $1$:
\begin{itemize}[nosep,leftmargin=*]
\item $\cos(\theta) = 1$ ocorre no ponto $(1,0)$ do ciclo (ângulo $0$, $2\pi$, ...), correspondente ao ponto mais próximo do eixo de rotação na direção considerada;
\item $\cos(\theta) = -1$ ocorre no ponto $(-1,0)$ do ciclo (ângulo $\pi$), o ponto diametralmente oposto.
\end{itemize}

Substituindo esses extremos em $h(t) = 12 - 10\cos(\theta)$:

\begin{itemize}[nosep,leftmargin=*]
\item Quando $\cos(\theta) = 1$: $h = 12 - 10(1) = 2$ (altura mínima, carrinho na base);
\item Quando $\cos(\theta) = -1$: $h = 12 - 10(-1) = 22$ (altura máxima, carrinho no topo).
\end{itemize}

Como o cosseno assume continuamente todos os valores entre $-1$ e $1$ ao longo de cada volta, $h(t)$ assume continuamente todos os valores entre $2$ e $22$. Logo, a \textbf{imagem é $[2, 22]$} metros.

\bigskip\noindent\textbf{Q029 · discursiva · média · entender · teórica · gabarito conferido}\par\nopagebreak
Considere a função $f(x) = 3\cos(2x) - 1$, definida para $x \in \mathbb{R}$.

No ciclo trigonométrico (circunferência de raio 1 centrada na origem), o valor de $\cos(\theta)$ corresponde à abscissa do ponto $P$ que se desloca sobre a circunferência conforme o ângulo $\theta$ varia.

Usando essa relação entre o ciclo trigonométrico e a representação de $f$ no plano cartesiano, determine, justificando cada resposta com base no comportamento do ponto $P$ sobre a circunferência:

a) o domínio de $f$;

b) o período de $f$;

c) a imagem de $f$.

\textbf{Gabarito.} Domínio: $\mathbb{R}$; Período: $\pi$; Imagem: $[-4, 2]$

\textbf{Resolução.} \textbf{a) Domínio}

No ciclo trigonométrico, o ângulo $\theta$ pode assumir qualquer valor real — positivo, negativo, maior que $2\pi$ etc. — pois o ponto $P$ simplesmente continua girando pela circunferência indefinidamente. Como $2x$ pode ser qualquer número real quando $x$ percorre $\mathbb{R}$, o cosseno de $2x$ está sempre definido. Logo:
$$D(f) = \mathbb{R}$$

\textbf{b) Período}

No ciclo trigonométrico, o ponto $P$ retorna à mesma posição sempre que o ângulo aumenta $2\pi$ radianos (uma volta completa). Aqui o ângulo que gira é $2x$, então uma volta completa ocorre quando:
$$2x = 2\pi \;\Rightarrow\; x = \pi$$
Assim, $\cos(2x)$ repete seus valores a cada $x = \pi$. As operações de multiplicar por $3$ e subtrair $1$ apenas reescalam e deslocam verticalmente o gráfico no plano cartesiano — elas não alteram o ritmo com que o ponto $P$ completa voltas no ciclo. Portanto o período de $f$ é:
$$T = \pi$$

\textbf{c) Imagem}

Como o ciclo trigonométrico tem raio $1$, a abscissa do ponto $P$ (isto é, $\cos(2x)$) varia sempre entre $-1$ e $1$:
$$-1 \le \cos(2x) \le 1$$
Multiplicando por $3$ (o que estica a projeção no plano cartesiano, correspondendo a um círculo 'equivalente' de raio 3 na análise da amplitude):
$$-3 \le 3\cos(2x) \le 3$$
Subtraindo $1$ (deslocamento vertical do gráfico no plano cartesiano):
$$-4 \le 3\cos(2x) - 1 \le 2$$
Logo:
$$Im(f) = [-4, 2]$$

\bigskip\noindent\textbf{Q032 · discursiva · difícil · entender · aplicada · gabarito conferido}\par\nopagebreak
Uma roda-gigante de um parque de diversões tem raio de 15 m, e seu eixo central está fixado a 17 m de altura em relação ao solo. A roda gira com velocidade angular constante, completando uma volta a cada 40 segundos. No instante t = 0 (em segundos, contado a partir da abertura do parque), uma cabine específica está exatamente no ponto mais baixo de sua trajetória circular. O parque mantém essa roda-gigante em funcionamento ininterrupto durante as 12 horas de seu horário de atendimento diário (das 9h às 21h).

Um estudante decide modelar a altura h(t), em metros, dessa cabine em relação ao solo, t segundos após a abertura do parque, representando simultaneamente o movimento da cabine no ciclo trigonométrico (posição angular ao longo da circunferência) e no plano cartesiano (gráfico de h em função de t).

a) Determine a lei da função h(t), explicitando como o ângulo percorrido no ciclo trigonométrico se relaciona com a altura registrada no gráfico cartesiano.

b) Indique o período dessa função e explique seu significado físico no contexto da roda-gigante.

c) Considerando que a roda-gigante funciona apenas durante o horário diário de operação do parque, determine o domínio de h(t) para um dia de funcionamento.

d) Determine a imagem de h(t), justificando o resultado a partir da comparação entre a variação do cosseno no ciclo trigonométrico (entre -1 e 1) e a amplitude observada no gráfico cartesiano.

\textbf{Gabarito.} h(t) = 17 - 15·cos($\pi$t/20) (equivalente a 17 + 15·sen($\pi$t/20 - $\pi$/2)); período = 40 s; domínio, para um dia de operação = [0, 43200] segundos; imagem = [2, 32] metros.

\textbf{Resolução.} \#\# Passo 1: Relacionar o movimento circular ao ciclo trigonométrico

A cabine descreve uma circunferência de raio $R=15$ m ao redor do centro da roda, que está a $17$ m do solo. Associamos a cada instante $t$ um ângulo $\theta(t)$ varrido a partir de uma posição de referência no ciclo trigonométrico.

Como a roda dá uma volta completa ($2\pi$ rad) a cada $40$ s, a velocidade angular é $\omega = \dfrac{2\pi}{40} = \dfrac{\pi}{20}$ rad/s, logo $\theta(t) = \dfrac{\pi}{20}t + \theta_0$.

\#\# Passo 2: Usar a condição inicial para determinar $\theta_0$

No ciclo trigonométrico, a ordenada de um ponto (multiplicada pelo raio) representa sua altura em relação ao centro da circunferência. Em $t=0$ a cabine está no ponto mais baixo, ou seja, na posição angular $\theta = -\dfrac{\pi}{2}$ (onde $\text{sen}\,\theta=-1$). Logo $\theta_0=-\dfrac{\pi}{2}$ e $\theta(t)=\dfrac{\pi}{20}t-\dfrac{\pi}{2}$.

\#\# Passo 3: Escrever h(t) transportando o ciclo para o plano cartesiano

A altura em relação ao solo é a altura do centro somada à projeção vertical multiplicada pelo raio:
$$h(t)=17+15\,\text{sen}\left(\dfrac{\pi}{20}t-\dfrac{\pi}{2}\right)$$
Usando $\text{sen}(x-\pi/2)=-\cos x$, obtemos a forma equivalente, mais simples de analisar no plano cartesiano:
$$h(t)=17-15\cos\left(\dfrac{\pi}{20}t\right)$$
Assim, cada volta completa no ciclo trigonométrico (variação de $2\pi$ no ângulo) corresponde, no plano cartesiano, a um ciclo completo do gráfico de $h(t)$.

\#\# Passo 4: Período

O período é o tempo necessário para o ângulo variar $2\pi$:
$$T=\dfrac{2\pi}{\pi/20}=40\text{ s}$$
Fisicamente, é o tempo de uma volta completa da roda-gigante, coincidindo com o dado do enunciado e confirmando o modelo.

\#\# Passo 5: Domínio

A roda funciona continuamente por $12$ horas $=43200$ s por dia, com $t$ contado a partir da abertura ($t=0$). Portanto, para um dia de funcionamento:
$$D=[0,\,43200]\ \text{(em segundos)}$$

\#\# Passo 6: Imagem

No ciclo trigonométrico, $\cos\theta$ varia continuamente entre $-1$ e $1$ conforme o ângulo percorre toda a circunferência. Como $43200/40=1080$ voltas completas cabem no intervalo de funcionamento, o ciclo é percorrido integralmente muitas vezes, garantindo que os valores extremos são de fato atingidos. Transportando essa variação para o plano cartesiano:
\begin{itemize}[nosep,leftmargin=*]
\item mínimo: $h_{min}=17-15(1)=2$ m (quando $\cos\theta=1$)
\item máximo: $h_{max}=17-15(-1)=32$ m (quando $\cos\theta=-1$)
\end{itemize}

Logo, a imagem é:
$$Im(h)=[2,\,32]\text{ (em metros)}$$

\bigskip\noindent\textbf{Q035 · múltipla escolha · fácil · entender · teórica · gabarito conferido}\par\nopagebreak
Ao girar um ponto sobre o ciclo trigonométrico (circunferência de raio 1 centrada na origem), sua ordenada varia continuamente entre $-1$ e $1$, repetindo os mesmos valores a cada volta completa de $2\pi$ radianos. Essa ordenada é justamente o valor de $\sin(x)$ representado no plano cartesiano. Considere agora a função $f(x) = 3\sin(x) - 1$, definida para todo número real $x$. Com base na relação entre o que ocorre no ciclo trigonométrico e o comportamento do gráfico de $f$ no plano cartesiano, assinale a alternativa que indica corretamente o domínio, o período e a imagem de $f$.
\begin{enumerate}[label=(\alph*)]
\item Domínio $\mathbb{R}$, período $6\pi$, imagem $[-4, 2]$
\item Domínio $\mathbb{R}$, período $2\pi$, imagem $[-4, 2]$
\item Domínio $\mathbb{R}$, período $2\pi$, imagem $[-2, 4]$
\item Domínio $\mathbb{R}$, período $2\pi$, imagem $[-1, 1]$
\end{enumerate}

\textbf{Gabarito.} Domínio $\mathbb{R}$, período $2\pi$, imagem $[-4, 2]$

\textbf{Resolução.} \textbf{Passo 1 — Domínio.} No ciclo trigonométrico, o ponto pode girar indefinidamente, para qualquer ângulo $x$ (positivo, negativo ou nulo), sem restrição. Logo, $\sin(x)$ está definido para todo $x$ real, e como $f(x)=3\sin(x)-1$ é obtida por operações que não restringem esse domínio, $D(f) = \mathbb{R}$.

\textbf{Passo 2 — Período.} No ciclo trigonométrico, a ordenada do ponto repete seus valores a cada volta completa, ou seja, a cada $2\pi$ radianos, pois o argumento de $\sin$ em $f(x)=3\sin(x)-1$ é simplesmente $x$ (não há fator multiplicando $x$ dentro do seno). Multiplicar o resultado do seno por $3$ e subtrair $1$ apenas estica e desloca verticalmente o gráfico no plano cartesiano — não altera a rapidez com que o ponto completa uma volta no ciclo. Portanto, o período continua sendo $T = 2\pi$.

\textbf{Passo 3 — Imagem.} No ciclo trigonométrico, a ordenada satisfaz $-1 \le \sin(x) \le 1$. Multiplicando por $3$: $-3 \le 3\sin(x) \le 3$. Subtraindo $1$ de cada parte da desigualdade (deslocamento vertical do gráfico no plano cartesiano): $-3-1 \le 3\sin(x)-1 \le 3-1$, ou seja, $-4 \le f(x) \le 2$. Logo, a imagem é $[-4, 2]$.

\textbf{Conclusão:} $D(f)=\mathbb{R}$, $T=2\pi$, $Im(f)=[-4,2]$.

\bigskip\noindent\textbf{Q038 · múltipla escolha · média · analisar · aplicada · gabarito conferido}\par\nopagebreak
Uma roda-gigante tem 15 m de raio e seu centro está fixado a 17 m do solo. Ela gira com velocidade angular constante, completando uma volta inteira a cada 40 segundos. Uma passageira embarca no ponto mais baixo da roda (o ponto da circunferência mais próximo do chão) no instante $t=0$, e a roda gira continuamente a partir daí. Associando o movimento circular da cabine ao ciclo trigonométrico — em que o ângulo percorrido a partir do ponto mais baixo corresponde a $\theta(t)=\dfrac{2\pi}{40}t$ — a altura $h(t)$, em metros, da cabine em relação ao solo, $t$ segundos após o embarque, é dada por
$$h(t) = 17 - 15\cos\left(\frac{2\pi}{40}t\right).$$
Considerando a correspondência entre o percurso da cabine no ciclo trigonométrico e o gráfico de $h(t)$ no plano cartesiano, identifique a alternativa que apresenta corretamente o período do movimento, o domínio de $h$ (levando em conta que $t$ representa o tempo decorrido desde o embarque, com a roda girando indefinidamente) e a imagem de $h$ (a faixa de alturas atingidas pela cabine).
\begin{enumerate}[label=(\alph*)]
\item Período de $\dfrac{\pi}{20}$ s; domínio $t \in [0, +\infty)$; imagem $[2, 32]$ metros.
\item Período de 40 s; domínio $t \in [0, +\infty)$; imagem $[-15, 15]$ metros.
\item Período de 40 s; domínio $t \in [0, 40]$; imagem $[2, 32]$ metros.
\item Período de 40 s; domínio $t \in [0, +\infty)$; imagem $[2, 32]$ metros.
\end{enumerate}

\textbf{Gabarito.} D

\textbf{Resolução.} \textbf{1. Relacionando o ciclo trigonométrico ao movimento}

No ciclo trigonométrico, o cosseno de um ângulo $\theta$ varia entre $-1$ e $1$ conforme o ponto percorre a circunferência de raio $1$. Aqui, a cabine percorre uma circunferência de raio $15$ m centrada a $17$ m do solo, e o ângulo percorrido no instante $t$ é $\theta(t) = \dfrac{2\pi}{40}t$. Como a passageira embarca no ponto mais baixo (onde $\cos\theta = 1$ dá a altura mínima), a altura é modelada por $h(t) = 17 - 15\cos\theta(t)$: o sinal negativo garante que, em $\theta=0$, tenhamos $h(0)=17-15=2$ m (ponto mais baixo).

\textbf{2. Período}

No ciclo trigonométrico, uma volta completa corresponde a $\theta$ variar de $0$ a $2\pi$. Isso ocorre quando
$$\frac{2\pi}{40}t = 2\pi \;\Rightarrow\; t = 40.$$
Ou seja, transportando o período fundamental $2\pi$ do cosseno (no ciclo) para o eixo $t$ (no plano cartesiano) através do fator de escala $\frac{2\pi}{40}$, obtém-se o período $T = 40$ segundos — o tempo de uma volta completa da roda.

\textbf{3. Domínio}

Como $t$ representa o tempo decorrido a partir do embarque e a roda continua girando indefinidamente (não para após uma volta), os valores possíveis de $t$ são todos os números reais não negativos:
$$\text{Dom}(h) = [0, +\infty).$$

\textbf{4. Imagem}

No ciclo trigonométrico, $\cos\theta \in [-1,1]$. Logo:
$$-1 \le \cos\theta(t) \le 1 \;\Rightarrow\; -15 \le -15\cos\theta(t) \le 15 \;\Rightarrow\; 2 \le 17-15\cos\theta(t) \le 32.$$
Portanto a imagem de $h$ é o intervalo $[2, 32]$ metros — a altura mínima ($2$ m, ponto mais baixo) e a máxima ($32$ m, ponto mais alto), obtidas somando/subtraindo o raio (amplitude) à altura do centro do ciclo.

\textbf{Conclusão:} Período $=40$ s; Domínio $=[0,+\infty)$; Imagem $=[2,32]$ m — alternativa D.

\bigskip\noindent\textbf{Q041 · múltipla escolha · difícil · entender · teórica · gabarito conferido}\par\nopagebreak
Um ponto P parte da posição (0, 5) no instante t = 0 e passa a se mover indefinidamente, no sentido anti-horário, sobre uma circunferência de centro (0, 3) e raio 2, completando uma volta completa a cada intervalo de tempo $\Delta$t = 4$\pi$ unidades de tempo. A ordenada de P, em função do tempo t, é dada por y(t) = 3 + 2cos(t/2).

Comparando a trajetória de P na circunferência (o ciclo trigonométrico deslocado e ampliado) com a representação gráfica de y(t) no plano cartesiano t × y, assinale a alternativa que apresenta corretamente o domínio, a imagem e o período dessa função.
\begin{enumerate}[label=(\alph*)]
\item Domínio = $\mathbb{R}$; Imagem = [1, 5]; Período = 4$\pi$
\item Domínio = [0, +$\infty$); Imagem = [1, 5]; Período = 2$\pi$
\item Domínio = [0, +$\infty$); Imagem = [1, 5]; Período = 4$\pi$
\item Domínio = [0, +$\infty$); Imagem = [-2, 2]; Período = 4$\pi$
\end{enumerate}

\textbf{Gabarito.} C) Domínio = [0, +$\infty$); Imagem = [1, 5]; Período = 4$\pi$

\textbf{Resolução.} \textbf{1) Relacionando o ciclo com a função}

Se P parte de $(0,5)$ e gira com ângulo $\theta(t) = \dfrac{\pi}{2} + \dfrac{t}{2}$, sua ordenada é $y(t) = 3 + 2\sin\theta(t) = 3 + 2\cos\left(\dfrac{t}{2}\right)$, coerente com o enunciado.

\textbf{2) Imagem — usando a circunferência}

Na circunferência, o centro tem ordenada $3$ e o raio é $2$; logo a ordenada de qualquer ponto da trajetória varia entre $3-2=1$ e $3+2=5$. Essa é exatamente a faixa vertical ocupada pela curva no plano cartesiano, pois o cosseno varia continuamente entre $-1$ e $1$:
$$y(t) = 3 + 2\cos(t/2) \in [3-2,\ 3+2] = [1,5].$$
Assim, $\text{Im}(y) = [1,5]$ (não $[-1,1]$ nem $[-2,2]$, que ignoram o deslocamento vertical do centro).

\textbf{3) Período — usando a volta completa no ciclo}

Uma volta completa no ciclo corresponde a um acréscimo de $2\pi$ no ângulo $\theta$. Como $\theta(t) = \pi/2 + t/2$, um acréscimo $\Delta t$ em $t$ gera um acréscimo $\Delta t/2$ em $\theta$. Igualando a uma volta completa:
$$\frac{\Delta t}{2} = 2\pi \quad\Rightarrow\quad \Delta t = 4\pi,$$
que é exatamente o valor dado no enunciado. Logo o período fundamental é $T = 4\pi$ (não $2\pi$, que seria o período de $\cos(t)$ sem levar em conta a 'lentidão' do giro, causada pelo fator $1/2$ no argumento).

\textbf{4) Domínio — comparando o ciclo (que se repete para sempre) com o gráfico cartesiano (que começa em t = 0)}

No ciclo trigonométrico puro, o ângulo $\theta$ poderia crescer ou decrescer indefinidamente, sugerindo (erradamente) que a função estaria definida para todo $t \in \mathbb{R}$. Mas o movimento de P \textbf{começa} em $t=0$: antes desse instante o ponto simplesmente não existia em movimento, e depois dele continua girando para sempre. Ou seja, o eixo do tempo no gráfico cartesiano não é bidirecional como o ângulo abstrato do ciclo — ele tem uma origem física. Portanto o domínio da função que descreve esse movimento real é:
$$D(y) = [0, +\infty).$$

\textbf{5) Conclusão}

$$D(y) = [0,+\infty), \quad \text{Im}(y) = [1,5], \quad T = 4\pi.$$
\subsection*{EM13MAT405}
{\small\itshape Reconhecer funções definidas por uma ou mais sentenças (como a tabela do Imposto de Renda, contas de luz, água, gás etc.), em suas representações algébrica e gráfica, convertendo essas representações de uma para outra e identificando domínios de validade, imagem, crescimento e decrescimento.}

\bigskip\noindent\textbf{Q002 · discursiva · fácil · entender · teórica · conferido em parte}\par\nopagebreak
Uma companhia de saneamento cobra a conta de água residencial mensal com base no consumo $x$, medido em metros cúbicos ($m^3$), segundo a lei abaixo, em que $C(x)$ é o valor da conta, em reais. Considere que $x$ pode assumir qualquer valor real não negativo (não necessariamente inteiro).

$$C(x) = \begin{cases} 15 + 2x, & \text{se } 0 \le x \le 10 \\ 35 + 4(x-10), & \text{se } x > 10 \end{cases}$$

a) Determine o domínio de validade dessa função, considerando o significado de $x$ no contexto.

b) Calcule o valor da conta para um consumo de exatamente $10\,m^3$ e para um consumo de $12\,m^3$, indicando em cada caso qual das duas sentenças da lei foi utilizada.

c) Determine a imagem da função $C$.

d) A função $C$ é crescente, decrescente ou nenhuma das duas coisas em todo o seu domínio? Justifique observando o comportamento de cada sentença.

e) Descreva (sem desenhar) o formato do gráfico de $C$: de que tipo são os dois trechos, se há salto (descontinuidade) no ponto de junção $x=10$, e como se compara a inclinação de um trecho com a do outro.

\textbf{Gabarito.} a) Domínio: $[0,+\infty)$. b) $C(10)=35$ (1ª sentença); $C(12)=43$ (2ª sentença). c) Imagem: $[15,+\infty)$. d) $C$ é crescente em todo o domínio. e) Duas semirretas que se encontram sem salto em $(10,35)$, sendo a segunda mais inclinada (coeficiente angular 4) que a primeira (coeficiente angular 2).

\textbf{Resolução.} \textbf{a) Domínio de validade}

Como $x$ representa um consumo de água, não pode ser negativo, e não há um limite superior estabelecido no problema (o consumo pode crescer indefinidamente). Logo, o domínio é

$$D = [0, +\infty) = \text{Interval}(0,\infty).$$

\textbf{b) Valores em $x=10$ e $x=12$ (mobilizando a mudança de sentença)}

Como $10$ satisfaz $0 \le x \le 10$, usamos a \textbf{primeira} sentença:
$$C(10) = 15 + 2\cdot 10 = 15 + 20 = 35.$$

Como $12 > 10$, usamos a \textbf{segunda} sentença:
$$C(12) = 35 + 4(12-10) = 35 + 4\cdot 2 = 35 + 8 = 43.$$

Observe que a mudança de regra em $x=10$ não gera salto: as duas sentenças coincidem exatamente nesse ponto ($C(10)=35$ pelas duas expressões), garantindo continuidade.

\textbf{c) Imagem}

No trecho $0 \le x \le 10$: $C(x) = 15+2x$ é linear crescente, variando de $C(0)=15$ até $C(10)=35$. Logo, nesse trecho, $C$ assume todos os valores em $[15,35]$.

No trecho $x>10$: $C(x) = 35+4(x-10)$ também é linear crescente, partindo de valores próximos (mas maiores) que $35$ e crescendo sem limite quando $x \to \infty$. Logo, nesse trecho, $C$ assume todos os valores em $(35, +\infty)$.

Unindo os dois conjuntos de valores:
$$\text{Im}(C) = [15,35] \cup (35,\infty) = [15,+\infty) = \text{Interval}(15,\infty).$$

\textbf{d) Crescimento}

No primeiro trecho, o coeficiente angular é $2>0$: função crescente.
No segundo trecho, o coeficiente angular é $4>0$: função crescente.
Como os dois trechos são crescentes e se encaixam continuamente em $x=10$ (sem salto para baixo), $C$ é \textbf{crescente em todo o domínio} $[0,+\infty)$.

\textbf{e) Formato do gráfico}

O gráfico é formado por dois segmentos de reta:
\begin{itemize}[nosep,leftmargin=*]
\item de $(0,15)$ até $(10,35)$, com inclinação $2$;
\item a partir de $(10,35)$, seguindo indefinidamente, com inclinação $4$.
\end{itemize}

Como os dois valores coincidem em $x=10$, não há salto (descontinuidade) no gráfico: as duas retas se encontram exatamente nesse ponto, formando um "bico" onde a reta muda de inclinação, ficando \textbf{mais inclinada} (mais íngreme) para $x>10$ do que para $0 \le x \le 10$.

\bigskip\noindent\textbf{Q004 · discursiva · média · entender · aplicada · conferido em parte}\par\nopagebreak
Uma companhia de saneamento calcula o valor mensal da conta de água de acordo com o consumo $x$ (em m³) do cliente, segundo a tabela abaixo:

\begin{center}
\begin{tabular}{|p{0.430\linewidth}|p{0.430\linewidth}|}
\hline
\textbf{Consumo mensal} & \textbf{Regra de cobrança} \\
\hline
$0 \le x \le 10$ & taxa mínima fixa de R\$\,20,00 \\
\hline
$10 < x \le 20$ & R\$\,20,00 mais R\$\,3,00 por cada m³ que exceder 10 \\
\hline
$x > 20$ & R\$\,50,00 mais R\$\,5,00 por cada m³ que exceder 20 \\
\hline
\end{tabular}
\end{center}

Seja $f(x)$ o valor, em reais, da conta de água correspondente a um consumo de $x$ m³.

a) Escreva a lei de formação algébrica de $f(x)$, especificando a expressão válida em cada uma das três faixas de consumo.

b) Qual é o domínio de validade de $f$, isto é, o conjunto de valores de $x$ para os quais a regra de cobrança faz sentido?

c) Uma família definiu que pode gastar, no máximo, R\$\,41,00 por mês com água, e sabe que seu consumo está entre 10 m³ e 20 m³. Qual é o maior consumo mensal, em m³, que essa família pode ter sem ultrapassar esse orçamento?

d) Esboce o gráfico de $f$ para $0 \le x \le 30$ e classifique o comportamento de $f$ (crescente, decrescente ou constante) em cada um dos três trechos.

e) Determine a imagem de $f$ restrita ao intervalo $0 \le x \le 30$.

\textbf{Gabarito.} a) $f(x)=20$ se $0\le x\le 10$; $f(x)=3x-10$ se $10<x\le 20$; $f(x)=5x-50$ se $x>20$. b) $D(f)=[0,+\infty)$. c) 17 m³. d) constante em $[0,10]$; crescente em $(10,20]$ e em $(20,30]$ (com inclinação maior nesse último trecho). e) Imagem $=[20,100]$.

\textbf{Resolução.} \textbf{a) Lei de formação}

Diretamente da tabela, escrevendo as expressões em função de $x$:

$$f(x)=\begin{cases}20, & 0\le x\le 10\\[2pt] 20+3(x-10), & 10< x\le 20\\[2pt] 50+5(x-20), & x> 20\end{cases}$$

Simplificando as duas últimas sentenças: $20+3(x-10)=3x-10$ e $50+5(x-20)=5x-50$.

$$f(x)=\begin{cases}20, & 0\le x\le 10\\ 3x-10, & 10< x\le 20\\ 5x-50, & x> 20\end{cases}$$

Verificação de continuidade: em $x=10$, $3(10)-10=20$ (coincide com a 1ª sentença); em $x=20$, $3(20)-10=50$ e $5(20)-50=50$ (coincidem). A função é contínua.

\textbf{b) Domínio de validade}

Como $x$ representa um consumo de água, não pode ser negativo, e não há limite superior estabelecido pela tarifação (a terceira faixa vale para qualquer consumo acima de 20 m³). Logo:

$$D(f) = \{x \in \mathbb{R} \mid x \ge 0\} = [0, +\infty)$$

\textbf{c) Consumo máximo dentro do orçamento}

Como o consumo da família está entre 10 e 20 m³, usa-se a segunda sentença:

$$20+3(x-10)=41$$
$$3(x-10)=21 \;\Rightarrow\; x-10=7 \;\Rightarrow\; x=17$$

Como $17 \in (10,20]$, a solução é válida nessa faixa. O consumo máximo é \textbf{17 m³}.

\textbf{d) Gráfico e comportamento}

\begin{itemize}[nosep,leftmargin=*]
\item Para $0\le x\le 10$: o gráfico é um segmento horizontal na altura $y=20$ — $f$ é \textbf{constante}.
\item Para $10< x\le 20$: o gráfico é um segmento de reta com coeficiente angular 3, subindo de $(10,20)$ até $(20,50)$ — $f$ é \textbf{crescente}.
\item Para $20< x\le 30$: o gráfico é um segmento de reta com coeficiente angular 5, subindo de $(20,50)$ até $(30,100)$ — $f$ é \textbf{crescente}, e mais inclinado que no trecho anterior.
\end{itemize}

O gráfico completo é uma linha poligonal contínua: um patamar horizontal seguido de dois segmentos crescentes com inclinações diferentes.

\textbf{e) Imagem no intervalo $[0,30]$}

No trecho constante, $f$ assume apenas o valor $20$. A partir de $x=10$, $f$ cresce continuamente (sem saltos, pois a função é contínua) até atingir, em $x=30$:

$$f(30)=5(30)-50=100$$

Como $f$ é contínua e não decresce em nenhum ponto de $[0,30]$, todos os valores entre $20$ e $100$ são atingidos, e nenhum valor fora desse intervalo é atingido. Logo:

$$\text{Im}(f|_{[0,30]}) = [20, 100]$$

\bigskip\noindent\textbf{Q006 · discursiva · difícil · entender · teórica · conferido em parte}\par\nopagebreak
Uma companhia de saneamento cobra a conta de água residencial de acordo com faixas de consumo mensal $m$, medido em metros cúbicos (m³), sendo o serviço válido para consumos de até $50\,m^3$. As regras de cobrança, em reais, são as seguintes:

\begin{center}
\begin{tabular}{|p{0.430\linewidth}|p{0.430\linewidth}|}
\hline
\textbf{Faixa de consumo} & \textbf{Regra de cobrança} \\
\hline
$0 \le m \le 10$ & Taxa fixa de manutenção de R\$\,5,00 mais R\$\,1,50 por m³ consumido \\
\hline
$10 < m \le 30$ & R\$\,20,00 fixos (referentes aos primeiros 10 m³) mais R\$\,2,00 por m³ que exceder 10 m³ \\
\hline
$30 < m \le 50$ & R\$\,60,00 fixos (referentes aos primeiros 30 m³) mais R\$\,4,00 por m³ que exceder 30 m³ \\
\hline
\end{tabular}
\end{center}

Seja $C(m)$ o valor da conta, em reais, para um consumo de $m$ metros cúbicos.

a) Escreva a lei algébrica de $C(m)$, definida por partes, especificando o domínio de validade de cada sentença.

b) Determine a imagem da função $C$ considerando todo o intervalo de consumo válido, de $0$ a $50\,m^3$.

c) Em certo mês, uma família pagou uma conta de R\$\,74,00. Determine o consumo desse mês, indicando claramente por que sentença da função esse valor deve ser calculado.

d) A conta de água é uma função crescente, decrescente ou constante no intervalo de consumo considerado? Justifique sua resposta levando em conta as três sentenças da função.

\textbf{Gabarito.} a) $C(m) = 1{,}5m+5$ para $0\le m\le10$; $C(m)=20+2(m-10)$ para $10<m\le30$; $C(m)=60+4(m-30)$ para $30<m\le50$. b) Imagem $=[5,140]$. c) $m=33{,}5\,m^3$ (calculado na terceira sentença, pois $74>60$). d) A função é crescente em todo o domínio $[0,50]$.

\textbf{Resolução.} \textbf{a) Lei algébrica por partes}

Traduzindo cada faixa da tabela:

$$C(m) = \begin{cases} 1{,}5m + 5, & 0 \le m \le 10 \\ 20 + 2(m-10), & 10 < m \le 30 \\ 60 + 4(m-30), & 30 < m \le 50 \end{cases}$$

Verificando a coerência nas transições: em $m=10$, a primeira sentença dá $1{,}5\times10+5=20$, que coincide com o valor inicial da segunda sentença. Em $m=30$, a segunda sentença dá $20+2\times20=60$, que coincide com o valor inicial da terceira. A função é, portanto, contínua.

\textbf{b) Imagem da função}

Cada sentença tem coeficiente angular positivo ($1{,}5$; $2$; $4$), logo $C$ é crescente em cada um dos três intervalos. Como os valores se encaixam continuamente nas transições, $C$ é crescente em todo o domínio $[0,50]$. Assim, o valor mínimo ocorre em $m=0$:
$$C(0) = 1{,}5(0)+5 = 5$$
e o valor máximo ocorre em $m=50$:
$$C(50) = 60+4(50-30) = 60+80 = 140$$

Como a função é contínua e crescente, sua imagem é todo o intervalo entre esses extremos:
$$\text{Im}(C) = [5, 140]$$

\textbf{c) Consumo correspondente a $C(m) = 74$}

Precisamos identificar em qual sentença esse valor ocorre. Na primeira sentença, $C(m)$ varia de $5$ a $20$ (pois $m \in [0,10]$); na segunda, de $20$ a $60$ (pois $m \in (10,30]$). Como $74 > 60$, o valor $74$ só pode pertencer à imagem da terceira sentença, válida para $30 < m \le 50$.

Resolvendo nessa sentença:
$$60 + 4(m-30) = 74$$
$$4(m-30) = 14$$
$$m - 30 = 3{,}5$$
$$m = 33{,}5$$

Como $33{,}5 \in (30, 50]$, a solução é válida: o consumo foi de $33{,}5\,m^3$.

\textbf{d) Crescimento}

Como cada uma das três sentenças é estritamente crescente (coeficientes $1{,}5$, $2$ e $4$, todos positivos) e os valores coincidem exatamente nos pontos de transição $m=10$ e $m=30$, a função $C(m)$ é \textbf{crescente} em todo o intervalo $[0,50]$: quanto maior o consumo, maior a conta. Note, porém, que a *taxa* de aumento por metro cúbico cresce a cada faixa ($1{,}5 \to 2 \to 4$ reais por m³), caracterizando uma tarifa progressiva.

\bigskip\noindent\textbf{Q008 · múltipla escolha · fácil · entender · aplicada · conferido em parte}\par\nopagebreak
Uma companhia de saneamento calcula o valor mensal da conta de água residencial, em reais, de acordo com o consumo mensal $x$ (em metros cúbicos), conforme a tabela abaixo:

\begin{center}
\begin{tabular}{|p{0.430\linewidth}|p{0.430\linewidth}|}
\hline
\textbf{Faixa de consumo} & \textbf{Valor cobrado} \\
\hline
$0 \le x \le 10$ & Taxa mínima fixa de R\$\,25,00 \\
\hline
$10 < x \le 25$ & R\$\,25,00 mais R\$\,2,50 por cada $m^3$ que exceder 10 $m^3$ \\
\hline
$x > 25$ & R\$\,62,50 mais R\$\,4,00 por cada $m^3$ que exceder 25 $m^3$ \\
\hline
\end{tabular}
\end{center}

Seja $V(x)$ o valor da conta, em reais, para um consumo de $x$ $m^3$, com $x \ge 0$. Qual das alternativas descreve corretamente o gráfico de $V(x)$?
\begin{enumerate}[label=(\alph*)]
\item Para $0\le x\le10$, segmento horizontal em $25$. Para $10<x\le25$, segmento crescente que parte de $(10;25)$ e chega a $(25;85)$, com inclinação $4$. Para $x>25$, semirreta que parte de $(25;85)$, com inclinação $2{,}5$.
\item Para $0\le x\le10$, o gráfico é um segmento horizontal na altura $25$. Para $10<x\le25$, um segmento crescente que parte do ponto $(10;25)$ e chega ao ponto $(25;62{,}5)$, com inclinação $2{,}5$. Para $x>25$, uma semirreta crescente que parte de $(25;62{,}5)$, com inclinação $4$, sem limite superior.
\item Para $0\le x\le10$, segmento horizontal em $25$. Para $10<x\le25$, segmento crescente que parte do ponto $(10;0)$ e chega a $(25;37{,}5)$, com inclinação $2{,}5$. Para $x>25$, semirreta que parte de $(25;37{,}5)$, com inclinação $4$.
\item Para $0\le x\le10$, o gráfico é um segmento crescente que parte da origem $(0;0)$ até $(10;25)$, com inclinação $2{,}5$; a partir daí, o gráfico segue como nas demais faixas descritas na tabela.
\end{enumerate}

\textbf{Gabarito.} B

\textbf{Resolução.} \textbf{Passo 1 — Escrever a função por partes.}

A partir da tabela:
$$V(x)=\begin{cases}25, & 0\le x\le 10\\[2pt] 25+2{,}5(x-10), & 10<x\le 25\\[2pt] 62{,}5+4(x-25), & x>25\end{cases}$$

\textbf{Passo 2 — Verificar os pontos de transição (continuidade).}

Em $x=10$: a primeira sentença dá $V(10)=25$. A segunda sentença, calculada em $x=10$, dá $25+2{,}5\cdot 0=25$. As duas coincidem, então o segundo trecho \textbf{parte do ponto $(10;25)$}, e não de $(10;0)$.

Em $x=25$: $V(25)=25+2{,}5\cdot(25-10)=25+37{,}5=62{,}5$. A terceira sentença, em $x=25$, dá $62{,}5+4\cdot0=62{,}5$. Logo o terceiro trecho \textbf{parte do ponto $(25;62{,}5)$}.

\textbf{Passo 3 — Identificar as inclinações de cada trecho.}

\begin{itemize}[nosep,leftmargin=*]
\item Para $0\le x\le10$: $V$ é constante (inclinação $0$) — segmento horizontal em $y=25$.
\item Para $10<x\le25$: coeficiente de $x$ é $2{,}5$ — segmento crescente de $(10;25)$ até $(25;62{,}5)$.
\item Para $x>25$: coeficiente de $x$ é $4$ — semirreta crescente a partir de $(25;62{,}5)$, com inclinação maior que a do trecho anterior (a conta fica mais cara por m³ para quem consome mais).
\end{itemize}

\textbf{Passo 4 — Comparar com as alternativas.}

Apenas a alternativa que mostra o segmento horizontal em $25$ até $x=10$, seguido do segmento que \textbf{parte de $(10;25)$} com inclinação $2{,}5$ até $(25;62{,}5)$, e depois a semirreta que \textbf{parte de $(25;62{,}5)$} com inclinação $4$, está correta. As demais erram ao ignorar a taxa fixa acumulada, ao trocar as inclinações ou ao tratar a primeira faixa como proporcional ao consumo.

\textbf{Conclusão:} a descrição correta é a da alternativa B.

\bigskip\noindent\textbf{Q010 · múltipla escolha · média · aplicar · teórica · conferido em parte}\par\nopagebreak
Uma companhia de saneamento cobra a tarifa de água residencial, em reais, em função do consumo mensal $x$ (em $m^3$), segundo a lei:

$$f(x) = \begin{cases} 15, & 0 \le x \le 5 \\ 15 + 4(x-5), & x > 5 \end{cases}$$

em que $x$ representa o volume consumido no mês (um número real não negativo, podendo assumir qualquer valor, inclusive não inteiro).

Assinale a alternativa que descreve corretamente o gráfico de $f$, seu domínio de validade, sua imagem e seu comportamento de crescimento/decrescimento.
\begin{enumerate}[label=(\alph*)]
\item O gráfico é um segmento horizontal em $y=15$ para $0\le x\le 5$, seguido de uma semirreta crescente de coeficiente angular $4$ partindo do ponto $(5,15)$, sem salto; domínio $[0,+\infty)$, imagem $[15,+\infty)$, crescente para $x>5$.
\item O gráfico é um segmento horizontal em $y=15$ para $0\le x\le 5$, seguido de uma reta decrescente de coeficiente angular $-4$ a partir de $(5,15)$, com domínio $[0,+\infty)$ e imagem $(-\infty,15]$.
\item O gráfico é um segmento horizontal em $y=15$ até $x=5$, seguido de um salto até $y=35$ e depois uma reta crescente de coeficiente angular $4$, pois para $x>5$ deve-se usar diretamente $f(x)=15+4x$.
\item O gráfico é apenas o segmento horizontal em $y=15$, pois a função só está definida para $0\le x\le 5$; fora desse intervalo, $f$ não possui domínio de validade.
\end{enumerate}

\textbf{Gabarito.} O gráfico é um segmento horizontal em $y=15$ para $0\le x\le 5$, seguido de uma semirreta crescente de coeficiente angular $4$ partindo de $(5,15)$, sem descontinuidade; domínio $[0,+\infty)$, imagem $[15,+\infty)$, e a função é crescente para $x>5$.

\textbf{Resolução.} \textbf{1. Simplificando a segunda sentença}

Para $x>5$: $f(x) = 15 + 4(x-5) = 15 + 4x - 20 = 4x - 5$.

\textbf{2. Verificando a continuidade no ponto de troca ($x=5$)}

\begin{itemize}[nosep,leftmargin=*]
\item Pela primeira sentença: $f(5) = 15$.
\item Pela segunda sentença (limite quando $x\to 5^+$): $4(5)-5 = 15$.
\end{itemize}

Como os dois valores coincidem, o gráfico \textbf{não tem salto} em $x=5$: a reta crescente parte exatamente do ponto $(5,15)$, sem descontinuidade.

\textbf{3. Forma do gráfico}

\begin{itemize}[nosep,leftmargin=*]
\item Para $0 \le x \le 5$: segmento \textbf{horizontal} em $y=15$ (a tarifa mínima, constante).
\item Para $x>5$: \textbf{semirreta crescente}, pois o coeficiente angular é $4>0$, partindo de $(5,15)$ e subindo indefinidamente.
\end{itemize}

\textbf{4. Domínio de validade}

Como $x$ é um consumo (não pode ser negativo e não há limite superior imposto), o domínio é $[0, +\infty)$.

\textbf{5. Imagem}

O menor valor de $f$ é $15$ (atingido em todo o trecho $[0,5]$), e para $x>5$ os valores crescem sem limite. Logo, a imagem é $[15, +\infty)$.

\textbf{6. Crescimento}

No intervalo $[0,5]$ a função é constante; no intervalo $[5,+\infty)$ ela é \textbf{crescente}, pois o coeficiente angular da reta é positivo (igual a 4).

\textbf{Conclusão:} o gráfico é um segmento horizontal em $y=15$ até $x=5$, seguido de uma semirreta crescente de coeficiente angular $4$ a partir de $(5,15)$, sem salto; domínio $[0,+\infty)$; imagem $[15,+\infty)$; crescente para $x>5$.

\bigskip\noindent\textbf{Q012 · múltipla escolha · difícil · entender · aplicada · sem conferência automática: confira o gabarito}\par\nopagebreak
Uma companhia de saneamento cobra a conta mensal de água residencial segundo a tabela abaixo, válida apenas para consumos entre 0 e 35 m³ por mês, pois essa é a capacidade máxima de registro do hidrômetro residencial utilizado:

\begin{itemize}[nosep,leftmargin=*]
\item Consumo de 0 até 10 m³: tarifa mínima fixa de R\$\,25,00, independente do volume efetivamente consumido.
\item Consumo acima de 10 m³ e até 25 m³: R\$\,25,00 mais R\$\,3,50 por m³ que exceder 10 m³.
\item Consumo acima de 25 m³ e até 35 m³: o valor cobrado na faixa anterior no limite de 25 m³ (ou seja, R\$\,77,50) mais R\$\,5,00 por m³ que exceder 25 m³.
\end{itemize}

Seja $x$ o consumo mensal em metros cúbicos, com $0 \le x \le 35$, e $C(x)$ o valor da conta, em reais, correspondente a esse consumo.

Qual das alternativas a seguir descreve corretamente o gráfico de $C$ e a imagem dessa função no domínio de validade dado?
\begin{enumerate}[label=(\alph*)]
\item A função é crescente em todo o domínio $[0,35]$, sem nenhum trecho constante, já que o valor da conta aumenta proporcionalmente ao consumo desde $x=0$; a imagem é $[25;\,127{,}5]$.
\item Como a última faixa de cobrança (acima de 25 m³) não tem um valor de tarifa diferente após um certo ponto, o consumo pode crescer indefinidamente e a imagem da função é o intervalo ilimitado $[25, +\infty)$.
\item O gráfico é um segmento horizontal em $C=25$ para $0 \le x \le 10$; em seguida uma semirreta de coeficiente angular $3{,}5$ unindo $(10;25)$ a $(25;77{,}5)$; e por fim uma semirreta de coeficiente angular $5$ unindo $(25;77{,}5)$ a $(35;127{,}5)$. A função é constante em $[0,10]$ e crescente em $(10,35]$, com imagem $[25;\,127{,}5]$.
\item O gráfico é constante em $[0,10]$ e depois crescente, mas como a tarifa de R\$\,3,50 e a de R\$\,5,00 incidem sobre todo o consumo da faixa (e não apenas sobre o excedente), o valor máximo é $C(35) = 162{,}5$, de modo que a imagem é $[25;\,162{,}5]$.
\end{enumerate}

\textbf{Gabarito.} Alternativa C — $C$ é constante ($=25$) em $[0,10]$, crescente em $(10,35]$ (inclinação $3{,}5$ até $x=25$ e depois $5$), com imagem $[25;\,127{,}5]$.

\textbf{Resolução.} \textbf{Passo 1 — Escrever a expressão algébrica por partes.}

Para $0 \le x \le 10$: $C(x) = 25$.

Para $10 < x \le 25$: $C(x) = 25 + 3{,}5\,(x-10) = 3{,}5x - 10$.

Para $25 < x \le 35$: $C(x) = 77{,}5 + 5\,(x-25) = 5x - 47{,}5$.

\textbf{Passo 2 — Verificar a continuidade nas mudanças de sentença.}

Em $x=10$: a 1ª sentença dá $25$; a 2ª dá $3{,}5(10)-10 = 25$. Coincidem.

Em $x=25$: a 2ª sentença dá $3{,}5(25)-10 = 77{,}5$; a 3ª dá $5(25)-47{,}5 = 77{,}5$. Coincidem.

Logo o gráfico não tem saltos: é um segmento horizontal seguido de duas semirretas com inclinações crescentes ($3{,}5$ e depois $5$), unidas continuamente.

\textbf{Passo 3 — Analisar crescimento.}

No trecho $[0,10]$, $C$ é constante. No trecho $(10,35]$, as inclinações $3{,}5$ e $5$ são positivas, logo $C$ é estritamente crescente nesse intervalo.

\textbf{Passo 4 — Determinar a imagem no domínio $[0,35]$.}

O valor mínimo ocorre em $x=0$ (ou em qualquer ponto de $[0,10]$): $C=25$.

O valor máximo ocorre em $x=35$: $C(35) = 5(35) - 47{,}5 = 175 - 47{,}5 = 127{,}5$.

Como $C$ é contínua e não decrescente em todo o domínio, ela assume \textbf{todos} os valores entre $25$ e $127{,}5$. Portanto a imagem é o intervalo $[25;\,127{,}5]$.

\textbf{Passo 5 — Comparar com as alternativas.}

A alternativa que descreve corretamente o trecho constante em $[0,10]$, o crescimento em $(10,35]$ com as inclinações corretas e a imagem $[25;\,127{,}5]$ é a alternativa \textbf{A}.

As demais alternativas cometem erros: ignorar o trecho de tarifa fixa (constante), aplicar a tarifa marginal sobre o consumo total em vez de apenas sobre o excedente, ou desconsiderar o limite superior do domínio de validade (35 m³) imposto pelo hidrômetro.
\subsection*{EM13MAT501}
{\small\itshape Investigar relações entre números expressos em tabelas para representá-los no plano cartesiano, identificando padrões e criando conjecturas para generalizar e expressar algebricamente essa generalização, reconhecendo quando essa representação é de função polinomial de 1º grau.}

\bigskip\noindent\textbf{Q014 · discursiva · fácil · criar · aplicada · gabarito conferido}\par\nopagebreak
Uma locadora de bicicletas cobra uma taxa fixa de utilização mais um valor proporcional ao tempo de uso. A tabela a seguir mostra o valor total pago por diferentes clientes, de acordo com o tempo de uso da bicicleta:

\begin{center}
\begin{tabular}{|p{0.430\linewidth}|p{0.430\linewidth}|}
\hline
\textbf{Tempo de uso (horas)} & \textbf{Valor total pago (R\$)} \\
\hline
1 & 8 \\
\hline
2 & 12 \\
\hline
3 & 16 \\
\hline
4 & 20 \\
\hline
\end{tabular}
\end{center}

a) Observando o padrão de variação entre o tempo de uso e o valor pago, escreva uma expressão algébrica que permita calcular o valor total V, em reais, pago por um cliente que usa a bicicleta durante t horas, para qualquer valor inteiro positivo de t.

b) Classifique o tipo de função obtida em (a), justificando sua resposta a partir da expressão algébrica encontrada.

\textbf{Gabarito.} V(t) = 4t + 4; a relação é uma função polinomial do 1º grau (afim), pois apresenta taxa de variação constante (4) e t aparece com expoente 1.

\textbf{Resolução.} \textbf{Passo 1 — Observar a variação do valor pago em relação ao tempo.}

Calculando a diferença entre valores consecutivos da tabela:

$12-8=4$

$16-12=4$

$20-16=4$

A cada hora adicional, o valor pago aumenta sempre a mesma quantia: $4$ reais. Isso indica uma variação \textbf{constante}, ou seja, uma taxa de variação fixa igual a $4$.

\textbf{Passo 2 — Escrever a lei geral.}

Como a variação é constante, o valor total deve ter a forma

$$V(t) = 4t + b$$

em que $b$ é o valor pago quando $t=0$ (a taxa fixa da locadora).

\textbf{Passo 3 — Determinar $b$ usando um ponto da tabela.}

Usando o par $(1, 8)$:

$$8 = 4(1) + b \implies b = 4$$

Verificando com outro ponto, por exemplo $(3,16)$:

$$4(3)+4 = 12+4 = 16 \ \checkmark$$

Logo, a lei que representa a situação é:

$$V(t) = 4t + 4$$

\textbf{Passo 4 — Classificar a função.}

A expressão $V(t) = 4t + 4$ tem a forma $V(t) = a\,t + b$, com $a = 4 \neq 0$ e $b = 4$, isto é, a variável $t$ aparece elevada apenas ao expoente 1. Portanto, trata-se de uma \textbf{função polinomial do 1º grau (função afim)}, com coeficiente angular (taxa de variação) $4$ e coeficiente linear $4$. Como $a=4>0$, a função é crescente, o que é coerente com o fato de o valor pago aumentar conforme o tempo de uso aumenta.

\bigskip\noindent\textbf{Q016 · discursiva · média · analisar · teórica · gabarito conferido}\par\nopagebreak
Observe a tabela abaixo, que relaciona valores de uma grandeza $x$ a valores de uma grandeza $y$, obtidos a partir de um mesmo padrão:

\begin{center}
\begin{tabular}{|p{0.172\linewidth}|p{0.172\linewidth}|p{0.172\linewidth}|p{0.172\linewidth}|p{0.172\linewidth}|}
\hline
\textbf{$x$} & \textbf{1} & \textbf{2} & \textbf{3} & \textbf{4} \\
\hline
$y$ & 5 & 8 & 11 & 14 \\
\hline
\end{tabular}
\end{center}

a) Analise como $y$ varia à medida que $x$ aumenta uma unidade e identifique o padrão presente na tabela.

b) Escreva a lei algébrica $y = f(x)$ que generaliza esse padrão para qualquer valor real de $x$.

c) Classifique o tipo de função obtida, justificando sua resposta a partir da lei encontrada.

\textbf{Gabarito.} A relação é uma função afim dada por $f(x) = 3x + 2$, pois a variação de $y$ é constante (igual a 3) para cada acréscimo unitário de $x$.

\textbf{Resolução.} \textbf{Passo 1 — Investigar o padrão de variação de $y$}

Calculando as diferenças sucessivas dos valores de $y$ correspondentes a acréscimos unitários de $x$:

$8-5=3$

$11-8=3$

$14-11=3$

A variação de $y$ é \textbf{constante e igual a 3} para cada aumento de uma unidade em $x$. Isso indica que os pontos $(x,y)$ pertencem a uma reta, ou seja, que a relação entre $x$ e $y$ pode ser descrita por uma função polinomial de 1º grau, da forma $y = ax + b$.

\textbf{Passo 2 — Determinar o coeficiente angular $a$}

Como a variação de $y$ por unidade de $x$ é constante e vale 3, temos $a = 3$ (esse valor corresponde à razão da progressão aritmética formada pelos valores de $y$).

\textbf{Passo 3 — Determinar o coeficiente linear $b$}

Usando o par $(1,5)$ na lei $y = 3x + b$:

$5 = 3(1) + b \Rightarrow b = 5 - 3 = 2$

\textbf{Passo 4 — Escrever e verificar a lei}

$f(x) = 3x + 2$

Verificando nos demais pontos:

$f(2) = 3(2)+2 = 8$ $\checkmark$

$f(3) = 3(3)+2 = 11$ $\checkmark$

$f(4) = 3(4)+2 = 14$ $\checkmark$

Todos os pares da tabela satisfazem a lei encontrada.

\textbf{Passo 5 — Classificação}

Como $f(x) = 3x + 2$ tem a forma $y = ax + b$ com $a = 3 \neq 0$, a função é uma \textbf{função polinomial de 1º grau (função afim)}, cujo gráfico é uma reta não horizontal com coeficiente angular $3$ e coeficiente linear $2$.

\bigskip\noindent\textbf{Q018 · discursiva · difícil · analisar · aplicada · gabarito conferido}\par\nopagebreak
Uma empresa de saneamento básico cobra mensalmente pelo consumo de água uma tarifa composta por um valor fixo somado a um valor proporcional ao consumo, medido em metros cúbicos (m³). A tabela a seguir mostra o valor total pago por três clientes que tiveram consumos diferentes num mesmo mês:

\begin{center}
\begin{tabular}{|p{0.430\linewidth}|p{0.430\linewidth}|}
\hline
\textbf{Consumo (m³)} & \textbf{Valor cobrado (R\$)} \\
\hline
3 & 21,50 \\
\hline
7 & 37,50 \\
\hline
12 & 57,50 \\
\hline
\end{tabular}
\end{center}

a) Analise os dados da tabela e explique por que o valor cobrado V pode ser descrito por uma função polinomial do 1º grau do consumo x.

b) Determine a lei algébrica que expressa V em função de x.

c) Qual é o valor da tarifa fixa mensal, isto é, a parcela que o cliente paga mesmo que não haja variação proporcional ao consumo?

\textbf{Gabarito.} A relação é uma função afim porque a razão $\Delta V/\Delta x$ é constante e igual a 4, mesmo com intervalos de consumo diferentes. A lei é $V(x) = 4x + 9{,}5$, e a tarifa fixa é R\$\,9,50.

\textbf{Resolução.} \textbf{1. Organizando os pares (consumo, valor):}

$(3;\,21{,}5)$, $(7;\,37{,}5)$, $(12;\,57{,}5)$

Como os consumos não estão igualmente espaçados (as diferenças são $4$ e $5$, não uma constante como em uma tabela usual), não basta olhar as diferenças brutas: é preciso calcular a \textbf{taxa de variação} $\dfrac{\Delta V}{\Delta x}$ entre pares consecutivos.

\textbf{2. Calculando as taxas de variação:}

Entre o 1º e o 2º cliente:
$$\Delta x_1 = 7-3=4, \qquad \Delta V_1 = 37{,}5-21{,}5=16$$
$$\frac{\Delta V_1}{\Delta x_1}=\frac{16}{4}=4$$

Entre o 2º e o 3º cliente:
$$\Delta x_2 = 12-7=5, \qquad \Delta V_2 = 57{,}5-37{,}5=20$$
$$\frac{\Delta V_2}{\Delta x_2}=\frac{20}{5}=4$$

\textbf{3. Reconhecendo a função afim:}

Como a razão $\dfrac{\Delta V}{\Delta x}$ é a mesma ($=4$) mesmo para intervalos de tamanhos diferentes, a variação de $V$ é diretamente proporcional à variação de $x$. Isso caracteriza uma \textbf{função polinomial do 1º grau (afim)}, do tipo $V(x)=ax+b$, com $a=4$ (a taxa constante encontrada).

\textbf{4. Determinando o coeficiente linear $b$:}

Usando o ponto $(3;\,21{,}5)$:
$$21{,}5 = 4\cdot 3 + b \;\Rightarrow\; 21{,}5 = 12 + b \;\Rightarrow\; b = 9{,}5$$

Logo:
$$V(x) = 4x + 9{,}5$$

\textbf{5. Verificação com o terceiro par:}

$$V(12) = 4\cdot 12 + 9{,}5 = 48+9{,}5 = 57{,}5 \;\checkmark$$

Confirma-se que a lei encontrada é consistente com todos os dados da tabela.

\textbf{6. Interpretando a tarifa fixa:}

A parte que não depende do consumo é o coeficiente linear $b$, ou seja, o valor de $V$ quando $x=0$:
$$V(0) = 9{,}5$$

Portanto, a tarifa fixa mensal é \textbf{R\$\,9,50}.

\bigskip\noindent\textbf{Q020 · múltipla escolha · fácil · criar · teórica · gabarito conferido}\par\nopagebreak
A tabela a seguir apresenta pares de valores (x, y) gerados por uma mesma regra de formação:

\begin{center}
\begin{tabular}{|p{0.143\linewidth}|p{0.143\linewidth}|p{0.143\linewidth}|p{0.143\linewidth}|p{0.143\linewidth}|p{0.143\linewidth}|}
\hline
\textbf{x} & \textbf{-2} & \textbf{-1} & \textbf{0} & \textbf{1} & \textbf{2} \\
\hline
y & 1 & 3 & 5 & 7 & 9 \\
\hline
\end{tabular}
\end{center}

a) Investigue os pares (x, y) da tabela, identifique um padrão entre os valores de x e y e escreva, com suas próprias palavras, a expressão algébrica y = f(x) que gera esses valores. Justifique por escrito como você chegou a essa expressão.

b) Com base no padrão observado, classifique o tipo de função obtida (por exemplo: constante, polinomial do 1º grau, polinomial do 2º grau etc.) e justifique sua classificação.

c) Assinale a alternativa que apresenta corretamente a lei de formação encontrada.

(Os itens a e b devem ser respondidos de forma discursiva e serão corrigidos separadamente, avaliando o processo de investigação e a justificativa apresentada; o item c será corrigido de forma objetiva.)
\begin{enumerate}[label=(\alph*)]
\item y = 2x + 5
\item y = 2x
\item y = 2x + 1
\item y = x² + 5
\end{enumerate}

\textbf{Gabarito.} y = 2x + 5 (alternativa A)

\textbf{Resolução.} \textbf{Item a) e b) — investigação do padrão}

Organizando as diferenças sucessivas dos valores de $y$ para cada aumento de 1 unidade em $x$:

$3-1=2,\quad 5-3=2,\quad 7-5=2,\quad 9-7=2$

A diferença é \textbf{constante e igual a 2} sempre que $x$ aumenta 1 unidade. Isso indica que a taxa de variação $\dfrac{\Delta y}{\Delta x}$ é constante, o que é a característica que define uma \textbf{função polinomial do 1º grau (função afim)}, do tipo $y = ax+b$, com $a\neq 0$.

Como a taxa de variação é $a=2$, e pelo par $(0,5)$ temos $y=5$ quando $x=0$, segue que $b=5$.

Logo, a lei de formação é:

$$y = 2x+5$$

Verificação com os demais pontos: $x=-2 \Rightarrow y=2(-2)+5=1$ $\checkmark$; $x=1 \Rightarrow y=2(1)+5=7$ $\checkmark$; $x=2 \Rightarrow y=2(2)+5=9$ $\checkmark$.

Como $a=2\neq 0$ e o expoente de $x$ é 1, confirma-se que se trata de uma \textbf{função polinomial do 1º grau}.

\textbf{Item c)}

A alternativa que apresenta corretamente essa lei é $y=2x+5$.

\bigskip\noindent\textbf{Q022 · múltipla escolha · média · criar · aplicada · gabarito conferido}\par\nopagebreak
Uma empresa de entregas rápidas registrou o valor cobrado (em reais) para diferentes distâncias percorridas por seus motoboys, conforme a tabela abaixo:

\begin{center}
\begin{tabular}{|p{0.430\linewidth}|p{0.430\linewidth}|}
\hline
\textbf{Distância d (km)} & \textbf{Valor cobrado V (R\$)} \\
\hline
2 & 12 \\
\hline
5 & 21 \\
\hline
8 & 30 \\
\hline
11 & 39 \\
\hline
\end{tabular}
\end{center}

Analisando os dados da tabela, qual expressão algébrica V(d) generaliza corretamente o padrão observado, e que tipo de função ela representa?
\begin{enumerate}[label=(\alph*)]
\item V(d) = 3d + 6 — função polinomial do 1º grau
\item V(d) = 6d — função polinomial do 1º grau
\item V(d) = 3d² + 6 — função polinomial do 2º grau
\item V(d) = 9d - 6 — função polinomial do 1º grau
\end{enumerate}

\textbf{Gabarito.} V(d) = 3d + 6 — função polinomial do 1º grau (afim)

\textbf{Resolução.} \textbf{Passo 1 — Organizar os dados e observar as variações}

Calculando a variação de $d$ e de $V$ entre linhas consecutivas da tabela:

$\Delta d = 5-2 = 3,\quad \Delta V = 21-12 = 9$

$\Delta d = 8-5 = 3,\quad \Delta V = 30-21 = 9$

$\Delta d = 11-8 = 3,\quad \Delta V = 39-30 = 9$

\textbf{Passo 2 — Calcular a taxa de variação}

Como $\Delta d$ é sempre igual e $\Delta V$ também é sempre igual, a razão $\dfrac{\Delta V}{\Delta d}$ é constante:

$\dfrac{\Delta V}{\Delta d} = \dfrac{9}{3} = 3$

Uma taxa de variação constante entre grandezas indica que a relação entre $d$ e $V$ é do tipo $V(d) = a\cdot d + b$, ou seja, uma \textbf{função polinomial do 1º grau (função afim)}, com $a = 3$.

\textbf{Passo 3 — Determinar o coeficiente linear $b$}

Usando o ponto $(2,12)$:

$12 = 3\cdot 2 + b \Rightarrow 12 = 6 + b \Rightarrow b = 6$

\textbf{Passo 4 — Escrever a lei de formação e verificar}

$V(d) = 3d + 6$

Verificando com os demais pontos:

\begin{itemize}[nosep,leftmargin=*]
\item $d=5$: $3(5)+6=21$ $\checkmark$
\item $d=8$: $3(8)+6=30$ $\checkmark$
\item $d=11$: $3(11)+6=39$ $\checkmark$
\end{itemize}

Todos os pontos satisfazem a expressão, confirmando que a relação é de \textbf{função polinomial do 1º grau}, com lei $V(d) = 3d + 6$.

\bigskip\noindent\textbf{Q024 · múltipla escolha · difícil · criar · teórica · gabarito conferido}\par\nopagebreak
A tabela abaixo relaciona valores de $x$ e os correspondentes valores de $y$, obtidos a partir de uma mesma regra de formação (os valores de $x$ não variam em intervalos iguais):

\begin{center}
\begin{tabular}{|p{0.172\linewidth}|p{0.172\linewidth}|p{0.172\linewidth}|p{0.172\linewidth}|p{0.172\linewidth}|}
\hline
\textbf{$x$} & \textbf{$-3$} & \textbf{$1$} & \textbf{$4$} & \textbf{$8$} \\
\hline
$y$ & $14$ & $2$ & $-7$ & $-19$ \\
\hline
\end{tabular}
\end{center}

Um estudante, analisando os pares ordenados $(x,y)$ da tabela, conjectura que essa relação pode ser representada por uma função polinomial do 1º grau. Assinale a alternativa que apresenta a lei de formação $f(x)$ que generaliza corretamente o padrão observado na tabela.
\begin{enumerate}[label=(\alph*)]
\item $f(x) = -11x + 13$
\item $f(x) = -12x + 14$
\item $f(x) = -3x + 5$
\item $f(x) = 3x - 1$
\end{enumerate}

\textbf{Gabarito.} f(x) = -3x + 5

\textbf{Resolução.} \textbf{Passo 1 — Verificar se a razão entre as variações é constante (condição para função afim).}

Como os valores de $x$ não estão igualmente espaçados, é preciso calcular a razão $\dfrac{\Delta y}{\Delta x}$ entre pares consecutivos da tabela, e não apenas $\Delta y$.

Entre $x=-3$ e $x=1$: $\Delta x = 1-(-3)=4$, $\Delta y = 2-14=-12$, logo $\dfrac{\Delta y}{\Delta x}=\dfrac{-12}{4}=-3$.

Entre $x=1$ e $x=4$: $\Delta x = 4-1=3$, $\Delta y = -7-2=-9$, logo $\dfrac{\Delta y}{\Delta x}=\dfrac{-9}{3}=-3$.

Entre $x=4$ e $x=8$: $\Delta x = 8-4=4$, $\Delta y = -19-(-7)=-12$, logo $\dfrac{\Delta y}{\Delta x}=\dfrac{-12}{4}=-3$.

\textbf{Passo 2 — Concluir o tipo de função.}

A razão $\dfrac{\Delta y}{\Delta x}$ é constante e igual a $-3$ para todos os pares consecutivos, mesmo com intervalos diferentes de $x$. Isso caracteriza uma \textbf{função polinomial do 1º grau (afim)}, da forma $f(x)=mx+b$, com $m=-3$.

\textbf{Passo 3 — Determinar o coeficiente linear $b$.}

Usando o ponto $(1,2)$: $2 = -3(1)+b \Rightarrow b = 2+3 = 5$.

Logo, $f(x) = -3x+5$.

\textbf{Passo 4 — Verificar com os demais pontos.}

$f(-3) = -3(-3)+5 = 9+5 = 14$ $\checkmark$

$f(4) = -3(4)+5 = -12+5 = -7$ $\checkmark$

$f(8) = -3(8)+5 = -24+5 = -19$ $\checkmark$

Todos os pares confirmam a lei $f(x) = -3x+5$.
\subsection*{EM13MAT502}
{\small\itshape Investigar relações entre números expressos em tabelas para representá-los no plano cartesiano, identificando padrões e criando conjecturas para generalizar e expressar algebricamente essa generalização, reconhecendo quando essa representação é de função polinomial de 2º grau do tipo y = ax².}

\bigskip\noindent\textbf{Q027 · discursiva · fácil · analisar · teórica · gabarito conferido}\par\nopagebreak
Observe a tabela abaixo, que relaciona valores de x com valores de y:

\begin{center}
\begin{tabular}{|p{0.172\linewidth}|p{0.172\linewidth}|p{0.172\linewidth}|p{0.172\linewidth}|p{0.172\linewidth}|}
\hline
\textbf{x} & \textbf{1} & \textbf{2} & \textbf{3} & \textbf{4} \\
\hline
y & 4 & 16 & 36 & 64 \\
\hline
\end{tabular}
\end{center}

a) Investigue como y varia em função de x, testando se há proporcionalidade direta, proporcionalidade com o quadrado de x, ou outra relação.

b) Escreva a lei algébrica y = f(x) que generaliza o padrão observado na tabela, e classifique o tipo de função obtida.

\textbf{Gabarito.} y = 4x²

\textbf{Resolução.} \textbf{Passo 1 — Testar proporcionalidade direta (y = ax):}

Se fosse y proporcional a x, a razão $y/x$ deveria ser constante. Calculando:

$\frac{4}{1}=4,\quad \frac{16}{2}=8,\quad \frac{36}{3}=12,\quad \frac{64}{4}=16$

Como essas razões não são constantes, y não é diretamente proporcional a x.

\textbf{Passo 2 — Testar proporcionalidade ao quadrado de x (y = ax²):}

Calculamos a razão $y/x^2$ para cada par:

$\frac{4}{1^2}=4,\quad \frac{16}{2^2}=4,\quad \frac{36}{3^2}=4,\quad \frac{64}{4^2}=4$

A razão é constante e igual a $4$ em todos os casos, o que indica que y é proporcional ao quadrado de x.

\textbf{Passo 3 — Escrever a lei geral:}

Como $\frac{y}{x^2}=4$ para todos os pares da tabela, concluímos que:

$$y = 4x^2$$

\textbf{Passo 4 — Verificação:}

$x=1: y=4(1)^2=4$ $\checkmark$
$x=2: y=4(2)^2=16$ $\checkmark$
$x=3: y=4(3)^2=36$ $\checkmark$
$x=4: y=4(4)^2=64$ $\checkmark$

\textbf{Conclusão:} A lei $y=4x^2$ é uma função polinomial do 2º grau do tipo $y=ax^2$ (sem termos de grau 1 ou constante), pois y é diretamente proporcional ao quadrado de x, com constante de proporcionalidade $a=4$.

\bigskip\noindent\textbf{Q030 · discursiva · média · analisar · aplicada · gabarito conferido}\par\nopagebreak
Um engenheiro testa a queda livre de uma pequena esfera solta de um drone e mede a distância total percorrida pela esfera (em metros) para diferentes tempos de queda (em segundos), obtendo a tabela abaixo:

\begin{center}
\begin{tabular}{|p{0.430\linewidth}|p{0.430\linewidth}|}
\hline
\textbf{Tempo $t$ (s)} & \textbf{Distância $d$ (m)} \\
\hline
1 & 5 \\
\hline
2 & 20 \\
\hline
3 & 45 \\
\hline
4 & 80 \\
\hline
\end{tabular}
\end{center}

a) Analisando os valores da tabela, identifique o padrão que relaciona $d$ e $t$ e escreva a lei algébrica que expressa $d$ em função de $t$.

b) Usando essa lei, determine a distância percorrida pela esfera após 5 segundos de queda.

\textbf{Gabarito.} a) $d(t) = 5t^2$ (a distância é proporcional ao quadrado do tempo). b) $d(5) = 125$ m.

\textbf{Resolução.} \textbf{Passo 1 — Investigar o padrão na tabela.}

Observe que $d$ não cresce na mesma proporção que $t$ (não é uma reta passando pela origem com razão constante $d/t$), pois:

$\dfrac{20}{5}=4$ mas $\dfrac{2}{1}=2$ — logo a relação não é linear simples ($d = kt$).

\textbf{Passo 2 — Testar a razão entre $d$ e $t^2$.}

Calculando $t^2$ para cada linha: $1, 4, 9, 16$.

Agora dividimos $d$ por $t^2$:

$\dfrac{5}{1}=5,\quad \dfrac{20}{4}=5,\quad \dfrac{45}{9}=5,\quad \dfrac{80}{16}=5$

A razão $d/t^2$ é constante e igual a $5$ em todos os casos.

\textbf{Passo 3 — Generalizar algebricamente.}

Como $\dfrac{d}{t^2}=5$ para todo par da tabela, concluímos que $d$ é diretamente proporcional ao quadrado de $t$, ou seja, a relação é do tipo $d = a\,t^2$, com $a=5$:

$$d(t) = 5t^2$$

Verificação: $d(1)=5(1)^2=5$; $d(2)=5(4)=20$; $d(3)=5(9)=45$; $d(4)=5(16)=80$. Todos conferem com a tabela.

\textbf{Passo 4 — Calcular a distância para $t=5$ s.}

$$d(5) = 5\cdot 5^2 = 5\cdot 25 = 125 \text{ m}$$

\bigskip\noindent\textbf{Q033 · discursiva · difícil · criar · teórica · gabarito conferido}\par\nopagebreak
Observe a tabela abaixo, que relaciona valores de $x$ e $y$:

\begin{center}
\begin{tabular}{|p{0.172\linewidth}|p{0.172\linewidth}|p{0.172\linewidth}|p{0.172\linewidth}|p{0.172\linewidth}|}
\hline
\textbf{$x$} & \textbf{$-3$} & \textbf{$-1$} & \textbf{$2$} & \textbf{$4$} \\
\hline
$y$ & $27$ & $3$ & $12$ & $48$ \\
\hline
\end{tabular}
\end{center}

Analisando os pares de valores, investigue como $y$ varia em função de $x$ e determine a expressão algébrica que generaliza essa relação para qualquer valor de $x$. Em seguida, classifique o tipo de função polinomial obtida, justificando sua resposta.

\textbf{Gabarito.} y = 3x², uma função quadrática do tipo y = ax² (com a = 3)

\textbf{Resolução.} \textbf{Passo 1 — Observar o espaçamento dos valores de $x$.}

Os valores de $x$ dados são $-3, -1, 2, 4$. As diferenças entre valores consecutivos de $x$ são $2, 3, 2$ — ou seja, \textbf{não são igualmente espaçados}. Isso significa que o método usual de calcular diferenças sucessivas de $y$ (útil quando $x$ forma uma PA) não pode ser aplicado diretamente. É preciso investigar a relação de outra forma.

\textbf{Passo 2 — Testar proporcionalidade direta ($y = bx$).}

Calculando a razão $\dfrac{y}{x}$ para cada par:

$$\frac{27}{-3}=-9,\quad \frac{3}{-1}=-3,\quad \frac{12}{2}=6,\quad \frac{48}{4}=12$$

As razões não são iguais, logo $y$ \textbf{não} é diretamente proporcional a $x$ — a relação não é do tipo $y=bx$.

\textbf{Passo 3 — Testar proporcionalidade ao quadrado de $x$ ($y = ax^2$).}

Calculando $x^2$ para cada valor: $9, 1, 4, 16$.

Calculando a razão $\dfrac{y}{x^2}$:

$$\frac{27}{9}=3,\quad \frac{3}{1}=3,\quad \frac{12}{4}=3,\quad \frac{48}{16}=3$$

Todas as razões são iguais a $3$! Isso indica que $y$ é proporcional ao quadrado de $x$, com constante de proporcionalidade $a=3$.

\textbf{Passo 4 — Escrever e verificar a generalização.}

A lei que expressa o padrão é:

$$y = 3x^2$$

Verificando com todos os pares da tabela:
\begin{itemize}[nosep,leftmargin=*]
\item $3\cdot(-3)^2 = 3\cdot 9 = 27$ $\checkmark$
\item $3\cdot(-1)^2 = 3\cdot 1 = 3$ $\checkmark$
\item $3\cdot(2)^2 = 3\cdot 4 = 12$ $\checkmark$
\item $3\cdot(4)^2 = 3\cdot 16 = 48$ $\checkmark$
\end{itemize}

\textbf{Passo 5 — Classificação.}

Como $y$ depende de $x$ através de $y=ax^2$ (com $a=3\neq 0$), trata-se de uma \textbf{função polinomial do 2º grau (função quadrática)} da forma particular $y=ax^2$, cujo gráfico é uma parábola com vértice na origem $(0,0)$.

\bigskip\noindent\textbf{Q036 · múltipla escolha · fácil · analisar · aplicada · gabarito conferido}\par\nopagebreak
Uma equipe de engenharia testa, em túnel de vento, a força de resistência do ar (F, em newtons) que atua sobre uma peça, variando a velocidade do ar (v, em m/s) incidente sobre ela. Os resultados obtidos estão na tabela abaixo:

\begin{center}
\begin{tabular}{|p{0.172\linewidth}|p{0.172\linewidth}|p{0.172\linewidth}|p{0.172\linewidth}|p{0.172\linewidth}|}
\hline
\textbf{v (m/s)} & \textbf{1} & \textbf{2} & \textbf{3} & \textbf{4} \\
\hline
F (N) & 3 & 12 & 27 & 48 \\
\hline
\end{tabular}
\end{center}

Analisando os dados da tabela, qual expressão algébrica representa corretamente a força de resistência F em função da velocidade v?
\begin{enumerate}[label=(\alph*)]
\item F = 3v
\item F = v² + 2
\item F = 3v²
\item F = 6v²
\end{enumerate}

\textbf{Gabarito.} F(v) = 3v²

\textbf{Resolução.} \textbf{Passo 1 — Testar proporcionalidade direta (linear):}

Calculamos a razão $\dfrac{F}{v}$ para cada par: $\dfrac{3}{1}=3$, $\dfrac{12}{2}=6$, $\dfrac{27}{3}=9$, $\dfrac{48}{4}=12$.

Como essa razão \textbf{não é constante}, a relação não é do tipo $F = av$ (não é linear/proporcional direta simples).

\textbf{Passo 2 — Analisar as diferenças sucessivas de F:}

$12-3=9$, $27-12=15$, $48-27=21$ (1ª diferença não é constante, confirmando que não é linear).

As diferenças dessas diferenças (2ª diferença): $15-9=6$ e $21-15=6$ — \textbf{constante}! Isso é a assinatura de uma função quadrática do tipo $y=ax^2+bx+c$.

\textbf{Passo 3 — Testar proporcionalidade ao quadrado de v:}

Calculamos $\dfrac{F}{v^2}$ para cada par: $\dfrac{3}{1}=3$, $\dfrac{12}{4}=3$, $\dfrac{27}{9}=3$, $\dfrac{48}{16}=3$.

A razão é \textbf{constante e igual a 3}, o que mostra que $F$ é diretamente proporcional a $v^2$, isto é, $F = 3v^2$ (sem termos lineares ou constantes, já que fisicamente quando $v=0$ não há resistência).

\textbf{Passo 4 — Verificar:}

$3(1)^2=3$ $\checkmark$, $3(2)^2=12$ $\checkmark$, $3(3)^2=27$ $\checkmark$, $3(4)^2=48$ $\checkmark$.

Logo, $F(v) = 3v^2$.

\bigskip\noindent\textbf{Q039 · múltipla escolha · média · criar · teórica · gabarito conferido}\par\nopagebreak
A tabela abaixo relaciona o comprimento $L$ do lado de um certo tipo de figura com a grandeza $A$ associada a ela, para quatro valores de $L$:

\begin{center}
\begin{tabular}{|p{0.430\linewidth}|p{0.430\linewidth}|}
\hline
\textbf{$L$} & \textbf{$A$} \\
\hline
1 & 3 \\
\hline
2 & 12 \\
\hline
3 & 27 \\
\hline
4 & 48 \\
\hline
\end{tabular}
\end{center}

Investigando os pares de valores dessa tabela, qual expressão algébrica, escrita em função de $L$, generaliza corretamente a relação entre $A$ e $L$ para todos os pares apresentados (e não apenas para alguns deles)?
\begin{enumerate}[label=(\alph*)]
\item $A = 3L^2$
\item $A = 3L$
\item $A = 9L - 6$
\item $A = L^2 + 2$
\end{enumerate}

\textbf{Gabarito.} A = 3L²

\textbf{Resolução.} \textbf{Passo 1 — Testar um modelo linear simples.}
Se a relação fosse do tipo $A = kL$, o valor de $k = A/L$ deveria ser constante. Calculando: $3/1=3$, $12/2=6$, $27/3=9$, $48/4=12$. Como $k$ não é constante, a relação \textbf{não é linear} (não é do tipo $A=kL$).

\textbf{Passo 2 — Observar como $A$ cresce em relação a $L$.}
As diferenças sucessivas de $A$ são: $12-3=9$, $27-12=15$, $48-27=21$. Essas diferenças não são constantes (o que confirmaria uma reta), mas crescem de forma regular: $15-9=6$ e $21-15=6$. Uma segunda diferença constante é a marca característica de uma função \textbf{quadrática}.

\textbf{Passo 3 — Testar um modelo do tipo $A = aL^2$.}
Como a razão $A/L$ não é constante, mas $A$ parece crescer com o quadrado de $L$, testamos a razão $A/L^2$: $3/1^2=3$, $12/2^2=3$, $27/3^2=3$, $48/4^2=3$. Essa razão é constante e igual a $3$ para todos os pares.

\textbf{Passo 4 — Escrever a generalização.}
Como $A/L^2 = 3$ para todos os pontos da tabela, concluímos que $A = 3L^2$.

\textbf{Verificação final:} $L=1\Rightarrow 3(1)^2=3$; $L=2\Rightarrow 3(4)=12$; $L=3\Rightarrow 3(9)=27$; $L=4\Rightarrow 3(16)=48$. Todos os pares conferem, confirmando que a relação é do tipo $y=ax^2$, com $a=3$.

\bigskip\noindent\textbf{Q042 · múltipla escolha · difícil · analisar · aplicada · gabarito conferido}\par\nopagebreak
Uma equipe de engenharia rodoviária testa, em pista seca, a distância de frenagem d (em metros) de um automóvel a partir do instante em que o motorista aciona os freios, para diferentes velocidades v (em km/h). Os resultados de quatro testes estão na tabela abaixo:

\begin{center}
\begin{tabular}{|p{0.172\linewidth}|p{0.172\linewidth}|p{0.172\linewidth}|p{0.172\linewidth}|p{0.172\linewidth}|}
\hline
\textbf{v (km/h)} & \textbf{10} & \textbf{20} & \textbf{30} & \textbf{40} \\
\hline
d (m) & 2 & 8 & 18 & 32 \\
\hline
\end{tabular}
\end{center}

Sabendo que a distância de frenagem depende apenas da velocidade e que o padrão observado nos testes se mantém para qualquer velocidade, qual expressão algébrica representa corretamente d em função de v?
\begin{enumerate}[label=(\alph*)]
\item $d(v) = \dfrac{v^2}{25}$
\item $d(v) = 0{,}6v - 4$
\item $d(v) = 0{,}2v$
\item $d(v) = \dfrac{v^2}{50}$
\end{enumerate}

\textbf{Gabarito.} d(v) = v²/50

\textbf{Resolução.} \textbf{Passo 1 — Testar se a relação é linear.}

As diferenças sucessivas de $d$ são: $8-2=6$, $18-8=10$, $32-18=14$. Como essas diferenças não são constantes, a relação \textbf{não é linear} (não é do tipo $d = mv + n$).

\textbf{Passo 2 — Testar as segundas diferenças.}

As diferenças das diferenças são: $10-6=4$ e $14-10=4$. Como as segundas diferenças são constantes (iguais a 4) para um passo constante de $v$ ($\Delta v = 10$), isso indica uma relação \textbf{quadrática}, ou seja, $d = av^2 + bv + c$.

\textbf{Passo 3 — Testar se a relação é do tipo $d = av^2$ (sem termos de grau menor).}

Calculamos a razão $\dfrac{d}{v^2}$ para cada par:

$\dfrac{2}{10^2} = \dfrac{2}{100} = \dfrac{1}{50}$

$\dfrac{8}{20^2} = \dfrac{8}{400} = \dfrac{1}{50}$

$\dfrac{18}{30^2} = \dfrac{18}{900} = \dfrac{1}{50}$

$\dfrac{32}{40^2} = \dfrac{32}{1600} = \dfrac{1}{50}$

Como a razão $\dfrac{d}{v^2}$ é constante e igual a $\dfrac{1}{50}$ em \textbf{todos} os pontos (e não passa por $v=0$ com $d\neq 0$ nem apresenta termo linear), concluímos que não há termos de grau 1 nem constante: a relação é exatamente do tipo $y = ax^2$.

\textbf{Passo 4 — Escrever a expressão geral.}

$$d(v) = \dfrac{v^2}{50}$$

Verificação final: $d(10)=\frac{100}{50}=2$, $d(20)=\frac{400}{50}=8$, $d(30)=\frac{900}{50}=18$, $d(40)=\frac{1600}{50}=32$ — todos coincidem com a tabela.
\subsection*{EM13MAT503}
{\small\itshape Investigar pontos de máximo ou de mínimo de funções quadráticas em contextos da Matemática Financeira ou da Cinemática, entre outros.}

\bigskip\noindent\textbf{Q050 · discursiva · difícil · aplicar · aplicada · conferido em parte}\par\nopagebreak
Uma agência organiza passeios de van para grupos de turistas. A van comporta no máximo 60 passageiros, e a agência só realiza o passeio se houver pelo menos 15 inscritos.

O preço da passagem depende do tamanho do grupo: para 15 passageiros, o preço é R\$\,90,00 por pessoa. A cada passageiro adicional além de 15, a agência reduz o preço da passagem em R\$\,1,00 para TODOS os integrantes do grupo (não apenas para o excedente), como forma de incentivar grupos maiores.

Além da receita das passagens, a agência tem um custo de R\$\,500,00 fixo pelo aluguel da van, mais R\$\,20,00 por passageiro (referentes a seguro e lanche).

Determine o número de passageiros que a agência deve buscar para cada passeio, de modo a maximizar seu lucro, e calcule o valor desse lucro máximo.

\textbf{Gabarito.} O lucro é máximo para 42 ou 43 passageiros, resultando em L = R\$\,1.306,00.

\textbf{Resolução.} \textbf{Passo 1 — Definir a variável e o preço por passageiro}

Seja $x$ o número de passageiros, com $x$ inteiro e $15 \le x \le 60$.

O preço inicial é R\$\,90,00 para $x=15$, e cai R\$\,1,00 para cada passageiro além de 15:
$$p(x) = 90 - 1\cdot(x-15) = 105 - x$$

\textbf{Passo 2 — Montar a receita}

$$R(x) = x \cdot p(x) = x(105-x) = 105x - x^2$$

\textbf{Passo 3 — Montar o custo}

$$C(x) = 500 + 20x$$

\textbf{Passo 4 — Montar o lucro}

$$L(x) = R(x) - C(x) = 105x - x^2 - 500 - 20x = -x^2 + 85x - 500$$

Essa é uma função quadrática com $a=-1<0$, portanto tem ponto de \textbf{máximo}.

\textbf{Passo 5 — Localizar o vértice (candidato contínuo)}

$$x_v = -\dfrac{b}{2a} = -\dfrac{85}{2\cdot(-1)} = \dfrac{85}{2} = 42{,}5$$

\textbf{Passo 6 — Ajustar ao domínio (número de passageiros é inteiro)}

Como $x$ deve ser um número inteiro de pessoas, o vértice $x=42{,}5$ não é um valor válido. Como a parábola é simétrica em torno de $x_v=42{,}5$, os dois inteiros mais próximos, $x=42$ e $x=43$, estão à mesma distância ($0{,}5$) do vértice e, portanto, produzem o \textbf{mesmo valor máximo} de $L(x)$ dentro do domínio inteiro.

\textbf{Passo 7 — Calcular o lucro nesses pontos}

$$L(42) = -(42)^2 + 85(42) - 500 = -1764 + 3570 - 500 = 1306$$
$$L(43) = -(43)^2 + 85(43) - 500 = -1849 + 3655 - 500 = 1306$$

Ambos os valores coincidem, confirmando a simetria.

\textbf{Passo 8 — Confirmar que é de fato o máximo no domínio permitido}

Como a parábola abre para baixo e $x=42$ e $x=43$ estão dentro do intervalo permitido $[15,60]$ e são os inteiros mais próximos do vértice contínuo, nenhum outro valor inteiro no domínio produz lucro maior (o lucro decresce à medida que $x$ se afasta de $42{,}5$ para qualquer lado).

\textbf{Conclusão}

A agência deve buscar \textbf{42 ou 43 passageiros} por passeio, obtendo um lucro máximo de \textbf{R\$\,1.306,00}.

\bigskip\noindent\textbf{Q058 · múltipla escolha · difícil · aplicar · teórica · conferido em parte}\par\nopagebreak
Uma agência de viagens vende pacotes turísticos a um preço unitário de $x$ reais. Um estudo de mercado mostrou que a quantidade de pacotes vendidos por mês, em função do preço, é dada por $q(x) = 1000 - 4x$ (em unidades). Por cláusula do contrato firmado com a operadora parceira, o preço do pacote deve obrigatoriamente estar entre R$\,200,00$ e R$\,300,00$, incluindo os extremos. Considerando a receita mensal $R(x) = x \cdot q(x)$, determine o preço que maximiza a receita da agência dentro dessa faixa contratual e o valor dessa receita máxima.
\begin{enumerate}[label=(\alph*)]
\item R\$\,200,00, com receita máxima de R\$\,40.000,00
\item R\$\,300,00, com receita máxima de R\$-60.000,00
\item R\$\,125,00, com receita máxima de R\$\,62.500,00
\item R\$\,250,00, com receita máxima de R\$\,0,00
\end{enumerate}

\textbf{Gabarito.} Preço de R\$\,200,00, com receita máxima de R\$\,40.000,00 (alternativa A)

\textbf{Resolução.} \textbf{Passo 1 — Modelar a receita.}

A receita mensal é $R(x) = x \cdot q(x) = x(1000 - 4x) = 1000x - 4x^2$.

Essa é uma função quadrática com $a = -4$, $b = 1000$, portanto a parábola tem concavidade voltada para baixo.

\textbf{Passo 2 — Encontrar o vértice da parábola (sem restrição).}

$$x_v = -\dfrac{b}{2a} = -\dfrac{1000}{2(-4)} = 125$$

$$R(125) = 1000(125) - 4(125)^2 = 125000 - 62500 = 62500$$

Sem restrições, a receita máxima ocorreria em $x = 125$, com $R = 62500$.

\textbf{Passo 3 — Verificar a restrição contratual.}

O preço deve pertencer ao intervalo $[200, 300]$. Como $x_v = 125 < 200$, o vértice \textbf{não pertence} a esse intervalo — logo o máximo da função \textbf{no domínio permitido} não é o vértice.

\textbf{Passo 4 — Analisar o comportamento da função no intervalo.}

Como $a = -4 < 0$, a função $R(x)$ é crescente para $x < 125$ e decrescente para $x > 125$. Todo o intervalo $[200, 300]$ está à direita do vértice, portanto $R(x)$ é \textbf{decrescente} em todo esse intervalo.

\textbf{Passo 5 — Concluir onde ocorre o máximo restrito.}

Se a função é decrescente em $[200, 300]$, o maior valor de $R(x)$ nesse intervalo ocorre no \textbf{extremo esquerdo}, isto é, em $x = 200$.

\textbf{Passo 6 — Calcular a receita máxima.}

$$R(200) = 1000(200) - 4(200)^2 = 200000 - 160000 = 40000$$

Para confirmar que realmente é o maior valor, comparamos com o outro extremo:

$$R(300) = 1000(300) - 4(300)^2 = 300000 - 360000 = -60000$$

Como $R(200) = 40000 > R(300) = -60000$, confirma-se que o máximo no intervalo ocorre em $x = 200$.

\textbf{Resposta:} o preço que maximiza a receita, respeitando a cláusula contratual, é \textbf{R$\,200,00$}, gerando uma receita máxima de \textbf{R$\,40.000,00$}.

\bigskip\noindent\textbf{Q063 · discursiva · média · aplicar · teórica · gabarito conferido}\par\nopagebreak
Uma rede de lanchonetes possui duas filiais, A e B, que vendem o mesmo produto por preços que podem ser ajustados mensalmente. Os gestores modelaram o lucro mensal de cada filial (em milhares de reais) em função do preço de venda unitário $p$ (em reais) pelas funções:

$L_A(p) = -2p^2 + 80p - 600$

$L_B(p) = -3p^2 + 132p - 1200$

A diretoria da rede definiu uma meta: para que a linha de produto continue sendo vendida em uma filial, o lucro mensal máximo dessa filial deve ser de pelo menos 240 mil reais.

a) Determine o preço de venda que maximiza o lucro de cada filial e o valor do lucro máximo correspondente.

b) Qual filial alcança o maior lucro máximo?

c) Com base na meta estabelecida pela diretoria, verifique se cada filial deve continuar vendendo o produto. Justifique sua resposta com os valores calculados.

\textbf{Gabarito.} Filial A: preço de R\$\,20, lucro máximo de 200 mil reais (não atinge a meta). Filial B: preço de R\$\,22, lucro máximo de 252 mil reais (atinge a meta). A Filial B tem o maior lucro máximo e é a única que deve continuar vendendo o produto.

\textbf{Resolução.} \textbf{Passo 1 — Lucro máximo da Filial A}

A função $L_A(p) = -2p^2 + 80p - 600$ é uma parábola com concavidade para baixo ($a = -2 < 0$), logo possui ponto de \textbf{máximo} no vértice.

Coordenada $p$ do vértice: $p_A = -\dfrac{b}{2a} = -\dfrac{80}{2(-2)} = -\dfrac{80}{-4} = 20$

Lucro máximo: $L_A(20) = -2(20)^2 + 80(20) - 600 = -800 + 1600 - 600 = 200$

Logo, a Filial A tem lucro máximo de \textbf{200 mil reais}, obtido com preço de \textbf{R\$\,20}.

\textbf{Passo 2 — Lucro máximo da Filial B}

A função $L_B(p) = -3p^2 + 132p - 1200$ também tem concavidade para baixo ($a=-3<0$).

$p_B = -\dfrac{132}{2(-3)} = -\dfrac{132}{-6} = 22$

$L_B(22) = -3(22)^2 + 132(22) - 1200 = -3(484) + 2904 - 1200 = -1452 + 2904 - 1200 = 252$

Logo, a Filial B tem lucro máximo de \textbf{252 mil reais}, obtido com preço de \textbf{R\$\,22}.

\textbf{Passo 3 — Comparação entre as filiais}

Como $252 > 200$, a \textbf{Filial B} alcança o maior lucro máximo.

\textbf{Passo 4 — Verificação da meta da diretoria}

A meta exige lucro máximo de pelo menos 240 mil reais.

\begin{itemize}[nosep,leftmargin=*]
\item Filial A: lucro máximo de 200 mil reais $< 240$ mil reais $\to$ \textbf{não atinge a meta}, devendo descontinuar o produto.
\item Filial B: lucro máximo de 252 mil reais $\geq 240$ mil reais $\to$ \textbf{atinge a meta}, podendo continuar vendendo o produto.
\end{itemize}

\bigskip\noindent\textbf{Q071 · múltipla escolha · média · aplicar · aplicada · gabarito conferido}\par\nopagebreak
Uma confecção fabrica um certo modelo de camisa a um custo de produção de R\$\,40,00 por unidade. Um estudo de mercado mostrou que, se o preço de venda for $p$ reais por unidade, a quantidade vendida mensalmente será $q(p) = 300 - 5p$ unidades, válida para $0 < p < 60$. O lucro mensal da confecção, em reais, é dado pelo produto entre o lucro obtido em cada unidade vendida (preço de venda menos custo de produção) e a quantidade vendida naquele mês. Qual deve ser o preço de venda $p$, em reais, para que o lucro mensal da confecção seja máximo?
\begin{enumerate}[label=(\alph*)]
\item R\$\,100,00
\item R\$\,30,00
\item R\$\,60,00
\item R\$\,50,00
\end{enumerate}

\textbf{Gabarito.} R\$\,50,00 (alternativa d)

\textbf{Resolução.} \textbf{1. Montar a função lucro}

O lucro por unidade vendida é $(p-40)$ reais, e a quantidade vendida é $q(p)=300-5p$. Logo, o lucro mensal é:

$$L(p) = (p-40)(300-5p)$$

Expandindo:

$$L(p) = 300p - 5p^2 - 12000 + 200p = -5p^2 + 500p - 12000$$

\textbf{2. Reconhecer a natureza da parábola}

$L(p)$ é uma função quadrática com $a = -5 < 0$, portanto a parábola tem concavidade voltada para baixo e possui um \textbf{ponto de máximo} no vértice.

\textbf{3. Calcular o preço que maximiza o lucro (coordenada $p$ do vértice)}

$$p_v = -\frac{b}{2a} = -\frac{500}{2\cdot(-5)} = -\frac{500}{-10} = 50$$

Como $0 < 50 < 60$, esse preço pertence ao domínio válido do problema.

\textbf{4. (Verificação) Calcular o lucro máximo correspondente}

$$L(50) = -5(50)^2 + 500(50) - 12000 = -12500 + 25000 - 12000 = 500$$

Ou seja, ao praticar o preço de R\$\,50,00, a confecção vende $q(50)=300-250=50$ unidades, com lucro unitário de R\$\,10,00, totalizando R\$\,500,00 de lucro mensal — o maior valor possível para essa função.

\textbf{Conclusão:} o preço que maximiza o lucro mensal é \textbf{R\$\,50,00}.

\bigskip\noindent\textbf{Q076 · discursiva · fácil · analisar · aplicada · gabarito conferido}\par\nopagebreak
Uma loja de roupas vende um certo modelo de jaqueta. O custo de aquisição de cada peça, pago ao fornecedor, é de R\$\,80,00. Uma pesquisa de mercado mostrou que, se a loja fixar o preço de venda em $p$ reais (com $80 < p < 150$), a quantidade de jaquetas vendidas por mês, em unidades, é dada por $q(p) = 300 - 2p$.

a) Escreva a expressão do lucro mensal $L(p)$ obtido com a venda das jaquetas, em função do preço $p$, sabendo que o lucro é o produto entre o lucro unitário (preço de venda menos custo de aquisição) e a quantidade vendida.

b) Determine o preço $p$ que maximiza o lucro mensal e calcule esse lucro máximo.

c) Confirme que esse preço realmente fornece o maior lucro possível, calculando o lucro para dois preços vizinhos — R\$\,114,00 e R\$\,116,00 — e comparando os resultados com o lucro obtido no preço ótimo encontrado no item (b).

d) A diretoria da loja quer saber para quais preços o lucro mensal ultrapassa R\$\,2.400,00. Determine esse intervalo de preços.

e) Suponha agora que o custo de aquisição da jaqueta suba de R\$\,80,00 para um valor genérico $c$ reais, mantendo-se a mesma relação de demanda $q(p) = 300-2p$. Escreva o lucro $L(p,c)$ em função de $p$ e $c$, obtenha o preço ótimo em função de $c$ e use essa expressão para determinar de quanto aumenta o preço ótimo se o custo subir de R\$\,80,00 para R\$\,90,00.

\textbf{Gabarito.} a) $L(p) = -2p^2+460p-24000$; b) $p=115$ reais, lucro máximo $= 2450$ reais; c) $L(114)=L(116)=2448 < 2450$, confirmando o máximo; d) $110 < p < 120$; e) $p^*(c) = 75 + c/2$; aumento de R\$\,5,00 no preço ótimo quando o custo sobe R\$\,10,00.

\textbf{Resolução.} \textbf{a) Construção da função lucro}

O lucro unitário é (preço $-$ custo) $= p - 80$. A quantidade vendida é $q(p) = 300-2p$. Logo:
$$L(p) = (p-80)(300-2p) = 300p - 2p^2 - 24000 + 160p = -2p^2 + 460p - 24000$$

com $80 < p < 150$.

\textbf{b) Preço ótimo e lucro máximo}

Como $a = -2 < 0$, a parábola tem concavidade para baixo, logo possui ponto de \textbf{máximo} no vértice:
$$p^* = -\frac{b}{2a} = -\frac{460}{2(-2)} = 115$$
$$L(115) = -2(115)^2 + 460(115) - 24000 = -26450 + 52900 - 24000 = 2450$$

O preço ótimo é \textbf{R\$\,115,00}, com lucro mensal máximo de \textbf{R\$\,2.450,00}.

\textbf{c) Confirmação pela comparação com preços vizinhos}

$$L(114) = -2(114)^2+460(114)-24000 = -25992+52440-24000 = 2448$$
$$L(116) = -2(116)^2+460(116)-24000 = -26912+53360-24000 = 2448$$

Como $L(114) = L(116) = 2448 < 2450 = L(115)$, o lucro cresce até $p=115$ e depois decresce — confirmando que $p=115$ é de fato o ponto de máximo (comportamento simétrico em torno do vértice).

\textbf{d) Intervalo em que o lucro supera R\$\,2.400,00}

$$-2p^2+460p-24000 > 2400 \Rightarrow -2p^2+460p-26400>0$$

Dividindo por $-2$ (invertendo a desigualdade):
$$p^2 - 230p + 13200 < 0$$

Raízes: $\Delta = 230^2 - 4(13200) = 52900-52800=100$, $\sqrt{\Delta}=10$
$$p = \frac{230\pm10}{2} \Rightarrow p=110 \text{ ou } p=120$$

Como a parábola $p^2-230p+13200$ tem concavidade para cima, ela é negativa entre as raízes:
$$110 < p < 120$$

Ou seja, o lucro supera R\$\,2.400,00 para preços entre R\$\,110,00 e R\$\,120,00.

\textbf{e) Lucro em função de $p$ e de um custo genérico $c$}

$$L(p,c) = (p-c)(300-2p) = -2p^2 + (300+2c)p - 300c$$

O vértice (preço ótimo) é:
$$p^*(c) = \frac{300+2c}{4} = 75 + \frac{c}{2}$$

Verificação: para $c=80$, $p^*=75+40=115$ (confere com o item b).

Para $c=90$: $p^*(90) = 75+45=120$.

A variação do preço ótimo é:
$$\Delta p^* = p^*(90)-p^*(80) = 120-115 = 5$$

Ou seja, se o custo de aquisição subir R\$\,10,00, o preço ótimo deve subir apenas \textbf{R\$\,5,00} — metade do aumento do custo, pois $p^*(c)=75+c/2$ tem coeficiente angular $1/2$ em relação a $c$.

\bigskip\noindent\textbf{Q084 · múltipla escolha · fácil · analisar · teórica · gabarito conferido}\par\nopagebreak
Dois estudantes lançam, ao mesmo tempo, foguetes de brinquedo verticalmente para cima a partir do solo. As alturas atingidas pelos foguetes A e B, em metros, em função do tempo $t$ (em segundos) após o lançamento, são dadas por:

$h_A(t) = -5t^2 + 20t + 2$

$h_B(t) = -4t^2 + 18t + 3$

Qual das alternativas a seguir indica corretamente qual foguete atinge a maior altura máxima, o valor dessa altura, o instante em que ela ocorre, e justifica por que esse valor corresponde de fato a um ponto de máximo (e não de mínimo)?
\begin{enumerate}[label=(\alph*)]
\item O foguete A atinge a maior altura máxima, pois seu coeficiente de $t^2$ (-5) é mais negativo que o do foguete B (-4), o que indicaria uma parábola 'mais alta'.
\item O foguete B atinge a maior altura máxima, 23,25 m, no instante t = 2,25 s; como o coeficiente de $t^2$ é negativo em ambas as funções, seus gráficos têm concavidade voltada para baixo e o vértice é ponto de máximo.
\item O foguete B atinge a maior altura máxima, pois sua altura inicial (3 m) é maior que a do foguete A (2 m).
\item O foguete B atinge a maior altura máxima, 20,25 m, no instante t = 2,25 s.
\end{enumerate}

\textbf{Gabarito.} O foguete B atinge a maior altura máxima, 23,25 m, no instante t = 2,25 s; como o coeficiente de t² é negativo em ambas as funções, seus gráficos têm concavidade voltada para baixo e o vértice é ponto de máximo.

\textbf{Resolução.} \textbf{Passo 1 — Reconhecer que os vértices são pontos de máximo.}

Em ambas as funções o coeficiente de $t^2$ é negativo ($a_A=-5$ e $a_B=-4$), logo os gráficos são parábolas com concavidade voltada para baixo. Isso garante que o vértice de cada parábola corresponde a um \textbf{ponto de máximo}, não de mínimo.

\textbf{Passo 2 — Calcular o instante do máximo de cada foguete.}

Para o foguete A: $t_A = -\dfrac{b}{2a} = -\dfrac{20}{2\cdot(-5)} = 2$ s

Para o foguete B: $t_B = -\dfrac{b}{2a} = -\dfrac{18}{2\cdot(-4)} = 2{,}25$ s

\textbf{Passo 3 — Calcular a altura máxima de cada foguete substituindo $t_v$ na função.}

$h_A(2) = -5(2)^2 + 20(2) + 2 = -20 + 40 + 2 = 22$ m

$h_B(2{,}25) = -4(2{,}25)^2 + 18(2{,}25) + 3 = -20{,}25 + 40{,}5 + 3 = 23{,}25$ m

\textbf{Passo 4 — Comparar os dois resultados.}

Como $23{,}25 > 22$, o foguete \textbf{B} atinge a maior altura máxima, igual a \textbf{23,25 m}, no instante \textbf{t = 2,25 s}, e esse valor é de fato um ponto de máximo porque $a_B<0$.
\subsection*{EM13MAT507}
{\small\itshape Identificar e associar sequências numéricas (PA) a funções afins de domínios discretos para análise de propriedades, incluindo dedução de algumas fórmulas e resolução de problemas.}

\bigskip\noindent\textbf{Q045 · discursiva · fácil · aplicar · teórica · gabarito conferido}\par\nopagebreak
Considere a progressão aritmética $(a_n)$ com primeiro termo $a_1 = 7$ e razão $r = 4$, definida para $n = 1, 2, 3, \dots$

a) Determine a lei de uma função afim $f(n)$ tal que $f(n) = a_n$ para todo $n$ natural positivo, e explicite qual é o domínio dessa função.

b) Utilizando a função $f$ obtida no item (a), calcule o vigésimo termo da progressão, isto é, $a_{20}$.

\textbf{Gabarito.} f(n) = 4n + 3, com domínio $\mathbb{N}^* = \{1,2,3,\dots\}$; $a_{20} = 83$

\textbf{Resolução.} \textbf{a) Associando a PA a uma função afim}

O termo geral de uma PA é dado por $a_n = a_1 + (n-1)r$.

Substituindo $a_1 = 7$ e $r = 4$:

$a_n = 7 + (n-1)\cdot 4 = 7 + 4n - 4 = 4n + 3$

Essa expressão tem a forma $f(n) = 4n + 3$, que é exatamente uma função afim $f(x) = mx + b$ com coeficiente angular $m = 4$ (igual à razão $r$) e coeficiente linear $b = 3$.

Como o índice $n$ representa a posição do termo na sequência, só faz sentido para $n = 1, 2, 3, \dots$. Logo, o domínio de $f$ não é todo o conjunto dos reais (como seria numa função afim contínua), mas sim o conjunto dos números naturais positivos:

$D(f) = \mathbb{N}^* = \{1, 2, 3, \dots\}$

Ou seja, $(a_n)$ é a "restrição" da função afim $f(n) = 4n+3$ ao domínio discreto $\mathbb{N}^*$ — os pontos $(n, a_n)$ estão todos sobre a reta $y = 4x + 3$, mas apenas nos valores inteiros positivos de $x$.

\textbf{b) Cálculo de $a_{20}$}

Basta calcular $f(20)$:

$f(20) = 4 \cdot 20 + 3 = 80 + 3 = 83$

Portanto, $a_{20} = 83$.

\bigskip\noindent\textbf{Q053 · múltipla escolha · fácil · analisar · aplicada · gabarito conferido}\par\nopagebreak
Uma academia cobra R\$\,70,00 pela mensalidade do primeiro mês de plano. A partir do segundo mês, o valor da mensalidade aumenta R\$\,8,00 em relação ao mês imediatamente anterior, e esse padrão se mantém indefinidamente. Chame de $n$ o número do mês (com $n = 1, 2, 3, \dots$) e de $f(n)$ o valor da mensalidade cobrada nesse mês. Construa a lei que expressa $f(n)$ em função de $n$, levando em conta o conjunto de valores que $n$ pode assumir, e use essa lei para determinar o valor da mensalidade no 15º mês.
\begin{enumerate}[label=(\alph*)]
\item R\$\,120,00
\item R\$\,190,00
\item R\$\,182,00
\item R\$\,1.050,00
\end{enumerate}

\textbf{Gabarito.} R\$\,182,00

\textbf{Resolução.} \textbf{Passo 1 — Reconhecer o padrão da sequência.}

Os valores das mensalidades são $70, 78, 86, 94, \dots$. A diferença entre um termo e o anterior é sempre a mesma, $8$, o que caracteriza uma \textbf{progressão aritmética (PA)} de primeiro termo $a_1 = 70$ e razão $r = 8$.

\textbf{Passo 2 — Obter o termo geral da PA.}

$$a_n = a_1 + (n-1)\cdot r = 70 + 8(n-1) = 8n + 62$$

\textbf{Passo 3 — Associar a PA a uma função afim.}

A expressão $a_n = 8n + 62$ tem a forma $f(n) = an + b$, com $a=8$ e $b=62$, que é exatamente a lei de uma \textbf{função afim}. Assim, a mensalidade no mês $n$ pode ser escrita como

$$f(n) = 8n + 62.$$

\textbf{Passo 4 — Identificar o domínio.}

Diferente de uma função afim genérica $f(x) = ax+b$ definida para todo $x$ real, aqui $n$ representa a *ordem do mês* — só faz sentido para valores inteiros positivos ($n=1,2,3,\dots$). Logo, o domínio de $f$ é o conjunto dos números naturais positivos, $\mathbb{N}^\ast$ (domínio \textbf{discreto}), e não o conjunto dos reais.

\textbf{Passo 5 — Calcular o valor pedido.}

$$f(15) = 8\cdot 15 + 62 = 120 + 62 = 182.$$

Portanto, a mensalidade no 15º mês é \textbf{R\$\,182,00}.

\bigskip\noindent\textbf{Q066 · discursiva · difícil · analisar · teórica · gabarito conferido}\par\nopagebreak
Duas progressões aritméticas de números naturais são dadas por:

$a_n = 4n - 1,\; n \in \mathbb{N}^*$ (ou seja, $3, 7, 11, 15, 19, 23, 27, \dots$)

$b_m = 6m - 3,\; m \in \mathbb{N}^*$ (ou seja, $3, 9, 15, 21, 27, \dots$)

Cada uma dessas progressões pode ser vista como a restrição ao conjunto dos números naturais positivos de uma função afim: $f(n) = 4n-1$ e $g(m) = 6m-3$, respectivamente, com $n, m$ percorrendo $\mathbb{N}^*$.

Alguns números aparecem nas duas listas (por exemplo, $3$, $15$ e $27$). Chame de $c_k$, com $k = 0, 1, 2, \dots$, o $(k+1)$-ésimo desses termos comuns, em ordem crescente.

a) Mostre que a sequência $(c_k)$ também é uma progressão aritmética e determine seu primeiro termo $c_0$ e sua razão.

b) Escreva a função afim $h: \mathbb{N} \to \mathbb{R}$ tal que $h(k) = c_k$ para todo $k \in \mathbb{N}$.

c) Quantos termos comuns às duas progressões existem entre $1$ e $2023$, inclusive?

\textbf{Gabarito.} a) $(c_k)$ é PA com $c_0=3$ e razão $12$. b) $h(k)=12k+3$, $k\in\mathbb{N}$. c) $169$ termos comuns.

\textbf{Resolução.} \textbf{Passo 1 — Igualar as duas funções afins nos naturais.}

Um número é termo comum às duas progressões quando existem $n, m \in \mathbb{N}^*$ tais que $f(n) = g(m)$, isto é:
$$4n - 1 = 6m - 3$$
$$4n - 6m = -2 \quad\Longrightarrow\quad 2n - 3m = -1$$

\textbf{Passo 2 — Analisar quando a solução é natural (uso do domínio discreto).}

Isolando $n$: $n = \dfrac{3m-1}{2}$. Como $n$ deve ser um número natural, $3m - 1$ precisa ser par, o que ocorre exatamente quando $m$ é ímpar (pois $3m$ tem a mesma paridade de $m$). Isso é uma restrição que só faz sentido porque o domínio de $f$ e $g$ é discreto — na reta real qualquer $m$ serviria, mas aqui só interessam os pares $(n,m)$ de naturais.

Escrevendo os ímpares como $m = 2k+1$, $k = 0,1,2,\dots$:
$$n = \frac{3(2k+1)-1}{2} = \frac{6k+2}{2} = 3k+1$$

\textbf{Passo 3 — Calcular o termo comum $c_k$.}

Substituindo em $f(n) = 4n-1$:
$$c_k = 4(3k+1) - 1 = 12k + 3$$

(Confirmação por $g$: $b_m = 6(2k+1)-3 = 12k+6-3 = 12k+3$ — mesma expressão.)

\textbf{Passo 4 — Reconhecer a nova PA.}

$c_k = 12k+3$ é uma função afim de $k$, logo $(c_k)$ é uma PA de primeiro termo $c_0 = 3$ e razão $12$ (note que $12 = \text{mmc}(4,6)$, resultado esperado quando se intersectam duas PAs).

\textbf{Passo 5 — Função afim associada.}

$$h: \mathbb{N} \to \mathbb{R}, \qquad h(k) = 12k+3$$
com domínio discreto $\mathbb{N} = \{0,1,2,\dots\}$, coeficiente angular (razão) $12$ e coeficiente linear $3$.

\textbf{Passo 6 — Contar os termos comuns até 2023.}

Queremos o maior $k$ natural tal que $h(k) \le 2023$:
$$12k+3 \le 2023 \;\Longrightarrow\; 12k \le 2020 \;\Longrightarrow\; k \le \frac{2020}{12} = 168{,}33\ldots$$

Como $k$ é natural, o maior valor possível é $k=168$. Verificando os limites:
$$h(168) = 12(168)+3 = 2019 \le 2023 \quad\checkmark$$
$$h(169) = 12(169)+3 = 2031 > 2023 \quad(\text{excede})$$

Os valores válidos de $k$ são $0, 1, 2, \dots, 168$, totalizando:
$$168 - 0 + 1 = 169 \text{ termos comuns}$$

\bigskip\noindent\textbf{Q074 · múltipla escolha · difícil · aplicar · aplicada · gabarito conferido}\par\nopagebreak
Uma empresa de transporte urbano inaugurou uma nova linha de ônibus. No primeiro mês de operação (mês 1), a linha transportou 340 passageiros. Devido à divulgação e à adesão gradual dos moradores, o número de passageiros transportados em cada mês superou o do mês anterior em exatamente 45 passageiros, formando uma progressão aritmética que se manteve mês após mês.

A diretoria da empresa sabe que essa evolução mensal pode ser descrita por uma função afim $f$, definida apenas para os números naturais $n \geq 1$ (o mês de operação), de modo que $f(n)$ forneça o número de passageiros transportados no mês $n$, coincidindo com os termos da progressão aritmética.

Com base nessa função afim, determine a partir de qual mês o número de passageiros transportados ultrapassará, pela primeira vez, 10.000 passageiros mensais.
\begin{enumerate}[label=(\alph*)]
\item Mês 16
\item Mês 217
\item Mês 215
\item Mês 216
\end{enumerate}

\textbf{Gabarito.} Mês 216 (alternativa D)

\textbf{Resolução.} \textbf{1. Identificando a progressão aritmética}

O número de passageiros no mês $n$ forma uma PA com primeiro termo $a_1 = 340$ e razão $r = 45$.

\textbf{2. Associando a PA à função afim}

O termo geral da PA é:
$$a_n = a_1 + (n-1)r = 340 + 45(n-1) = 45n + 295$$

Como a evolução mensal é descrita por uma função afim $f$ que coincide com a PA nos naturais, temos:
$$f(n) = 45n + 295, \quad n \in \mathbb{N}, \; n \geq 1$$

Note que $f$ é a restrição aos naturais da função afim $g(x) = 45x + 295$ definida em $\mathbb{R}$: a reta que passa pelos pontos $(1,340)$, $(2,385)$, $(3,430)$, etc. É essa correspondência entre a PA e a reta que permite tratar o problema algebricamente como uma inequação, mesmo sabendo que só fazem sentido valores naturais de $n$.

\textbf{3. Resolvendo a inequação (em $\mathbb{R}$, para localizar a região)}

Queremos o menor $n$ tal que $f(n) > 10000$:
$$45n + 295 > 10000$$
$$45n > 9705$$
$$n > \frac{9705}{45} = 215{,}67\ldots$$

\textbf{4. Voltando ao domínio discreto}

Como $n$ representa um número de mês, deve ser um número natural. O menor natural maior que $215{,}67$ é $n = 216$ (e não $215$, pois $215 < 215{,}67$).

\textbf{5. Verificação}

$$a_{215} = 45(215) + 295 = 9675 + 295 = 9970 < 10000$$
$$a_{216} = 45(216) + 295 = 9720 + 295 = 10015 > 10000$$

De fato, o mês 215 ainda não ultrapassa 10.000, e o mês 216 é o primeiro a ultrapassar.

\textbf{Resposta: mês 216.}

\bigskip\noindent\textbf{Q079 · discursiva · média · analisar · aplicada · gabarito conferido}\par\nopagebreak
Uma marcenaria iniciou a fabricação de cadeiras artesanais em regime de produção mensal constante. No primeiro mês de operação ($n=1$) foram fabricadas 120 cadeiras. A partir do segundo mês, a quantidade produzida a cada mês aumenta sempre do mesmo valor fixo em relação ao mês anterior, de modo que no quarto mês ($n=4$) a produção atingiu 165 cadeiras. Sabe-se que esse ritmo de crescimento se mantém constante indefinidamente, e que $n$ representa o número do mês de produção ($n = 1, 2, 3, \dots$).

a) Determine a razão da progressão aritmética formada pela produção mensal e escreva a lei de uma função afim $f(n)$ que forneça a quantidade de cadeiras fabricadas no mês $n$, indicando explicitamente o domínio dessa função.

b) Usando a função $f(n)$ obtida, determine em que mês a produção mensal será de 300 cadeiras.

c) Calcule o total de cadeiras fabricadas do 1º ao 10º mês.

\textbf{Gabarito.} a) $r = 15$; $f(n) = 15n + 105$, com $n \in \mathbb{N}^* = \{1,2,3,\dots\}$. b) 13º mês. c) 1875 cadeiras.

\textbf{Resolução.} \textbf{a) Associando a PA à função afim}

A produção mensal forma uma progressão aritmética, pois cresce sempre da mesma quantidade fixa a cada mês. Temos $a_1 = 120$ e $a_4 = 165$.

Como $a_4 = a_1 + 3r$:
$$165 = 120 + 3r \implies 3r = 45 \implies r = 15$$

O termo geral da PA é:
$$a_n = a_1 + (n-1)r = 120 + 15(n-1) = 15n + 105$$

Como $a_n$ é da forma $f(n) = 15n + 105$, uma expressão linear em $n$ com coeficiente angular $15$ e coeficiente linear $105$, essa é exatamente a lei de uma \textbf{função afim}. A diferença essencial em relação a uma função afim de domínio contínuo é que aqui a variável $n$ representa o número do mês, que só assume valores naturais positivos. Logo:
$$f(n) = 15n + 105, \quad n \in \mathbb{N}^* = \{1, 2, 3, \dots\}$$

ou seja, o gráfico de $f$ é formado por pontos isolados (discretos) sobre a reta $y = 15x + 105$, e não por uma reta contínua.

\textbf{b) Determinando o mês em que a produção é 300}

Queremos $n$ tal que $f(n) = 300$:
$$15n + 105 = 300$$
$$15n = 195$$
$$n = 13$$

Como $13 \in \mathbb{N}^*$, essa solução é compatível com o domínio discreto da função: a produção de 300 cadeiras ocorre no \textbf{13º mês}.

\textbf{c) Soma da produção do 1º ao 10º mês}

Usando a fórmula da soma dos $n$ primeiros termos de uma PA:
$$S_{10} = \frac{n}{2}\left(2a_1 + (n-1)r\right) = \frac{10}{2}\left(2\cdot 120 + 9\cdot 15\right)$$
$$S_{10} = 5\left(240 + 135\right) = 5 \cdot 375 = 1875$$

Foram fabricadas \textbf{1875 cadeiras} do 1º ao 10º mês.

\bigskip\noindent\textbf{Q087 · múltipla escolha · média · aplicar · teórica · gabarito conferido}\par\nopagebreak
Considere a progressão aritmética $(a_n)$ com primeiro termo $a_1 = 7$ e razão $r = 4$: $(7, 11, 15, 19, 23, \dots)$. Essa PA pode ser vista como a restrição da função afim $f(x) = 4x + 3$ ao conjunto dos números naturais não nulos, de modo que $a_n = f(n)$ para todo $n \geq 1$ (o gráfico da PA é o conjunto de pontos discretos $(n, a_n)$ sobre a reta que representa $f$).

Qual é o menor termo dessa PA que é maior do que 100?
\begin{enumerate}[label=(\alph*)]
\item 100
\item 99
\item 103
\item 107
\end{enumerate}

\textbf{Gabarito.} 103

\textbf{Resolução.} \textbf{Passo 1 — Associar a PA à função afim.}
O termo geral da PA é $a_n = a_1 + (n-1)r = 7 + 4(n-1) = 4n + 3$. Logo, a PA corresponde à função afim $f(x) = 4x+3$ restrita aos naturais não nulos, isto é, $a_n = f(n)$, com $n \in \{1, 2, 3, \dots\}$.

\textbf{Passo 2 — Resolver a condição sobre os números reais.}
Queremos o menor termo maior que 100, ou seja, o menor $n$ tal que $f(n) > 100$. Primeiro resolvemos a inequação como se $x$ fosse real:
$$4x + 3 > 100 \implies 4x > 97 \implies x > 24{,}25.$$

\textbf{Passo 3 — Voltar ao domínio discreto da PA.}
Como $n$ representa a posição de um termo, ele precisa ser um número \textbf{natural}. O menor número natural maior que $24{,}25$ é $n = 25$ (e não $24$, que não satisfaz a inequação, nem $24{,}25$, que não é índice de termo algum).

\textbf{Passo 4 — Calcular o termo.}
$$a_{25} = f(25) = 4(25) + 3 = 103.$$

Como verificação: $a_{24} = 4(24)+3 = 99 < 100$ e $a_{25} = 103 > 100$, confirmando que $a_{25}=103$ é o menor termo da PA maior que 100.
\subsection*{EM13MAT508}
{\small\itshape Identificar e associar sequências numéricas (PG) a funções exponenciais de domínios discretos para análise de propriedades, incluindo dedução de algumas fórmulas e resolução de problemas.}

\bigskip\noindent\textbf{Q048 · discursiva · média · analisar · teórica · gabarito conferido}\par\nopagebreak
Considere a progressão geométrica $(a_n)$, com $n \in \mathbb{N}^{*}$ (ou seja, $n = 1, 2, 3, \dots$), definida por $a_1 = 4$ e razão $q = 2$.

a) Determine a lei de uma função exponencial $f:\mathbb{R}\to\mathbb{R}$ tal que $f(n) = a_n$ para todo $n \in \mathbb{N}^{*}$, explicitando como o primeiro termo e a razão da PG se tornam os parâmetros dessa função.

b) Compare os domínios de $(a_n)$ e de $f$, explicando por que a progressão pode ser vista como a restrição da função $f$ a um subconjunto discreto de $\mathbb{R}$.

c) Determine o valor real de $x$ para o qual $f(x) = 24$. Em seguida, usando o domínio da PG, decida (com justificativa) se existe algum termo $a_n$ da progressão igual a 24.

\textbf{Gabarito.} a) $f(x) = 4\cdot 2^{x-1}$ (equivalente a $2^{x+1}$), com $a_1=4$ e $q=2$ tornando-se, respectivamente, o valor inicial e a base da exponencial. b) $(a_n)$ tem domínio discreto $\mathbb{N}^{*}$ e é a restrição de $f$, que tem domínio contínuo $\mathbb{R}$, ao conjunto $\mathbb{N}^{*}$. c) $x=\log_2 12\approx 3{,}585$, que não é natural; logo não existe termo da PG igual a 24 (situa-se entre $a_3=16$ e $a_4=32$).

\textbf{Resolução.} \textbf{a) Lei da função exponencial}

O termo geral da PG é $a_n = a_1 \cdot q^{\,n-1} = 4\cdot 2^{\,n-1}$.

Para obter uma função definida em todo $\mathbb{R}$ que coincida com $a_n$ quando $x=n$ é natural, basta trocar $n$ por uma variável real $x$, mantendo o primeiro termo como o valor em $x=1$ e a razão como a base da potência:
$$f(x) = 4\cdot 2^{\,x-1}.$$
Os parâmetros de $f$ são exatamente os da PG: $a_1=4$ é o valor de $f$ em $x=1$, e $q=2$ é a base exponencial (a razão de crescimento por unidade de $x$).

\textbf{b) Comparação dos domínios}

A sequência $(a_n)$ só está definida para $n \in \mathbb{N}^{*} = \{1,2,3,\dots\}$, um conjunto \textbf{discreto}. Já a função $f(x)=4\cdot2^{x-1}$ está definida para \textbf{todo} $x \in \mathbb{R}$, formando um conjunto contínuo. Como $f(n) = 4\cdot 2^{n-1} = a_n$ para cada $n\in\mathbb{N}^{*}$, a progressão geométrica é exatamente a restrição de $f$ ao subconjunto discreto $\mathbb{N}^{*}\subset\mathbb{R}$: os pontos do gráfico de $(a_n)$ são pontos isolados sobre a curva contínua de $f$.

\textbf{c) Resolvendo $f(x)=24$}

$$4\cdot 2^{x-1} = 24 \implies 2^{x-1} = 6 \implies x - 1 = \log_2 6 \implies x = 1 + \log_2 6.$$

Como $1+\log_2 6 = \log_2 2 + \log_2 6 = \log_2 12$, temos
$$x = \log_2 12 \approx 3{,}585.$$

Esse valor não é um número natural. Como o domínio da PG é $\mathbb{N}^{*}$ (apenas os inteiros positivos), e $x \notin \mathbb{N}^{*}$, \textbf{não existe} termo $a_n$ igual a 24. De fato, calculando os termos vizinhos: $a_3 = 4\cdot2^2=16$ e $a_4=4\cdot2^3=32$, confirmando que $24$ fica entre dois termos consecutivos, sem corresponder a nenhum índice natural.

\bigskip\noindent\textbf{Q056 · múltipla escolha · média · aplicar · aplicada · gabarito conferido}\par\nopagebreak
Um biólogo monitora o crescimento de uma colônia de bactérias em um experimento de laboratório, registrando a quantidade de indivíduos ao final de cada hora, a partir da 1ª hora de observação. Os dados coletados nas quatro primeiras horas foram:

\begin{center}
\begin{tabular}{|p{0.430\linewidth}|p{0.430\linewidth}|}
\hline
\textbf{Hora (n)} & \textbf{Quantidade de bactérias} \\
\hline
1 & 100 \\
\hline
2 & 300 \\
\hline
3 & 900 \\
\hline
4 & 2700 \\
\hline
\end{tabular}
\end{center}

Essa quantidade cresce segundo o mesmo padrão observado nessas quatro medições, formando uma progressão geométrica. O biólogo deseja expressar o número de bactérias por meio de uma função exponencial $f(n)$, definida para $n$ natural, $n \geq 1$, que reproduza exatamente os valores registrados na tabela.

Assinale a alternativa que apresenta corretamente essa função $f(n)$ e o número de bactérias previsto para a 7ª hora de observação.
\begin{enumerate}[label=(\alph*)]
\item $f(n) = 100 \cdot 3^{n-1}$; na 7ª hora haverá $109300$ bactérias, calculadas pela soma de todos os termos da progressão até a 7ª hora.
\item $f(n) = 100 + 200(n-1)$; na 7ª hora haverá $1300$ bactérias.
\item $f(n) = 100 \cdot 3^{n-1}$; na 7ª hora haverá $72900$ bactérias.
\item $f(n) = 100 \cdot 3^{n}$; na 7ª hora haverá $218700$ bactérias.
\end{enumerate}

\textbf{Gabarito.} f(n) = 100·3\textasciicircum{}(n-1); f(7) = 72900 bactérias

\textbf{Resolução.} \textbf{Passo 1 — Reconhecer a progressão geométrica.}

Os valores $100, 300, 900, 2700$ satisfazem $\dfrac{300}{100} = \dfrac{900}{300} = \dfrac{2700}{900} = 3$, logo constituem uma PG de primeiro termo $a_1 = 100$ e razão $q = 3$.

\textbf{Passo 2 — Termo geral da PG.}

$a_n = a_1 \cdot q^{\,n-1} = 100 \cdot 3^{\,n-1}$

\textbf{Passo 3 — Associar a PG à função exponencial de domínio discreto.}

A sequência $(a_n)$ é a restrição, aos números naturais $n \geq 1$, da função exponencial $f(n) = 100 \cdot 3^{\,n-1}$. Ou seja, $f(1)=100$, $f(2)=300$, $f(3)=900$, $f(4)=2700$, coincidindo com os dados da tabela — a PG é exatamente a "amostragem" dessa função exponencial nos pontos inteiros.

\textbf{Passo 4 — Calcular a previsão para a 7ª hora.}

$f(7) = 100 \cdot 3^{7-1} = 100 \cdot 3^{6} = 100 \cdot 729 = 72900$

\textbf{Conclusão:} $f(n) = 100 \cdot 3^{n-1}$ e, na 7ª hora, há $72900$ bactérias.

\bigskip\noindent\textbf{Q061 · discursiva · fácil · aplicar · aplicada · gabarito conferido · corrigida após a avaliação docente}\par\nopagebreak
Um biólogo inicia uma cultura de bactérias com 200 indivíduos. A cada hora que passa, a quantidade de bactérias na cultura dobra em relação à hora anterior. Seja $a_n$ o número de bactérias existentes ao final da $n$-ésima hora de observação, com $n = 1, 2, 3, \dots$ (isto é, $a_1$ é a quantidade ao final da 1ª hora, $a_2$ ao final da 2ª hora, e assim por diante).

a) Mostre que a sequência $(a_n)$ é uma progressão geométrica, indicando o primeiro termo e a razão.

b) Escreva a lei $f(n)$ de uma função exponencial que descreva exatamente essa progressão geométrica, especificando qual é o domínio dessa função e por que ele deve ser desse tipo (e não todos os números reais).

c) Usando a fórmula do termo geral da PG (ou, equivalentemente, a função $f(n)$ obtida no item b), determine quantas bactérias haverá ao final da 6ª hora de observação.

\textbf{Gabarito.} a) PG de razão $q=2$ e $a_1=400$ (as 200 iniciais já dobraram ao final da 1ª hora); b) $f(n)=400\cdot 2^{n-1}=200\cdot 2^{n}$, com domínio $D=\{1,2,3,\dots\}$ (naturais, pois $n$ conta horas); c) $a_6 = f(6) = 12800$ bactérias.

\textbf{Resolução.} \textbf{a) Identificando a PG}

Como a quantidade dobra a cada hora, temos $a_{n+1} = 2\cdot a_n$ para todo $n \geq 1$. Isso significa que a razão entre termos consecutivos é constante e igual a $q = 2$. Logo, $(a_n)$ é uma progressão geométrica de razão $q=2$.

O primeiro termo exige atenção: as 200 bactérias são a quantidade \textbf{no início} da observação, e $a_1$ é a quantidade \textbf{ao final} da 1ª hora, quando ela já dobrou uma vez. Assim, $a_1 = 2 \cdot 200 = 400$.

\textbf{b) Associando a PG a uma função exponencial}

O termo geral da PG é dado por
$$a_n = a_1 \cdot q^{\,n-1} = 400 \cdot 2^{\,n-1} = 200 \cdot 2^{\,n}.$$

Essa expressão é exatamente a lei de uma função exponencial
$$f(n) = 200 \cdot 2^{\,n},$$
em que a base da exponencial ($2$) coincide com a razão da PG e o coeficiente ($200$) é a quantidade inicial — o valor que a curva assumiria em $n = 0$, instante que não pertence à sequência.

Como $n$ representa o número da hora de observação (1ª hora, 2ª hora, etc.), $n$ só pode assumir valores inteiros positivos: $n \in \{1, 2, 3, \dots\}$. Assim, o domínio de $f$ é o conjunto dos números naturais (a partir de 1), $D = \mathbb{N}^* = \{1,2,3,\dots\}$, e \textbf{não} todo o conjunto dos reais — não faz sentido, por exemplo, falar em "1,5 hora de observação" nesse contexto discreto de contagem de termos da sequência. Por isso $f$ é uma função exponencial de domínio discreto que apenas "amostra" a curva exponencial contínua $y = 200\cdot 2^{x}$ nos pontos inteiros.

\textbf{c) Cálculo para a 6ª hora}

Usando $n=6$:
$$f(6) = a_6 = 200 \cdot 2^{6} = 200 \cdot 64 = 12800.$$

Logo, ao final da 6ª hora haverá \textbf{12800 bactérias}.

\bigskip\noindent\textbf{Q069 · múltipla escolha · fácil · analisar · teórica · gabarito conferido}\par\nopagebreak
Considere a sequência $(a_n)$, definida para $n = 0, 1, 2, 3, \ldots$, com primeiro termo $a_0 = 5$ e razão $q = 3$ (uma progressão geométrica). Um pesquisador registra o número de bactérias em uma cultura a cada hora exata, de modo que $a_n$ é a contagem observada na hora $n$. Ele percebe que esses valores coincidem com os da função exponencial $f(t) = 5 \cdot 3^t$ para $t = n$, ou seja, $f(n) = a_n$ para todo $n$ natural.

Ricardo, colega do pesquisador, argumenta: "Como $f(t)=5\cdot 3^t$ está definida para qualquer número real $t$, o valor $f(2{,}5) = 5\cdot 3^{2{,}5}$ é a contagem real de bactérias no instante intermediário entre a segunda e a terceira hora."

Avalie a afirmação de Ricardo e assinale a alternativa correta.
\begin{enumerate}[label=(\alph*)]
\item Ricardo está errado, mas pelo motivo errado: ele deveria ter estimado o valor em $t=2{,}5$ fazendo a média aritmética simples entre $a_2=45$ e $a_3=135$ (interpolação linear), e não usando a fórmula exponencial $5\cdot3^{2,5}$.
\item Ricardo está errado: a progressão geométrica $(a_n)$ está definida apenas para os índices naturais $n=0,1,2,3,\ldots$, que correspondem às horas exatas em que a contagem foi feita. A função $f(t)=5\cdot3^t$ é uma extensão contínua que apenas interpola os termos da PG nos pontos inteiros; o valor $f(2{,}5)$ existe matematicamente, mas não corresponde a nenhum termo da sequência, pois o fenômeno só foi definido (e medido) nos instantes discretos.
\item Ricardo está certo, porque, em princípio, o número de bactérias poderia ser medido continuamente ao longo do tempo; logo $f(2{,}5)$ representa a contagem real esperada nesse instante, ainda que o pesquisador só tenha registrado valores a cada hora.
\item Ricardo está certo, pois a PG e a função exponencial são exatamente a mesma coisa: como $f$ está definida para todo $t$ real, qualquer valor calculado por $f$ é automaticamente um termo válido da progressão.
\end{enumerate}

\textbf{Gabarito.} B

\textbf{Resolução.} \textbf{1. Identificando o que é a PG e o que é a função exponencial.}

A sequência $(a_n)$ é uma progressão geométrica: $a_0=5$, $q=3$, então $a_n = 5\cdot 3^n$, mas \textbf{apenas para $n = 0,1,2,3,\ldots$} (índices naturais), pois cada termo corresponde a uma contagem feita em uma hora exata.

A função $f(t) = 5\cdot 3^t$ tem domínio real ($t$ pode ser qualquer número real), e ela foi construída de modo que $f(n) = a_n$ para todo $n$ natural — ou seja, $f$ é uma \textbf{extensão contínua} dos termos discretos da PG.

\textbf{2. Calculando alguns termos para comparação.}

$a_0 = 5$, $a_1 = 15$, $a_2 = 45$, $a_3 = 135$.

De fato, $f(0)=5$, $f(1)=15$, $f(2)=45$, $f(3)=135$: os valores coincidem exatamente nos pontos inteiros.

\textbf{3. Analisando $f(2{,}5)$.}

Como $f$ tem domínio real, o número $f(2{,}5)=5\cdot 3^{2,5}\approx 77{,}9$ existe matematicamente. Porém, a progressão $(a_n)$ — que é o modelo que de fato representa o fenômeno medido (contagens feitas hora a hora) — \textbf{só está definida para índices naturais}. Não existe "$a_{2,5}$": o pesquisador não mediu bactérias em $t=2{,}5$, e o modelo discreto não faz essa previsão.

\textbf{4. Conclusão.}

$f(t)=5\cdot3^t$ é apenas a curva contínua que passa pelos pontos $(n, a_n)$ para interpolar visualmente a PG; o valor $f(2{,}5)$ é um artifício matemático de extensão da fórmula, sem correspondência com um termo real da sequência, que só existe para $n$ natural. Logo, Ricardo está errado.

\textbf{Gabarito: alternativa B.}

\bigskip\noindent\textbf{Q082 · discursiva · difícil · aplicar · aplicada · gabarito conferido}\par\nopagebreak
Um biólogo inicia, no instante t = 0 (medido em horas), uma cultura com 500 bactérias em condições ideais de crescimento. Ele verifica que, a cada 4 horas, a quantidade de bactérias triplica. Por questões práticas, o técnico responsável só registra a contagem exatamente nos instantes t = 0, 4, 8, 12, 16, ... horas, obtendo assim uma sequência numérica de valores registrados, que pode ser indexada por n = 0, 1, 2, 3, ...

a) Escreva o termo geral $a_n$ da progressão geométrica formada pelos valores registrados pelo técnico (em função de n).

b) Sabendo que, na realidade, a população de bactérias cresce continuamente (não apenas nos instantes de registro), escreva uma função exponencial $f(t)$, com $t \geq 0$ representando o tempo em horas desde o início do experimento, tal que os valores registrados pelo técnico sejam exatamente os valores de $f$ nos instantes t = 4n, ou seja, $a_n = f(4n)$ para todo n natural.

c) Usando a função $f$, determine em que instante t (em horas, com uma casa decimal) a população da cultura atinge 20000 bactérias. Esse instante corresponde a um dos momentos em que o técnico faz um registro (múltiplo de 4 horas)? Justifique.

\textbf{Gabarito.} a) $a_n = 500\cdot 3^n$; b) $f(t) = 500\cdot 3^{t/4}$, com $f(4n)=a_n$; c) $t = 4\log_3 40 \approx 13{,}4$ horas, que não é múltiplo de 4, logo não coincide com nenhum registro do técnico.

\textbf{Resolução.} \textbf{a) Termo geral da PG}

Os valores registrados formam uma progressão geométrica de primeiro termo $a_0 = 500$ e razão $q = 3$ (pois a quantidade triplica a cada registro). Logo:
$$a_n = 500 \cdot 3^n, \quad n = 0,1,2,3,\dots$$

Conferindo: $a_0=500$, $a_1=1500$, $a_2=4500$, $a_3=13500$, $a_4=40500$.

\textbf{b) Função exponencial associada}

A PG descreve apenas os instantes discretos $t=0,4,8,12,\dots$, ou seja, $t = 4n \Rightarrow n = t/4$. Para obter uma função contínua $f(t)$ que "estenda" a PG a qualquer instante real, basta substituir $n$ por $t/4$ na fórmula do termo geral:
$$f(t) = 500 \cdot 3^{t/4}, \quad t \geq 0$$

Verificação da associação PG–função exponencial: para $t = 4n$,
$$f(4n) = 500\cdot 3^{4n/4} = 500\cdot 3^n = a_n.$$
Ou seja, a sequência $(a_n)$ é exatamente a restrição de $f$ ao domínio discreto $\{0,4,8,12,\dots\}$, enquanto $f$ está definida para todo $t$ real não negativo (domínio contínuo). Isso confirma que a PG é o "retrato" de uma função exponencial amostrada em intervalos regulares.

Conferência pontual: $f(8) = 500\cdot 3^{8/4} = 500\cdot 3^2 = 4500 = a_2$. $\checkmark$

\textbf{c) Instante em que a população atinge 20000 bactérias}

Resolvendo $f(t) = 20000$:
$$500\cdot 3^{t/4} = 20000 \;\Rightarrow\; 3^{t/4} = 40$$
Aplicando logaritmo (base 3, ou natural):
$$\frac{t}{4} = \log_3 40 = \frac{\ln 40}{\ln 3} \approx \frac{3{,}689}{1{,}099} \approx 3{,}358$$
$$t \approx 4 \times 3{,}358 \approx 13{,}4 \text{ horas}$$

Como $t \approx 13{,}4$ não é múltiplo de 4, esse instante \textbf{não coincide} com nenhum registro do técnico. De fato, $t=13{,}4$ está entre os registros $t=12$ (onde $a_3 = 13500$) e $t=16$ (onde $a_4 = 40500$), confirmando que 20000 bactérias é um valor atingido *entre* duas contagens discretas — algo que só a função exponencial contínua $f(t)$ permite calcular, pois a PG só fornece informação nos instantes $t=4n$.

\bigskip\noindent\textbf{Q090 · múltipla escolha · difícil · aplicar · teórica · gabarito conferido}\par\nopagebreak
Uma progressão geométrica $(a_n)$, definida apenas para $n \in \mathbb{N}^*$ (isto é, $n = 1, 2, 3, \dots$), tem $a_1 = 4$ e $a_4 = 32$. Sabe-se que essa PG é a restrição ao conjunto $\{1,2,3,\dots\}$ de uma função exponencial $f:\mathbb{Z}\to\mathbb{R}$, dada por $f(n) = A\cdot B^n$, tal que $f(n) = a_n$ para todo $n \geq 1$. Determine o valor de $f(0)$, isto é, o valor que a função exponencial assume fora do domínio original da progressão.
\begin{enumerate}[label=(\alph*)]
\item $f(0) = \dfrac{1}{2}$
\item $f(0) = 4$
\item $f(0) = 8$
\item $f(0) = 2$
\end{enumerate}

\textbf{Gabarito.} f(0) = 2

\textbf{Resolução.} \textbf{Passo 1 — Termo geral da PG.}

Como $(a_n)$ é uma progressão geométrica de razão $q$, seu termo geral é $a_n = a_1\cdot q^{n-1}$.

\textbf{Passo 2 — Determinar a razão $q$.}

Usando $a_4 = a_1\cdot q^3$:

$$32 = 4\cdot q^3 \Rightarrow q^3 = 8 \Rightarrow q = 2.$$

\textbf{Passo 3 — Associar a PG à função exponencial.}

Queremos escrever $a_n = A\cdot B^n$ (com domínio estendido a todo $\mathbb{Z}$), coincidindo com $a_n = a_1\cdot q^{n-1}$ para $n\geq 1$. Reescrevendo:

$$a_1\cdot q^{n-1} = \left(\dfrac{a_1}{q}\right)\cdot q^{n}.$$

Comparando com $A\cdot B^n$, temos $B = q = 2$ e $A = \dfrac{a_1}{q} = \dfrac{4}{2} = 2$.

\textbf{Passo 4 — Verificação.}

$f(1) = A\cdot B = 2\cdot 2 = 4 = a_1$ $\checkmark$

$f(4) = A\cdot B^4 = 2\cdot 16 = 32 = a_4$ $\checkmark$

A função $f(n) = 2\cdot 2^n$ realmente coincide com a PG para todo $n\geq 1$, sendo sua extensão natural a todos os inteiros (domínio discreto ampliado).

\textbf{Passo 5 — Calcular $f(0)$.}

$$f(0) = A\cdot B^0 = A = 2.$$

Portanto, $f(0) = 2$, valor que não pertence ao domínio original da PG (que começa em $n=1$), mas é obtido naturalmente pela função exponencial associada.
\end{document}